% Môn: Vật lý
% Lớp: 12
% Chương: Chương II. Khí lí tưởng
% Bài: L12C2 Bài 11. Phương trình trạng thái của khí lí tưởng
% Số câu: 243
% App đọc file này trực tiếp — sửa rồi Commit trên GitHub, bấm Đồng bộ GitHub .tex

% L12C2 Bài 11. Phương trình trạng thái của khí lí tưởng

% ===== Câu 1 =====
% ID: L12C2B11-01-DS
% Mức: TH
\dangbt{Bài toán thực tế dùng phương trình trạng thái}
\begin{ex}
% ID: L12C2B11-01-DS
% Mức: TH
Xét các phát biểu sau trong dạng “Bài toán thực tế dùng phương trình trạng thái”.
\choiceTF
{\True Phương trình Clapeyron của khí lí tưởng là — phát biểu đúng là: $pV=nRT$.}
{Phương trình Clapeyron của khí lí tưởng là — phát biểu $pV=RT/n$ là đúng.}
{\True Với lượng khí xác định, phương trình trạng thái giữa hai trạng thái là — phát biểu đúng là: $p_1V_1/T_1=p_2V_2/T_2$.}
{Với lượng khí xác định, phương trình trạng thái giữa hai trạng thái là — phát biểu $p_1/T_1=p_2/T_2$ trong mọi quá trình là đúng.}
\loigiai{
\begin{itemchoice}
\item (Đúng) Phương trình Clapeyron của khí lí tưởng là — phát biểu đúng là: $pV=nRT$. (Vì): Một lượng khí lí tưởng thỏa $pV=nRT$. b) Sai: trái với hệ thức hoặc điều kiện đã nêu. c) Đúng: Do $pV/T=nR$ không đổi. d) Sai: trái với hệ thức hoặc điều kiện đã nêu
\item (Sai) Phương trình Clapeyron của khí lí tưởng là — phát biểu $pV=RT/n$ là đúng.
\item (Đúng) Với lượng khí xác định, phương trình trạng thái giữa hai trạng thái là — phát biểu đúng là: $p_1V_1/T_1=p_2V_2/T_2$.
\item (Sai) Với lượng khí xác định, phương trình trạng thái giữa hai trạng thái là — phát biểu $p_1/T_1=p_2/T_2$ trong mọi quá trình là đúng.
\end{itemchoice}
}
\end{ex}

% ===== Câu 2 =====
% ID: L12C2B11-01-TL
% Mức: TH
\begin{ex}
% ID: L12C2B11-01-TL
% Mức: TH
Trong dạng “Bài toán thực tế dùng phương trình trạng thái”, Trình bày cơ sở lí thuyết của dạng “Bài toán thực tế dùng phương trình trạng thái” và nêu điều kiện áp dụng.
\choiceTF

\loigiai{
Nêu đúng đại lượng không đổi, hệ thức đặc trưng và yêu cầu dùng nhiệt độ kelvin, áp suất tuyệt đối khi có liên quan.
}
\end{ex}

% ===== Câu 3 =====
% ID: L12C2B11-01-TLN
% Mức: TH
\begin{ex}
% ID: L12C2B11-01-TLN
% Mức: TH
Trong dạng “Bài toán thực tế dùng phương trình trạng thái”, Có 1 mol khí ở 310 $K$ chiếm 3 lít. Tính áp suất theo kPa, lấy R=8,31.
\shortans{858,7}
\loigiai{
$p=nRT/V=858,7$ kPa.
}
\end{ex}

% ===== Câu 4 =====
% ID: L12C2B11-01-TN
% Mức: VD
\dangbt{Biến đổi trạng thái của một lượng khí xác định}
\begin{ex}
% ID: L12C2B11-01-TN
% Mức: VD
Nồi áp suất có van dạng lỗ tròn diện tích $S=1 \mathrm{cm}^2$, bị ép bởi lò xo độ cứng $k=1300 \mathrm{N/m}$ và \emph{bị nén một đoạn} $x$ (hằng số). Ban đầu ở áp suất khí quyển $p_0=1,013\times10^5 \mathrm{Pa}$ và $27^\circ\mathrm{C}$. Hỏi nhiệt độ để van vừa mở?
\choice
{\True $117^\circ\mathrm{C}$}
{$390^\circ\mathrm{C}$}
{$17^\circ\mathrm{C}$}
{$87^\circ\mathrm{C}$}
\loigiai{
Van mở khi áp suất trong nồi thắng được lực lò xo:
 $p_\text{mở}=p_0+\frac{kx}{S}.$
 Đẳng tích (xem nồi gần như kín, $V$ không đổi):
 $\frac{p_\text{mở}}{p_0}=\frac{T_\text{mở}}{T_0}\quad (T_0=300 \mathrm{K}).$
 Đề gốc thiếu số cụ thể của $x$. Với $\displaystyle x=0,00234 \mathrm{m}=0,234 \mathrm{cm}$ (giá trị chuẩn thường dùng để khớp đáp án),
  $\frac{kx}{S}=\frac{1300\cdot0,00234}{1\times10^{-4}}\approx3,0\times10^4\mathrm{Pa}\Rightarrow$ \frac{p_\text{mở}}{p_0} $\approx1,30.$ 
 Suy ra 
 $T_\text{mở}\approx1,30 T_0\approx1,30\cdot300=390 \mathrm{K}=117^\circ\mathrm{C}.$
}
\end{ex}

% ===== Câu 5 =====
% ID: L12C2B11-02-DS
% Mức: TH
\dangbt{Bài toán thực tế dùng phương trình trạng thái}
\begin{ex}
% ID: L12C2B11-02-DS
% Mức: TH
Xét các phát biểu sau trong dạng “Bài toán thực tế dùng phương trình trạng thái”.
\choiceTF
{\True Có 1 mol khí ở 290 $K$ chiếm 3 lít. Lấy R= $8,31\,\text{J}$ /(mol.K). Áp suất gần bằng — phát biểu đúng là: $803,3\,\text{kPa}$.}
{Có 1 mol khí ở 290 $K$ chiếm 3 lít. Lấy R= $8,31\,\text{J}$ /(mol.K). Áp suất gần bằng — phát biểu $80,33\,\text{kPa}$ là đúng.}
{\True Khí chuyển từ $p_1=105$ kPa, $V_1=5 lít,$ T_1=300 $K đến$ V_2=15 $lít,$ T_2=380 $K.$ p_2 $gần bằng — phát biểu đúng là: 44,33 kPa.$}
{Khí chuyển từ $p_1=105$ kPa, $V_1=5 lít,$ T_1=300 $K đến$ V_2=15 $lít,$ T_2=380 $K.$ p_2 $gần bằng — phát biểu 82,89 kPa là đúng.$}
\loigiai{
\begin{itemchoice}
\item (Đúng) Có 1 mol khí ở 290 $K$ chiếm 3 lít. Lấy R= $8,31\,\text{J}$ /(mol.K). Áp suất gần bằng — phát biểu đúng là: $803,3\,\text{kPa}$. (Vì): $p=nRT/V=803,3$ kPa. b) Sai: trái với hệ thức hoặc điều kiện đã nêu. c) Đúng: $p_2=p_1V_1T_2/(V_2T_1)=44,33$ kPa. d) Sai: trái với hệ thức hoặc điều kiện đã nêu
\item (Sai) Có 1 mol khí ở 290 $K$ chiếm 3 lít. Lấy R= $8,31\,\text{J}$ /(mol.K). Áp suất gần bằng — phát biểu $80,33\,\text{kPa}$ là đúng.
\item (Đúng) Khí chuyển từ $p_1=105$ kPa, $V_1=5 lít,$ T_1=300 $K đến$ V_2=15 $lít,$ T_2=380 $K.$ p_2 $gần bằng — phát biểu đúng là: 44,33 kPa.$.
\item (Sai) Khí chuyển từ $p_1=105$ kPa, $V_1=5 lít,$ T_1=300 $K đến$ V_2=15 $lít,$ T_2=380 $K.$ p_2 $gần bằng — phát biểu 82,89 kPa là đúng.$.
\end{itemchoice}
}
\end{ex}

% ===== Câu 6 =====
% ID: L12C2B11-02-TL
% Mức: TH
\begin{ex}
% ID: L12C2B11-02-TL
% Mức: TH
Trong dạng “Bài toán thực tế dùng phương trình trạng thái”, Mô tả dạng đồ thị đặc trưng và cách nhận biết quá trình trong dạng “Bài toán thực tế dùng phương trình trạng thái”.
\choiceTF

\loigiai{
Xác định các trục, đại lượng không đổi và suy ra hình dạng đồ thị từ hệ thức vật lí tương ứng.
}
\end{ex}

% ===== Câu 7 =====
% ID: L12C2B11-02-TLN
% Mức: TH
\begin{ex}
% ID: L12C2B11-02-TLN
% Mức: TH
Trong dạng “Bài toán thực tế dùng phương trình trạng thái”, Khí có $110\,\text{kPa}$, 5 lít ở 320 K. Đun đẳng tích đến 380 K. Tính áp suất sau theo kPa.
\shortans{130,63}
\loigiai{
$p_2=p_1T_2/T_1=130,63$ kPa.
}
\end{ex}

% ===== Câu 8 =====
% ID: L12C2B11-02-TN
% Mức: NB
\begin{ex}
% ID: L12C2B11-02-TN
% Mức: NB
Trong dạng “Bài toán thực tế dùng phương trình trạng thái”, Phương trình Clapeyron của khí lí tưởng là
\choice
{\True $pV=nRT$}
{$pV=RT/n$}
{$p/T=nV$}
{$p=nR/(VT)$}
\loigiai{
Đáp án A. Một lượng khí lí tưởng thỏa $pV=nRT$.
}
\end{ex}

% ===== Câu 9 =====
% ID: L12C2B11-03-DS
% Mức: VD
\begin{ex}
% ID: L12C2B11-03-DS
% Mức: VD
Xét các phát biểu sau trong dạng “Bài toán thực tế dùng phương trình trạng thái”.
\choiceTF
{\True Trong phương trình khí lí tưởng, nhiệt độ phải dùng — phát biểu đúng là: kelvin.}
{Trong phương trình khí lí tưởng, nhiệt độ phải dùng — phát biểu độ $C$ là đúng.}
{\True Ở thể tích không đổi, áp suất của lượng khí xác định tỉ lệ — phát biểu đúng là: thuận với nhiệt độ tuyệt đối.}
{Ở thể tích không đổi, áp suất của lượng khí xác định tỉ lệ — phát biểu thuận với bình phương nhiệt độ là đúng.}
\loigiai{
\begin{itemchoice}
\item (Đúng) Trong phương trình khí lí tưởng, nhiệt độ phải dùng — phát biểu đúng là: kelvin. (Vì): $T$ phải là nhiệt độ tuyệt đối. b) Sai: trái với hệ thức hoặc điều kiện đã nêu. c) Đúng: Từ $pV=nRT$, khi $V$ cố định thì $p/T$ không đổi. d) Sai: trái với hệ thức hoặc điều kiện đã nêu
\item (Sai) Trong phương trình khí lí tưởng, nhiệt độ phải dùng — phát biểu độ $C$ là đúng.
\item (Đúng) Ở thể tích không đổi, áp suất của lượng khí xác định tỉ lệ — phát biểu đúng là: thuận với nhiệt độ tuyệt đối.
\item (Sai) Ở thể tích không đổi, áp suất của lượng khí xác định tỉ lệ — phát biểu thuận với bình phương nhiệt độ là đúng.
\end{itemchoice}
}
\end{ex}

% ===== Câu 10 =====
% ID: L12C2B11-03-TL
% Mức: VD
\begin{bt}
% ID: L12C2B11-03-TL
% Mức: VD
Trong dạng “Bài toán thực tế dùng phương trình trạng thái”, Khí có $120\,\text{kPa}$, 7 lít ở 280 $K$; chuyển đến $480\,\text{kPa}$ và 340 K. Tính thể tích sau theo lít.
\loigiai{
$V_2=p_1V_1T_2/(p_2T_1)=2,13$ lít.
}
\end{bt}

% ===== Câu 11 =====
% ID: L12C2B11-03-TLN
% Mức: VD
\begin{ex}
% ID: L12C2B11-03-TLN
% Mức: VD
Trong dạng “Bài toán thực tế dùng phương trình trạng thái”, Khí có $120\,\text{kPa}$, 7 lít ở 280 $K$; chuyển đến $480\,\text{kPa}$ và 340 K. Tính thể tích sau theo lít.
\shortans{2,13}
\loigiai{
$V_2=p_1V_1T_2/(p_2T_1)=2,13$ lít.
}
\end{ex}

% ===== Câu 12 =====
% ID: L12C2B11-03-TN
% Mức: NB
\begin{ex}
% ID: L12C2B11-03-TN
% Mức: NB
Trong dạng “Bài toán thực tế dùng phương trình trạng thái”, Với lượng khí xác định, phương trình trạng thái giữa hai trạng thái là
\choice
{\True $p_1V_1/T_1=p_2V_2/T_2$}
{$p_1V_1T_1=p_2V_2T_2$}
{$p_1/T_1=p_2/T_2$ trong mọi quá trình}
{$V_1/T_1=V_2/T_2$ trong mọi quá trình}
\loigiai{
Đáp án A. Do $pV/T=nR$ không đổi.
}
\end{ex}

% ===== Câu 13 =====
% ID: L12C2B11-04-DS
% Mức: VD
\begin{ex}
% ID: L12C2B11-04-DS
% Mức: VD
Xét các phát biểu sau trong dạng “Bài toán thực tế dùng phương trình trạng thái”.
\choiceTF
{\True Số mol khí được tính từ khối lượng m và khối lượng mol $M$ bằng — phát biểu đúng là: $n=m/M$.}
{Số mol khí được tính từ khối lượng m và khối lượng mol $M$ bằng — phát biểu $n=mM$ là đúng.}
{\True Khối lượng riêng của khí lí tưởng thỏa — phát biểu đúng là: $\rho=pM/(RT)$.}
{Khối lượng riêng của khí lí tưởng thỏa — phát biểu $\rho=pRT/M$ là đúng.}
\loigiai{
\begin{itemchoice}
\item (Đúng) Số mol khí được tính từ khối lượng m và khối lượng mol $M$ bằng — phát biểu đúng là: $n=m/M$. (Vì): Theo định nghĩa số mol, $n=m/M$. b) Sai: trái với hệ thức hoặc điều kiện đã nêu. c) Đúng: Từ $pV=(m/M)RT$ và $\rho=m/V$. d) Sai: trái với hệ thức hoặc điều kiện đã nêu
\item (Sai) Số mol khí được tính từ khối lượng m và khối lượng mol $M$ bằng — phát biểu $n=mM$ là đúng.
\item (Đúng) Khối lượng riêng của khí lí tưởng thỏa — phát biểu đúng là: $\rho=pM/(RT)$.
\item (Sai) Khối lượng riêng của khí lí tưởng thỏa — phát biểu $\rho=pRT/M$ là đúng.
\end{itemchoice}
}
\end{ex}

% ===== Câu 14 =====
% ID: L12C2B11-04-TL
% Mức: VD
\begin{bt}
% ID: L12C2B11-04-TL
% Mức: VD
Trong dạng “Bài toán thực tế dùng phương trình trạng thái”, 1
\loigiai{
$n=m/M=28/28=1$ mol.
}
\end{bt}

% ===== Câu 15 =====
% ID: L12C2B11-04-TLN
% Mức: VD
\begin{ex}
% ID: L12C2B11-04-TLN
% Mức: VD
Trong dạng “Bài toán thực tế dùng phương trình trạng thái”, Một khí có khối lượng $28\,\text{g}$, khối lượng mol $28\,\text{g}$ /mol. Tính số mol.
\shortans{1}
\loigiai{
$n=m/M=28/28=1$ mol.
}
\end{ex}

% ===== Câu 16 =====
% ID: L12C2B11-04-TN
% Mức: TH
\begin{ex}
% ID: L12C2B11-04-TN
% Mức: TH
Trong dạng “Bài toán thực tế dùng phương trình trạng thái”, Có 1 mol khí ở 290 $K$ chiếm 3 lít. Lấy R= $8,31\,\text{J}$ /(mol.K). Áp suất gần bằng
\choice
{\True $803,3\,\text{kPa}$}
{$80,33\,\text{kPa}$}
{$8033\,\text{kPa}$}
{$0,03\,\text{kPa}$}
\loigiai{
Đáp án A. $p=nRT/V=803,3$ kPa.
}
\end{ex}

% ===== Câu 17 =====
% ID: L12C2B11-05-DS
% Mức: VDC
\begin{ex}
% ID: L12C2B11-05-DS
% Mức: VDC
Xét các phát biểu sau trong dạng “Bài toán thực tế dùng phương trình trạng thái”.
\choiceTF
{\True Hai bình cùng p, $V$, $T$ chứa khí lí tưởng khác loại thì có — phát biểu đúng là: cùng số phân tử.}
{Hai bình cùng p, $V$, $T$ chứa khí lí tưởng khác loại thì có — phát biểu cùng khối lượng là đúng.}
{\True Khi giải bài toán khí lí tưởng, áp suất đưa vào phương trình phải là — phát biểu đúng là: áp suất tuyệt đối.}
{Khi giải bài toán khí lí tưởng, áp suất đưa vào phương trình phải là — phát biểu áp suất do chất lỏng בלבד là đúng.}
\loigiai{
\begin{itemchoice}
\item (Đúng) Hai bình cùng p, $V$, $T$ chứa khí lí tưởng khác loại thì có — phát biểu đúng là: cùng số phân tử. (Vì): Theo $n=pV/(RT)$, số mol và số phân tử bằng nhau. b) Sai: trái với hệ thức hoặc điều kiện đã nêu. c) Đúng: Phương trình trạng thái sử dụng áp suất tuyệt đối. d) Sai: trái với hệ thức hoặc điều kiện đã nêu
\item (Sai) Hai bình cùng p, $V$, $T$ chứa khí lí tưởng khác loại thì có — phát biểu cùng khối lượng là đúng.
\item (Đúng) Khi giải bài toán khí lí tưởng, áp suất đưa vào phương trình phải là — phát biểu đúng là: áp suất tuyệt đối.
\item (Sai) Khi giải bài toán khí lí tưởng, áp suất đưa vào phương trình phải là — phát biểu áp suất do chất lỏng בלבד là đúng.
\end{itemchoice}
}
\end{ex}

% ===== Câu 18 =====
% ID: L12C2B11-05-TL
% Mức: VD
\begin{ex}
% ID: L12C2B11-05-TL
% Mức: VD
Trong dạng “Bài toán thực tế dùng phương trình trạng thái”, Từ $pV=nRT$, khi $T$, $V$ cố định thì $p\sim n$.
\choiceTF

\loigiai{
Từ $pV=nRT$, khi $T$, $V$ cố định thì $p\sim n$.
}
\end{ex}

% ===== Câu 19 =====
% ID: L12C2B11-05-TLN
% Mức: VD
\begin{ex}
% ID: L12C2B11-05-TLN
% Mức: VD
Trong dạng “Bài toán thực tế dùng phương trình trạng thái”, Ở cùng $T$ và $V$, áp suất tăng 3 lần. Số mol tăng bao nhiêu lần?
\shortans{3}
\loigiai{
Từ $pV=nRT$, khi $T$, $V$ cố định thì $p\sim n$.
}
\end{ex}

% ===== Câu 20 =====
% ID: L12C2B11-05-TN
% Mức: TH
\begin{ex}
% ID: L12C2B11-05-TN
% Mức: TH
Trong dạng “Bài toán thực tế dùng phương trình trạng thái”, Khí chuyển từ $p_1=105$ kPa, $V_1=5 lít,$ T_1=300 $K đến$ V_2=15 $lít,$ T_2=380 $K.$ p_2 $gần bằng$
\choice
{\True $44,33\,\text{kPa}$}
{$399\,\text{kPa}$}
{$82,89\,\text{kPa}$}
{$105\,\text{kPa}$}
\loigiai{
Đáp án A. $p_2=p_1V_1T_2/(V_2T_1)=44,33$ kPa.
}
\end{ex}

% ===== Câu 21 =====
% ID: L12C2B11-06-DS
% Mức: TH
\dangbt{Nhận biết phương trình trạng thái của khí lí tưởng}
\begin{ex}
% ID: L12C2B11-06-DS
% Mức: TH
Xét các phát biểu sau trong dạng “Nhận biết phương trình trạng thái của khí lí tưởng”.
\choiceTF
{\True Phương trình Clapeyron của khí lí tưởng là — phát biểu đúng là: $pV=nRT$.}
{Phương trình Clapeyron của khí lí tưởng là — phát biểu $pV=RT/n$ là đúng.}
{\True Với lượng khí xác định, phương trình trạng thái giữa hai trạng thái là — phát biểu đúng là: $p_1V_1/T_1=p_2V_2/T_2$.}
{Với lượng khí xác định, phương trình trạng thái giữa hai trạng thái là — phát biểu $p_1/T_1=p_2/T_2$ trong mọi quá trình là đúng.}
\loigiai{
\begin{itemchoice}
\item (Đúng) Phương trình Clapeyron của khí lí tưởng là — phát biểu đúng là: $pV=nRT$. (Vì): Một lượng khí lí tưởng thỏa $pV=nRT$. b) Sai: trái với hệ thức hoặc điều kiện đã nêu. c) Đúng: Do $pV/T=nR$ không đổi. d) Sai: trái với hệ thức hoặc điều kiện đã nêu
\item (Sai) Phương trình Clapeyron của khí lí tưởng là — phát biểu $pV=RT/n$ là đúng.
\item (Đúng) Với lượng khí xác định, phương trình trạng thái giữa hai trạng thái là — phát biểu đúng là: $p_1V_1/T_1=p_2V_2/T_2$.
\item (Sai) Với lượng khí xác định, phương trình trạng thái giữa hai trạng thái là — phát biểu $p_1/T_1=p_2/T_2$ trong mọi quá trình là đúng.
\end{itemchoice}
}
\end{ex}

% ===== Câu 22 =====
% ID: L12C2B11-06-TL
% Mức: VDC
\dangbt{Bài toán thực tế dùng phương trình trạng thái}
\begin{bt}
% ID: L12C2B11-06-TL
% Mức: VDC
Trong dạng “Bài toán thực tế dùng phương trình trạng thái”, undefined
\loigiai{
$V_2/V_1=(T_2/T_1)/(p_2/p_1)=1$.
}
\end{bt}

% ===== Câu 23 =====
% ID: L12C2B11-06-TLN
% Mức: VDC
\begin{ex}
% ID: L12C2B11-06-TLN
% Mức: VDC
Trong dạng “Bài toán thực tế dùng phương trình trạng thái”, Khí chuyển trạng thái sao cho p tăng 4 lần và $T$ tăng 4 lần. Thể tích thay đổi bao nhiêu lần?
\shortans{1}
\loigiai{
$V_2/V_1=(T_2/T_1)/(p_2/p_1)=1$.
}
\end{ex}

% ===== Câu 24 =====
% ID: L12C2B11-06-TN
% Mức: TH
\begin{ex}
% ID: L12C2B11-06-TN
% Mức: TH
Trong dạng “Bài toán thực tế dùng phương trình trạng thái”, Trong phương trình khí lí tưởng, nhiệt độ phải dùng
\choice
{\True kelvin}
{độ $C$}
{độ $F$}
{thang tùy ý}
\loigiai{
Đáp án A. $T$ phải là nhiệt độ tuyệt đối.
}
\end{ex}

% ===== Câu 25 =====
% ID: L12C2B11-07-DS
% Mức: TH
\dangbt{Nhận biết phương trình trạng thái của khí lí tưởng}
\begin{ex}
% ID: L12C2B11-07-DS
% Mức: TH
Xét các phát biểu sau trong dạng “Nhận biết phương trình trạng thái của khí lí tưởng”.
\choiceTF
{\True Có 4 mol khí ở 320 $K$ chiếm 5 lít. Lấy R= $8,31\,\text{J}$ /(mol.K). Áp suất gần bằng — phát biểu đúng là: $2127,36\,\text{kPa}$.}
{Có 4 mol khí ở 320 $K$ chiếm 5 lít. Lấy R= $8,31\,\text{J}$ /(mol.K). Áp suất gần bằng — phát biểu $212,74\,\text{kPa}$ là đúng.}
{\True Khí chuyển từ $p_1=125$ kPa, $V_1=7 lít,$ T_1=250 $K đến$ V_2=14 $lít,$ T_2=330 $K.$ p_2 $gần bằng — phát biểu đúng là: 82,5 kPa.$}
{Khí chuyển từ $p_1=125$ kPa, $V_1=7 lít,$ T_1=250 $K đến$ V_2=14 $lít,$ T_2=330 $K.$ p_2 $gần bằng — phát biểu 94,7 kPa là đúng.$}
\loigiai{
\begin{itemchoice}
\item (Đúng) Có 4 mol khí ở 320 $K$ chiếm 5 lít. Lấy R= $8,31\,\text{J}$ /(mol.K). Áp suất gần bằng — phát biểu đúng là: $2127,36\,\text{kPa}$. (Vì): $p=nRT/V=2127,36$ kPa. b) Sai: trái với hệ thức hoặc điều kiện đã nêu. c) Đúng: $p_2=p_1V_1T_2/(V_2T_1)=82,5$ kPa. d) Sai: trái với hệ thức hoặc điều kiện đã nêu
\item (Sai) Có 4 mol khí ở 320 $K$ chiếm 5 lít. Lấy R= $8,31\,\text{J}$ /(mol.K). Áp suất gần bằng — phát biểu $212,74\,\text{kPa}$ là đúng.
\item (Đúng) Khí chuyển từ $p_1=125$ kPa, $V_1=7 lít,$ T_1=250 $K đến$ V_2=14 $lít,$ T_2=330 $K.$ p_2 $gần bằng — phát biểu đúng là: 82,5 kPa.$.
\item (Sai) Khí chuyển từ $p_1=125$ kPa, $V_1=7 lít,$ T_1=250 $K đến$ V_2=14 $lít,$ T_2=330 $K.$ p_2 $gần bằng — phát biểu 94,7 kPa là đúng.$.
\end{itemchoice}
}
\end{ex}

% ===== Câu 26 =====
% ID: L12C2B11-07-TL
% Mức: TH
\begin{ex}
% ID: L12C2B11-07-TL
% Mức: TH
Trong dạng “Nhận biết phương trình trạng thái của khí lí tưởng”, Trình bày cơ sở lí thuyết của dạng “Nhận biết phương trình trạng thái của khí lí tưởng” và nêu điều kiện áp dụng.
\choiceTF

\loigiai{
Nêu đúng đại lượng không đổi, hệ thức đặc trưng và yêu cầu dùng nhiệt độ kelvin, áp suất tuyệt đối khi có liên quan.
}
\end{ex}

% ===== Câu 27 =====
% ID: L12C2B11-07-TLN
% Mức: TH
\begin{ex}
% ID: L12C2B11-07-TLN
% Mức: TH
Trong dạng “Nhận biết phương trình trạng thái của khí lí tưởng”, Có 2 mol khí ở 310 $K$ chiếm 4 lít. Tính áp suất theo kPa, lấy R=8,31.
\shortans{1288,05}
\loigiai{
$p=nRT/V=1288,05$ kPa.
}
\end{ex}

% ===== Câu 28 =====
% ID: L12C2B11-07-TN
% Mức: VD
\dangbt{Bài toán thực tế dùng phương trình trạng thái}
\begin{ex}
% ID: L12C2B11-07-TN
% Mức: VD
Trong dạng “Bài toán thực tế dùng phương trình trạng thái”, Ở thể tích không đổi, áp suất của lượng khí xác định tỉ lệ
\choice
{\True thuận với nhiệt độ tuyệt đối}
{nghịch với nhiệt độ}
{thuận với bình phương nhiệt độ}
{không phụ thuộc nhiệt độ}
\loigiai{
Đáp án A. Từ $pV=nRT$, khi $V$ cố định thì $p/T$ không đổi.
}
\end{ex}

% ===== Câu 29 =====
% ID: L12C2B11-08-DS
% Mức: VD
\dangbt{Nhận biết phương trình trạng thái của khí lí tưởng}
\begin{ex}
% ID: L12C2B11-08-DS
% Mức: VD
Xét các phát biểu sau trong dạng “Nhận biết phương trình trạng thái của khí lí tưởng”.
\choiceTF
{\True Trong phương trình khí lí tưởng, nhiệt độ phải dùng — phát biểu đúng là: kelvin.}
{Trong phương trình khí lí tưởng, nhiệt độ phải dùng — phát biểu độ $C$ là đúng.}
{\True Ở thể tích không đổi, áp suất của lượng khí xác định tỉ lệ — phát biểu đúng là: thuận với nhiệt độ tuyệt đối.}
{Ở thể tích không đổi, áp suất của lượng khí xác định tỉ lệ — phát biểu thuận với bình phương nhiệt độ là đúng.}
\loigiai{
\begin{itemchoice}
\item (Đúng) Trong phương trình khí lí tưởng, nhiệt độ phải dùng — phát biểu đúng là: kelvin. (Vì): $T$ phải là nhiệt độ tuyệt đối. b) Sai: trái với hệ thức hoặc điều kiện đã nêu. c) Đúng: Từ $pV=nRT$, khi $V$ cố định thì $p/T$ không đổi. d) Sai: trái với hệ thức hoặc điều kiện đã nêu
\item (Sai) Trong phương trình khí lí tưởng, nhiệt độ phải dùng — phát biểu độ $C$ là đúng.
\item (Đúng) Ở thể tích không đổi, áp suất của lượng khí xác định tỉ lệ — phát biểu đúng là: thuận với nhiệt độ tuyệt đối.
\item (Sai) Ở thể tích không đổi, áp suất của lượng khí xác định tỉ lệ — phát biểu thuận với bình phương nhiệt độ là đúng.
\end{itemchoice}
}
\end{ex}

% ===== Câu 30 =====
% ID: L12C2B11-08-TL
% Mức: TH
\begin{ex}
% ID: L12C2B11-08-TL
% Mức: TH
Trong dạng “Nhận biết phương trình trạng thái của khí lí tưởng”, Mô tả dạng đồ thị đặc trưng và cách nhận biết quá trình trong dạng “Nhận biết phương trình trạng thái của khí lí tưởng”.
\choiceTF

\loigiai{
Xác định các trục, đại lượng không đổi và suy ra hình dạng đồ thị từ hệ thức vật lí tương ứng.
}
\end{ex}

% ===== Câu 31 =====
% ID: L12C2B11-08-TLN
% Mức: TH
\begin{ex}
% ID: L12C2B11-08-TLN
% Mức: TH
Trong dạng “Nhận biết phương trình trạng thái của khí lí tưởng”, Khí có $120\,\text{kPa}$, 6 lít ở 320 K. Đun đẳng tích đến 380 K. Tính áp suất sau theo kPa.
\shortans{142,5}
\loigiai{
$p_2=p_1T_2/T_1=142,5$ kPa.
}
\end{ex}

% ===== Câu 32 =====
% ID: L12C2B11-08-TN
% Mức: VD
\dangbt{Bài toán thực tế dùng phương trình trạng thái}
\begin{ex}
% ID: L12C2B11-08-TN
% Mức: VD
Trong dạng “Bài toán thực tế dùng phương trình trạng thái”, Số mol khí được tính từ khối lượng m và khối lượng mol $M$ bằng
\choice
{\True $n=m/M$}
{$n=mM$}
{$n=M/m$}
{$n=m+M$}
\loigiai{
Đáp án A. Theo định nghĩa số mol, $n=m/M$.
}
\end{ex}

% ===== Câu 33 =====
% ID: L12C2B11-09-DS
% Mức: VD
\dangbt{Nhận biết phương trình trạng thái của khí lí tưởng}
\begin{ex}
% ID: L12C2B11-09-DS
% Mức: VD
Xét các phát biểu sau trong dạng “Nhận biết phương trình trạng thái của khí lí tưởng”.
\choiceTF
{\True Số mol khí được tính từ khối lượng m và khối lượng mol $M$ bằng — phát biểu đúng là: $n=m/M$.}
{Số mol khí được tính từ khối lượng m và khối lượng mol $M$ bằng — phát biểu $n=mM$ là đúng.}
{\True Khối lượng riêng của khí lí tưởng thỏa — phát biểu đúng là: $\rho=pM/(RT)$.}
{Khối lượng riêng của khí lí tưởng thỏa — phát biểu $\rho=pRT/M$ là đúng.}
\loigiai{
\begin{itemchoice}
\item (Đúng) Số mol khí được tính từ khối lượng m và khối lượng mol $M$ bằng — phát biểu đúng là: $n=m/M$. (Vì): Theo định nghĩa số mol, $n=m/M$. b) Sai: trái với hệ thức hoặc điều kiện đã nêu. c) Đúng: Từ $pV=(m/M)RT$ và $\rho=m/V$. d) Sai: trái với hệ thức hoặc điều kiện đã nêu
\item (Sai) Số mol khí được tính từ khối lượng m và khối lượng mol $M$ bằng — phát biểu $n=mM$ là đúng.
\item (Đúng) Khối lượng riêng của khí lí tưởng thỏa — phát biểu đúng là: $\rho=pM/(RT)$.
\item (Sai) Khối lượng riêng của khí lí tưởng thỏa — phát biểu $\rho=pRT/M$ là đúng.
\end{itemchoice}
}
\end{ex}

% ===== Câu 34 =====
% ID: L12C2B11-09-TL
% Mức: VD
\begin{bt}
% ID: L12C2B11-09-TL
% Mức: VD
Trong dạng “Nhận biết phương trình trạng thái của khí lí tưởng”, Khí có $130\,\text{kPa}$, 8 lít ở 280 $K$; chuyển đến $260\,\text{kPa}$ và 340 K. Tính thể tích sau theo lít.
\loigiai{
$V_2=p_1V_1T_2/(p_2T_1)=4,86$ lít.
}
\end{bt}

% ===== Câu 35 =====
% ID: L12C2B11-09-TLN
% Mức: VD
\begin{ex}
% ID: L12C2B11-09-TLN
% Mức: VD
Trong dạng “Nhận biết phương trình trạng thái của khí lí tưởng”, Khí có $130\,\text{kPa}$, 8 lít ở 280 $K$; chuyển đến $260\,\text{kPa}$ và 340 K. Tính thể tích sau theo lít.
\shortans{4,86}
\loigiai{
$V_2=p_1V_1T_2/(p_2T_1)=4,86$ lít.
}
\end{ex}

% ===== Câu 36 =====
% ID: L12C2B11-09-TN
% Mức: VD
\dangbt{Bài toán thực tế dùng phương trình trạng thái}
\begin{ex}
% ID: L12C2B11-09-TN
% Mức: VD
Trong dạng “Bài toán thực tế dùng phương trình trạng thái”, Khối lượng riêng của khí lí tưởng thỏa
\choice
{\True $\rho=pM/(RT)$}
{$\rho=RT/(pM)$}
{$\rho=pRT/M$}
{$\rho=p/(MRT)$}
\loigiai{
Đáp án A. Từ $pV=(m/M)RT$ và $\rho=m/V$.
}
\end{ex}

% ===== Câu 37 =====
% ID: L12C2B11-10-DS
% Mức: VDC
\dangbt{Nhận biết phương trình trạng thái của khí lí tưởng}
\begin{ex}
% ID: L12C2B11-10-DS
% Mức: VDC
Xét các phát biểu sau trong dạng “Nhận biết phương trình trạng thái của khí lí tưởng”.
\choiceTF
{\True Hai bình cùng p, $V$, $T$ chứa khí lí tưởng khác loại thì có — phát biểu đúng là: cùng số phân tử.}
{Hai bình cùng p, $V$, $T$ chứa khí lí tưởng khác loại thì có — phát biểu cùng khối lượng là đúng.}
{\True Khi giải bài toán khí lí tưởng, áp suất đưa vào phương trình phải là — phát biểu đúng là: áp suất tuyệt đối.}
{Khi giải bài toán khí lí tưởng, áp suất đưa vào phương trình phải là — phát biểu áp suất do chất lỏng בלבד là đúng.}
\loigiai{
\begin{itemchoice}
\item (Đúng) Hai bình cùng p, $V$, $T$ chứa khí lí tưởng khác loại thì có — phát biểu đúng là: cùng số phân tử. (Vì): Theo $n=pV/(RT)$, số mol và số phân tử bằng nhau. b) Sai: trái với hệ thức hoặc điều kiện đã nêu. c) Đúng: Phương trình trạng thái sử dụng áp suất tuyệt đối. d) Sai: trái với hệ thức hoặc điều kiện đã nêu
\item (Sai) Hai bình cùng p, $V$, $T$ chứa khí lí tưởng khác loại thì có — phát biểu cùng khối lượng là đúng.
\item (Đúng) Khi giải bài toán khí lí tưởng, áp suất đưa vào phương trình phải là — phát biểu đúng là: áp suất tuyệt đối.
\item (Sai) Khi giải bài toán khí lí tưởng, áp suất đưa vào phương trình phải là — phát biểu áp suất do chất lỏng בלבד là đúng.
\end{itemchoice}
}
\end{ex}

% ===== Câu 38 =====
% ID: L12C2B11-10-TL
% Mức: VD
\begin{bt}
% ID: L12C2B11-10-TL
% Mức: VD
Trong dạng “Nhận biết phương trình trạng thái của khí lí tưởng”, 2
\loigiai{
$n=m/M=56/28=2$ mol.
}
\end{bt}

% ===== Câu 39 =====
% ID: L12C2B11-10-TLN
% Mức: VD
\begin{ex}
% ID: L12C2B11-10-TLN
% Mức: VD
Trong dạng “Nhận biết phương trình trạng thái của khí lí tưởng”, Một khí có khối lượng $56\,\text{g}$, khối lượng mol $28\,\text{g}$ /mol. Tính số mol.
\shortans{2}
\loigiai{
$n=m/M=56/28=2$ mol.
}
\end{ex}

% ===== Câu 40 =====
% ID: L12C2B11-10-TN
% Mức: VDC
\dangbt{Bài toán thực tế dùng phương trình trạng thái}
\begin{ex}
% ID: L12C2B11-10-TN
% Mức: VDC
Trong dạng “Bài toán thực tế dùng phương trình trạng thái”, Hai bình cùng p, $V$, $T$ chứa khí lí tưởng khác loại thì có
\choice
{\True cùng số phân tử}
{cùng khối lượng}
{cùng khối lượng riêng}
{khác số mol}
\loigiai{
Đáp án A. Theo $n=pV/(RT)$, số mol và số phân tử bằng nhau.
}
\end{ex}

% ===== Câu 41 =====
% ID: L12C2B11-11-DS
% Mức: TH
\dangbt{Tính các thông số trạng thái p, V, T}
\begin{ex}
% ID: L12C2B11-11-DS
% Mức: TH
Xét các phát biểu sau trong dạng “Tính các thông số trạng thái p, $V$, T”.
\choiceTF
{\True Phương trình Clapeyron của khí lí tưởng là — phát biểu đúng là: $pV=nRT$.}
{Phương trình Clapeyron của khí lí tưởng là — phát biểu $pV=RT/n$ là đúng.}
{\True Với lượng khí xác định, phương trình trạng thái giữa hai trạng thái là — phát biểu đúng là: $p_1V_1/T_1=p_2V_2/T_2$.}
{Với lượng khí xác định, phương trình trạng thái giữa hai trạng thái là — phát biểu $p_1/T_1=p_2/T_2$ trong mọi quá trình là đúng.}
\loigiai{
\begin{itemchoice}
\item (Đúng) Phương trình Clapeyron của khí lí tưởng là — phát biểu đúng là: $pV=nRT$. (Vì): Một lượng khí lí tưởng thỏa $pV=nRT$. b) Sai: trái với hệ thức hoặc điều kiện đã nêu. c) Đúng: Do $pV/T=nR$ không đổi. d) Sai: trái với hệ thức hoặc điều kiện đã nêu
\item (Sai) Phương trình Clapeyron của khí lí tưởng là — phát biểu $pV=RT/n$ là đúng.
\item (Đúng) Với lượng khí xác định, phương trình trạng thái giữa hai trạng thái là — phát biểu đúng là: $p_1V_1/T_1=p_2V_2/T_2$.
\item (Sai) Với lượng khí xác định, phương trình trạng thái giữa hai trạng thái là — phát biểu $p_1/T_1=p_2/T_2$ trong mọi quá trình là đúng.
\end{itemchoice}
}
\end{ex}

% ===== Câu 42 =====
% ID: L12C2B11-11-TL
% Mức: VD
\dangbt{Nhận biết phương trình trạng thái của khí lí tưởng}
\begin{ex}
% ID: L12C2B11-11-TL
% Mức: VD
Trong dạng “Nhận biết phương trình trạng thái của khí lí tưởng”, Từ $pV=nRT$, khi $T$, $V$ cố định thì $p\sim n$.
\choiceTF

\loigiai{
Từ $pV=nRT$, khi $T$, $V$ cố định thì $p\sim n$.
}
\end{ex}

% ===== Câu 43 =====
% ID: L12C2B11-11-TLN
% Mức: VD
\begin{ex}
% ID: L12C2B11-11-TLN
% Mức: VD
Trong dạng “Nhận biết phương trình trạng thái của khí lí tưởng”, Ở cùng $T$ và $V$, áp suất tăng 4 lần. Số mol tăng bao nhiêu lần?
\shortans{4}
\loigiai{
Từ $pV=nRT$, khi $T$, $V$ cố định thì $p\sim n$.
}
\end{ex}

% ===== Câu 44 =====
% ID: L12C2B11-11-TN
% Mức: VDC
\dangbt{Bài toán thực tế dùng phương trình trạng thái}
\begin{ex}
% ID: L12C2B11-11-TN
% Mức: VDC
Trong dạng “Bài toán thực tế dùng phương trình trạng thái”, Khi giải bài toán khí lí tưởng, áp suất đưa vào phương trình phải là
\choice
{\True áp suất tuyệt đối}
{áp suất dư}
{áp suất do chất lỏng בלבד}
{áp suất tương đối tùy ý}
\loigiai{
Đáp án A. Phương trình trạng thái sử dụng áp suất tuyệt đối.
}
\end{ex}

% ===== Câu 45 =====
% ID: L12C2B11-12-DS
% Mức: TH
\dangbt{Tính các thông số trạng thái p, V, T}
\begin{ex}
% ID: L12C2B11-12-DS
% Mức: TH
Xét các phát biểu sau trong dạng “Tính các thông số trạng thái p, $V$, T”.
\choiceTF
{\True Có 1 mol khí ở 250 $K$ chiếm 6 lít. Lấy R= $8,31\,\text{J}$ /(mol.K). Áp suất gần bằng — phát biểu đúng là: $346,25\,\text{kPa}$.}
{Có 1 mol khí ở 250 $K$ chiếm 6 lít. Lấy R= $8,31\,\text{J}$ /(mol.K). Áp suất gần bằng — phát biểu $34,63\,\text{kPa}$ là đúng.}
{\True Khí chuyển từ $p_1=130$ kPa, $V_1=8 lít,$ T_1=260 $K đến$ V_2=24 $lít,$ T_2=340 $K.$ p_2 $gần bằng — phát biểu đúng là: 56,67 kPa.$}
{Khí chuyển từ $p_1=130$ kPa, $V_1=8 lít,$ T_1=260 $K đến$ V_2=24 $lít,$ T_2=340 $K.$ p_2 $gần bằng — phát biểu 99,41 kPa là đúng.$}
\loigiai{
\begin{itemchoice}
\item (Đúng) Có 1 mol khí ở 250 $K$ chiếm 6 lít. Lấy R= $8,31\,\text{J}$ /(mol.K). Áp suất gần bằng — phát biểu đúng là: $346,25\,\text{kPa}$. (Vì): $p=nRT/V=346,25$ kPa. b) Sai: trái với hệ thức hoặc điều kiện đã nêu. c) Đúng: $p_2=p_1V_1T_2/(V_2T_1)=56,67$ kPa. d) Sai: trái với hệ thức hoặc điều kiện đã nêu
\item (Sai) Có 1 mol khí ở 250 $K$ chiếm 6 lít. Lấy R= $8,31\,\text{J}$ /(mol.K). Áp suất gần bằng — phát biểu $34,63\,\text{kPa}$ là đúng.
\item (Đúng) Khí chuyển từ $p_1=130$ kPa, $V_1=8 lít,$ T_1=260 $K đến$ V_2=24 $lít,$ T_2=340 $K.$ p_2 $gần bằng — phát biểu đúng là: 56,67 kPa.$.
\item (Sai) Khí chuyển từ $p_1=130$ kPa, $V_1=8 lít,$ T_1=260 $K đến$ V_2=24 $lít,$ T_2=340 $K.$ p_2 $gần bằng — phát biểu 99,41 kPa là đúng.$.
\end{itemchoice}
}
\end{ex}

% ===== Câu 46 =====
% ID: L12C2B11-12-TL
% Mức: VDC
\dangbt{Nhận biết phương trình trạng thái của khí lí tưởng}
\begin{bt}
% ID: L12C2B11-12-TL
% Mức: VDC
Trong dạng “Nhận biết phương trình trạng thái của khí lí tưởng”, undefined
\loigiai{
$V_2/V_1=(T_2/T_1)/(p_2/p_1)=1$.
}
\end{bt}

% ===== Câu 47 =====
% ID: L12C2B11-12-TLN
% Mức: VDC
\begin{ex}
% ID: L12C2B11-12-TLN
% Mức: VDC
Trong dạng “Nhận biết phương trình trạng thái của khí lí tưởng”, Khí chuyển trạng thái sao cho p tăng 2 lần và $T$ tăng 2 lần. Thể tích thay đổi bao nhiêu lần?
\shortans{1}
\loigiai{
$V_2/V_1=(T_2/T_1)/(p_2/p_1)=1$.
}
\end{ex}

% ===== Câu 48 =====
% ID: L12C2B11-12-TN
% Mức: NB
\begin{ex}
% ID: L12C2B11-12-TN
% Mức: NB
Trong dạng “Nhận biết phương trình trạng thái của khí lí tưởng”, Phương trình Clapeyron của khí lí tưởng là
\choice
{\True $pV=nRT$}
{$pV=RT/n$}
{$p/T=nV$}
{$p=nR/(VT)$}
\loigiai{
Đáp án A. Một lượng khí lí tưởng thỏa $pV=nRT$.
}
\end{ex}

% ===== Câu 49 =====
% ID: L12C2B11-13-DS
% Mức: VD
\dangbt{Tính các thông số trạng thái p, V, T}
\begin{ex}
% ID: L12C2B11-13-DS
% Mức: VD
Xét các phát biểu sau trong dạng “Tính các thông số trạng thái p, $V$, T”.
\choiceTF
{\True Trong phương trình khí lí tưởng, nhiệt độ phải dùng — phát biểu đúng là: kelvin.}
{Trong phương trình khí lí tưởng, nhiệt độ phải dùng — phát biểu độ $C$ là đúng.}
{\True Ở thể tích không đổi, áp suất của lượng khí xác định tỉ lệ — phát biểu đúng là: thuận với nhiệt độ tuyệt đối.}
{Ở thể tích không đổi, áp suất của lượng khí xác định tỉ lệ — phát biểu thuận với bình phương nhiệt độ là đúng.}
\loigiai{
\begin{itemchoice}
\item (Đúng) Trong phương trình khí lí tưởng, nhiệt độ phải dùng — phát biểu đúng là: kelvin. (Vì): $T$ phải là nhiệt độ tuyệt đối. b) Sai: trái với hệ thức hoặc điều kiện đã nêu. c) Đúng: Từ $pV=nRT$, khi $V$ cố định thì $p/T$ không đổi. d) Sai: trái với hệ thức hoặc điều kiện đã nêu
\item (Sai) Trong phương trình khí lí tưởng, nhiệt độ phải dùng — phát biểu độ $C$ là đúng.
\item (Đúng) Ở thể tích không đổi, áp suất của lượng khí xác định tỉ lệ — phát biểu đúng là: thuận với nhiệt độ tuyệt đối.
\item (Sai) Ở thể tích không đổi, áp suất của lượng khí xác định tỉ lệ — phát biểu thuận với bình phương nhiệt độ là đúng.
\end{itemchoice}
}
\end{ex}

% ===== Câu 50 =====
% ID: L12C2B11-13-TL
% Mức: TH
\begin{ex}
% ID: L12C2B11-13-TL
% Mức: TH
Trong dạng “Tính các thông số trạng thái p, $V$, T”, Trình bày cơ sở lí thuyết của dạng “Tính các thông số trạng thái p, $V$, T” và nêu điều kiện áp dụng.
\choiceTF

\loigiai{
Nêu đúng đại lượng không đổi, hệ thức đặc trưng và yêu cầu dùng nhiệt độ kelvin, áp suất tuyệt đối khi có liên quan.
}
\end{ex}

% ===== Câu 51 =====
% ID: L12C2B11-13-TLN
% Mức: TH
\begin{ex}
% ID: L12C2B11-13-TLN
% Mức: TH
Trong dạng “Tính các thông số trạng thái p, $V$, T”, Có 3 mol khí ở 320 $K$ chiếm 5 lít. Tính áp suất theo kPa, lấy R=8,31.
\shortans{1595,52}
\loigiai{
$p=nRT/V=1595,52$ kPa.
}
\end{ex}

% ===== Câu 52 =====
% ID: L12C2B11-13-TN
% Mức: NB
\dangbt{Nhận biết phương trình trạng thái của khí lí tưởng}
\begin{ex}
% ID: L12C2B11-13-TN
% Mức: NB
Trong dạng “Nhận biết phương trình trạng thái của khí lí tưởng”, Với lượng khí xác định, phương trình trạng thái giữa hai trạng thái là
\choice
{\True $p_1V_1/T_1=p_2V_2/T_2$}
{$p_1V_1T_1=p_2V_2T_2$}
{$p_1/T_1=p_2/T_2$ trong mọi quá trình}
{$V_1/T_1=V_2/T_2$ trong mọi quá trình}
\loigiai{
Đáp án A. Do $pV/T=nR$ không đổi.
}
\end{ex}

% ===== Câu 53 =====
% ID: L12C2B11-14-DS
% Mức: VD
\dangbt{Tính các thông số trạng thái p, V, T}
\begin{ex}
% ID: L12C2B11-14-DS
% Mức: VD
Xét các phát biểu sau trong dạng “Tính các thông số trạng thái p, $V$, T”.
\choiceTF
{\True Số mol khí được tính từ khối lượng m và khối lượng mol $M$ bằng — phát biểu đúng là: $n=m/M$.}
{Số mol khí được tính từ khối lượng m và khối lượng mol $M$ bằng — phát biểu $n=mM$ là đúng.}
{\True Khối lượng riêng của khí lí tưởng thỏa — phát biểu đúng là: $\rho=pM/(RT)$.}
{Khối lượng riêng của khí lí tưởng thỏa — phát biểu $\rho=pRT/M$ là đúng.}
\loigiai{
\begin{itemchoice}
\item (Đúng) Số mol khí được tính từ khối lượng m và khối lượng mol $M$ bằng — phát biểu đúng là: $n=m/M$. (Vì): Theo định nghĩa số mol, $n=m/M$. b) Sai: trái với hệ thức hoặc điều kiện đã nêu. c) Đúng: Từ $pV=(m/M)RT$ và $\rho=m/V$. d) Sai: trái với hệ thức hoặc điều kiện đã nêu
\item (Sai) Số mol khí được tính từ khối lượng m và khối lượng mol $M$ bằng — phát biểu $n=mM$ là đúng.
\item (Đúng) Khối lượng riêng của khí lí tưởng thỏa — phát biểu đúng là: $\rho=pM/(RT)$.
\item (Sai) Khối lượng riêng của khí lí tưởng thỏa — phát biểu $\rho=pRT/M$ là đúng.
\end{itemchoice}
}
\end{ex}

% ===== Câu 54 =====
% ID: L12C2B11-14-TL
% Mức: TH
\begin{ex}
% ID: L12C2B11-14-TL
% Mức: TH
Trong dạng “Tính các thông số trạng thái p, $V$, T”, Mô tả dạng đồ thị đặc trưng và cách nhận biết quá trình trong dạng “Tính các thông số trạng thái p, $V$, T”.
\choiceTF

\loigiai{
Xác định các trục, đại lượng không đổi và suy ra hình dạng đồ thị từ hệ thức vật lí tương ứng.
}
\end{ex}

% ===== Câu 55 =====
% ID: L12C2B11-14-TLN
% Mức: TH
\begin{ex}
% ID: L12C2B11-14-TLN
% Mức: TH
Trong dạng “Tính các thông số trạng thái p, $V$, T”, Khí có $130\,\text{kPa}$, 7 lít ở 280 K. Đun đẳng tích đến 340 K. Tính áp suất sau theo kPa.
\shortans{157,86}
\loigiai{
$p_2=p_1T_2/T_1=157,86$ kPa.
}
\end{ex}

% ===== Câu 56 =====
% ID: L12C2B11-14-TN
% Mức: TH
\dangbt{Nhận biết phương trình trạng thái của khí lí tưởng}
\begin{ex}
% ID: L12C2B11-14-TN
% Mức: TH
Trong dạng “Nhận biết phương trình trạng thái của khí lí tưởng”, Có 4 mol khí ở 320 $K$ chiếm 5 lít. Lấy R= $8,31\,\text{J}$ /(mol.K). Áp suất gần bằng
\choice
{\True $2127,36\,\text{kPa}$}
{$212,74\,\text{kPa}$}
{$21273,6\,\text{kPa}$}
{$0,1\,\text{kPa}$}
\loigiai{
Đáp án A. $p=nRT/V=2127,36$ kPa.
}
\end{ex}

% ===== Câu 57 =====
% ID: L12C2B11-15-DS
% Mức: VDC
\dangbt{Tính các thông số trạng thái p, V, T}
\begin{ex}
% ID: L12C2B11-15-DS
% Mức: VDC
Xét các phát biểu sau trong dạng “Tính các thông số trạng thái p, $V$, T”.
\choiceTF
{\True Hai bình cùng p, $V$, $T$ chứa khí lí tưởng khác loại thì có — phát biểu đúng là: cùng số phân tử.}
{Hai bình cùng p, $V$, $T$ chứa khí lí tưởng khác loại thì có — phát biểu cùng khối lượng là đúng.}
{\True Khi giải bài toán khí lí tưởng, áp suất đưa vào phương trình phải là — phát biểu đúng là: áp suất tuyệt đối.}
{Khi giải bài toán khí lí tưởng, áp suất đưa vào phương trình phải là — phát biểu áp suất do chất lỏng בלבד là đúng.}
\loigiai{
\begin{itemchoice}
\item (Đúng) Hai bình cùng p, $V$, $T$ chứa khí lí tưởng khác loại thì có — phát biểu đúng là: cùng số phân tử. (Vì): Theo $n=pV/(RT)$, số mol và số phân tử bằng nhau. b) Sai: trái với hệ thức hoặc điều kiện đã nêu. c) Đúng: Phương trình trạng thái sử dụng áp suất tuyệt đối. d) Sai: trái với hệ thức hoặc điều kiện đã nêu
\item (Sai) Hai bình cùng p, $V$, $T$ chứa khí lí tưởng khác loại thì có — phát biểu cùng khối lượng là đúng.
\item (Đúng) Khi giải bài toán khí lí tưởng, áp suất đưa vào phương trình phải là — phát biểu đúng là: áp suất tuyệt đối.
\item (Sai) Khi giải bài toán khí lí tưởng, áp suất đưa vào phương trình phải là — phát biểu áp suất do chất lỏng בלבד là đúng.
\end{itemchoice}
}
\end{ex}

% ===== Câu 58 =====
% ID: L12C2B11-15-TL
% Mức: VD
\begin{bt}
% ID: L12C2B11-15-TL
% Mức: VD
Trong dạng “Tính các thông số trạng thái p, $V$, T”, Khí có $140\,\text{kPa}$, 3 lít ở 290 $K$; chuyển đến $420\,\text{kPa}$ và 350 K. Tính thể tích sau theo lít.
\loigiai{
$V_2=p_1V_1T_2/(p_2T_1)=1,21$ lít.
}
\end{bt}

% ===== Câu 59 =====
% ID: L12C2B11-15-TLN
% Mức: VD
\begin{ex}
% ID: L12C2B11-15-TLN
% Mức: VD
Trong dạng “Tính các thông số trạng thái p, $V$, T”, Khí có $140\,\text{kPa}$, 3 lít ở 290 $K$; chuyển đến $420\,\text{kPa}$ và 350 K. Tính thể tích sau theo lít.
\shortans{1,21}
\loigiai{
$V_2=p_1V_1T_2/(p_2T_1)=1,21$ lít.
}
\end{ex}

% ===== Câu 60 =====
% ID: L12C2B11-15-TN
% Mức: TH
\dangbt{Nhận biết phương trình trạng thái của khí lí tưởng}
\begin{ex}
% ID: L12C2B11-15-TN
% Mức: TH
Trong dạng “Nhận biết phương trình trạng thái của khí lí tưởng”, Khí chuyển từ $p_1=125$ kPa, $V_1=7 lít,$ T_1=250 $K đến$ V_2=14 $lít,$ T_2=330 $K.$ p_2 $gần bằng$
\choice
{\True $82,5\,\text{kPa}$}
{$330\,\text{kPa}$}
{$94,7\,\text{kPa}$}
{$125\,\text{kPa}$}
\loigiai{
Đáp án A. $p_2=p_1V_1T_2/(V_2T_1)=82,5$ kPa.
}
\end{ex}

% ===== Câu 61 =====
% ID: L12C2B11-16-DS
% Mức: TH
\dangbt{Xác định số mol, khối lượng và khối lượng riêng của khí}
\begin{ex}
% ID: L12C2B11-16-DS
% Mức: TH
Xét các phát biểu sau trong dạng “Xác định số mol, khối lượng và khối lượng riêng của khí”.
\choiceTF
{\True Phương trình Clapeyron của khí lí tưởng là — phát biểu đúng là: $pV=nRT$.}
{Phương trình Clapeyron của khí lí tưởng là — phát biểu $pV=RT/n$ là đúng.}
{\True Với lượng khí xác định, phương trình trạng thái giữa hai trạng thái là — phát biểu đúng là: $p_1V_1/T_1=p_2V_2/T_2$.}
{Với lượng khí xác định, phương trình trạng thái giữa hai trạng thái là — phát biểu $p_1/T_1=p_2/T_2$ trong mọi quá trình là đúng.}
\loigiai{
\begin{itemchoice}
\item (Đúng) Phương trình Clapeyron của khí lí tưởng là — phát biểu đúng là: $pV=nRT$. (Vì): Một lượng khí lí tưởng thỏa $pV=nRT$. b) Sai: trái với hệ thức hoặc điều kiện đã nêu. c) Đúng: Do $pV/T=nR$ không đổi. d) Sai: trái với hệ thức hoặc điều kiện đã nêu
\item (Sai) Phương trình Clapeyron của khí lí tưởng là — phát biểu $pV=RT/n$ là đúng.
\item (Đúng) Với lượng khí xác định, phương trình trạng thái giữa hai trạng thái là — phát biểu đúng là: $p_1V_1/T_1=p_2V_2/T_2$.
\item (Sai) Với lượng khí xác định, phương trình trạng thái giữa hai trạng thái là — phát biểu $p_1/T_1=p_2/T_2$ trong mọi quá trình là đúng.
\end{itemchoice}
}
\end{ex}

% ===== Câu 62 =====
% ID: L12C2B11-16-TL
% Mức: VD
\dangbt{Tính các thông số trạng thái p, V, T}
\begin{bt}
% ID: L12C2B11-16-TL
% Mức: VD
Trong dạng “Tính các thông số trạng thái p, $V$, T”, 3
\loigiai{
$n=m/M=84/28=3$ mol.
}
\end{bt}

% ===== Câu 63 =====
% ID: L12C2B11-16-TLN
% Mức: VD
\begin{ex}
% ID: L12C2B11-16-TLN
% Mức: VD
Trong dạng “Tính các thông số trạng thái p, $V$, T”, Một khí có khối lượng $84\,\text{g}$, khối lượng mol $28\,\text{g}$ /mol. Tính số mol.
\shortans{3}
\loigiai{
$n=m/M=84/28=3$ mol.
}
\end{ex}

% ===== Câu 64 =====
% ID: L12C2B11-16-TN
% Mức: TH
\dangbt{Nhận biết phương trình trạng thái của khí lí tưởng}
\begin{ex}
% ID: L12C2B11-16-TN
% Mức: TH
Trong dạng “Nhận biết phương trình trạng thái của khí lí tưởng”, Trong phương trình khí lí tưởng, nhiệt độ phải dùng
\choice
{\True kelvin}
{độ $C$}
{độ $F$}
{thang tùy ý}
\loigiai{
Đáp án A. $T$ phải là nhiệt độ tuyệt đối.
}
\end{ex}

% ===== Câu 65 =====
% ID: L12C2B11-17-DS
% Mức: TH
\dangbt{Xác định số mol, khối lượng và khối lượng riêng của khí}
\begin{ex}
% ID: L12C2B11-17-DS
% Mức: TH
Xét các phát biểu sau trong dạng “Xác định số mol, khối lượng và khối lượng riêng của khí”.
\choiceTF
{\True Có 3 mol khí ở 270 $K$ chiếm 8 lít. Lấy R= $8,31\,\text{J}$ /(mol.K). Áp suất gần bằng — phát biểu đúng là: $841,39\,\text{kPa}$.}
{Có 3 mol khí ở 270 $K$ chiếm 8 lít. Lấy R= $8,31\,\text{J}$ /(mol.K). Áp suất gần bằng — phát biểu $84,14\,\text{kPa}$ là đúng.}
{\True Khí chuyển từ $p_1=95$ kPa, $V_1=3 lít,$ T_1=280 $K đến$ V_2=15 $lít,$ T_2=360 $K.$ p_2 $gần bằng — phát biểu đúng là: 24,43 kPa.$}
{Khí chuyển từ $p_1=95$ kPa, $V_1=3 lít,$ T_1=280 $K đến$ V_2=15 $lít,$ T_2=360 $K.$ p_2 $gần bằng — phát biểu 73,89 kPa là đúng.$}
\loigiai{
\begin{itemchoice}
\item (Đúng) Có 3 mol khí ở 270 $K$ chiếm 8 lít. Lấy R= $8,31\,\text{J}$ /(mol.K). Áp suất gần bằng — phát biểu đúng là: $841,39\,\text{kPa}$. (Vì): $p=nRT/V=841,39$ kPa. b) Sai: trái với hệ thức hoặc điều kiện đã nêu. c) Đúng: $p_2=p_1V_1T_2/(V_2T_1)=24,43$ kPa. d) Sai: trái với hệ thức hoặc điều kiện đã nêu
\item (Sai) Có 3 mol khí ở 270 $K$ chiếm 8 lít. Lấy R= $8,31\,\text{J}$ /(mol.K). Áp suất gần bằng — phát biểu $84,14\,\text{kPa}$ là đúng.
\item (Đúng) Khí chuyển từ $p_1=95$ kPa, $V_1=3 lít,$ T_1=280 $K đến$ V_2=15 $lít,$ T_2=360 $K.$ p_2 $gần bằng — phát biểu đúng là: 24,43 kPa.$.
\item (Sai) Khí chuyển từ $p_1=95$ kPa, $V_1=3 lít,$ T_1=280 $K đến$ V_2=15 $lít,$ T_2=360 $K.$ p_2 $gần bằng — phát biểu 73,89 kPa là đúng.$.
\end{itemchoice}
}
\end{ex}

% ===== Câu 66 =====
% ID: L12C2B11-17-TL
% Mức: VD
\dangbt{Tính các thông số trạng thái p, V, T}
\begin{ex}
% ID: L12C2B11-17-TL
% Mức: VD
Trong dạng “Tính các thông số trạng thái p, $V$, T”, Từ $pV=nRT$, khi $T$, $V$ cố định thì $p\sim n$.
\choiceTF

\loigiai{
Từ $pV=nRT$, khi $T$, $V$ cố định thì $p\sim n$.
}
\end{ex}

% ===== Câu 67 =====
% ID: L12C2B11-17-TLN
% Mức: VD
\begin{ex}
% ID: L12C2B11-17-TLN
% Mức: VD
Trong dạng “Tính các thông số trạng thái p, $V$, T”, Ở cùng $T$ và $V$, áp suất tăng 2 lần. Số mol tăng bao nhiêu lần?
\shortans{2}
\loigiai{
Từ $pV=nRT$, khi $T$, $V$ cố định thì $p\sim n$.
}
\end{ex}

% ===== Câu 68 =====
% ID: L12C2B11-17-TN
% Mức: VD
\dangbt{Nhận biết phương trình trạng thái của khí lí tưởng}
\begin{ex}
% ID: L12C2B11-17-TN
% Mức: VD
Trong dạng “Nhận biết phương trình trạng thái của khí lí tưởng”, Ở thể tích không đổi, áp suất của lượng khí xác định tỉ lệ
\choice
{\True thuận với nhiệt độ tuyệt đối}
{nghịch với nhiệt độ}
{thuận với bình phương nhiệt độ}
{không phụ thuộc nhiệt độ}
\loigiai{
Đáp án A. Từ $pV=nRT$, khi $V$ cố định thì $p/T$ không đổi.
}
\end{ex}

% ===== Câu 69 =====
% ID: L12C2B11-18-DS
% Mức: VD
\dangbt{Xác định số mol, khối lượng và khối lượng riêng của khí}
\begin{ex}
% ID: L12C2B11-18-DS
% Mức: VD
Xét các phát biểu sau trong dạng “Xác định số mol, khối lượng và khối lượng riêng của khí”.
\choiceTF
{\True Trong phương trình khí lí tưởng, nhiệt độ phải dùng — phát biểu đúng là: kelvin.}
{Trong phương trình khí lí tưởng, nhiệt độ phải dùng — phát biểu độ $C$ là đúng.}
{\True Ở thể tích không đổi, áp suất của lượng khí xác định tỉ lệ — phát biểu đúng là: thuận với nhiệt độ tuyệt đối.}
{Ở thể tích không đổi, áp suất của lượng khí xác định tỉ lệ — phát biểu thuận với bình phương nhiệt độ là đúng.}
\loigiai{
\begin{itemchoice}
\item (Đúng) Trong phương trình khí lí tưởng, nhiệt độ phải dùng — phát biểu đúng là: kelvin. (Vì): $T$ phải là nhiệt độ tuyệt đối. b) Sai: trái với hệ thức hoặc điều kiện đã nêu. c) Đúng: Từ $pV=nRT$, khi $V$ cố định thì $p/T$ không đổi. d) Sai: trái với hệ thức hoặc điều kiện đã nêu
\item (Sai) Trong phương trình khí lí tưởng, nhiệt độ phải dùng — phát biểu độ $C$ là đúng.
\item (Đúng) Ở thể tích không đổi, áp suất của lượng khí xác định tỉ lệ — phát biểu đúng là: thuận với nhiệt độ tuyệt đối.
\item (Sai) Ở thể tích không đổi, áp suất của lượng khí xác định tỉ lệ — phát biểu thuận với bình phương nhiệt độ là đúng.
\end{itemchoice}
}
\end{ex}

% ===== Câu 70 =====
% ID: L12C2B11-18-TL
% Mức: VDC
\dangbt{Tính các thông số trạng thái p, V, T}
\begin{bt}
% ID: L12C2B11-18-TL
% Mức: VDC
Trong dạng “Tính các thông số trạng thái p, $V$, T”, undefined
\loigiai{
$V_2/V_1=(T_2/T_1)/(p_2/p_1)=1$.
}
\end{bt}

% ===== Câu 71 =====
% ID: L12C2B11-18-TLN
% Mức: VDC
\begin{ex}
% ID: L12C2B11-18-TLN
% Mức: VDC
Trong dạng “Tính các thông số trạng thái p, $V$, T”, Khí chuyển trạng thái sao cho p tăng 3 lần và $T$ tăng 3 lần. Thể tích thay đổi bao nhiêu lần?
\shortans{1}
\loigiai{
$V_2/V_1=(T_2/T_1)/(p_2/p_1)=1$.
}
\end{ex}

% ===== Câu 72 =====
% ID: L12C2B11-18-TN
% Mức: VD
\dangbt{Nhận biết phương trình trạng thái của khí lí tưởng}
\begin{ex}
% ID: L12C2B11-18-TN
% Mức: VD
Trong dạng “Nhận biết phương trình trạng thái của khí lí tưởng”, Số mol khí được tính từ khối lượng m và khối lượng mol $M$ bằng
\choice
{\True $n=m/M$}
{$n=mM$}
{$n=M/m$}
{$n=m+M$}
\loigiai{
Đáp án A. Theo định nghĩa số mol, $n=m/M$.
}
\end{ex}

% ===== Câu 73 =====
% ID: L12C2B11-19-DS
% Mức: VD
\dangbt{Xác định số mol, khối lượng và khối lượng riêng của khí}
\begin{ex}
% ID: L12C2B11-19-DS
% Mức: VD
Xét các phát biểu sau trong dạng “Xác định số mol, khối lượng và khối lượng riêng của khí”.
\choiceTF
{\True Số mol khí được tính từ khối lượng m và khối lượng mol $M$ bằng — phát biểu đúng là: $n=m/M$.}
{Số mol khí được tính từ khối lượng m và khối lượng mol $M$ bằng — phát biểu $n=mM$ là đúng.}
{\True Khối lượng riêng của khí lí tưởng thỏa — phát biểu đúng là: $\rho=pM/(RT)$.}
{Khối lượng riêng của khí lí tưởng thỏa — phát biểu $\rho=pRT/M$ là đúng.}
\loigiai{
\begin{itemchoice}
\item (Đúng) Số mol khí được tính từ khối lượng m và khối lượng mol $M$ bằng — phát biểu đúng là: $n=m/M$. (Vì): Theo định nghĩa số mol, $n=m/M$. b) Sai: trái với hệ thức hoặc điều kiện đã nêu. c) Đúng: Từ $pV=(m/M)RT$ và $\rho=m/V$. d) Sai: trái với hệ thức hoặc điều kiện đã nêu
\item (Sai) Số mol khí được tính từ khối lượng m và khối lượng mol $M$ bằng — phát biểu $n=mM$ là đúng.
\item (Đúng) Khối lượng riêng của khí lí tưởng thỏa — phát biểu đúng là: $\rho=pM/(RT)$.
\item (Sai) Khối lượng riêng của khí lí tưởng thỏa — phát biểu $\rho=pRT/M$ là đúng.
\end{itemchoice}
}
\end{ex}

% ===== Câu 74 =====
% ID: L12C2B11-19-TL
% Mức: TH
\begin{ex}
% ID: L12C2B11-19-TL
% Mức: TH
Trong dạng “Xác định số mol, khối lượng và khối lượng riêng của khí”, Trình bày cơ sở lí thuyết của dạng “Xác định số mol, khối lượng và khối lượng riêng của khí” và nêu điều kiện áp dụng.
\choiceTF

\loigiai{
Nêu đúng đại lượng không đổi, hệ thức đặc trưng và yêu cầu dùng nhiệt độ kelvin, áp suất tuyệt đối khi có liên quan.
}
\end{ex}

% ===== Câu 75 =====
% ID: L12C2B11-19-TLN
% Mức: TH
\begin{ex}
% ID: L12C2B11-19-TLN
% Mức: TH
Trong dạng “Xác định số mol, khối lượng và khối lượng riêng của khí”, Có 2 mol khí ở 290 $K$ chiếm 7 lít. Tính áp suất theo kPa, lấy R=8,31.
\shortans{688,54}
\loigiai{
$p=nRT/V=688,54$ kPa.
}
\end{ex}

% ===== Câu 76 =====
% ID: L12C2B11-19-TN
% Mức: VD
\dangbt{Nhận biết phương trình trạng thái của khí lí tưởng}
\begin{ex}
% ID: L12C2B11-19-TN
% Mức: VD
Trong dạng “Nhận biết phương trình trạng thái của khí lí tưởng”, Khối lượng riêng của khí lí tưởng thỏa
\choice
{\True $\rho=pM/(RT)$}
{$\rho=RT/(pM)$}
{$\rho=pRT/M$}
{$\rho=p/(MRT)$}
\loigiai{
Đáp án A. Từ $pV=(m/M)RT$ và $\rho=m/V$.
}
\end{ex}

% ===== Câu 77 =====
% ID: L12C2B11-20-DS
% Mức: VDC
\dangbt{Xác định số mol, khối lượng và khối lượng riêng của khí}
\begin{ex}
% ID: L12C2B11-20-DS
% Mức: VDC
Xét các phát biểu sau trong dạng “Xác định số mol, khối lượng và khối lượng riêng của khí”.
\choiceTF
{\True Hai bình cùng p, $V$, $T$ chứa khí lí tưởng khác loại thì có — phát biểu đúng là: cùng số phân tử.}
{Hai bình cùng p, $V$, $T$ chứa khí lí tưởng khác loại thì có — phát biểu cùng khối lượng là đúng.}
{\True Khi giải bài toán khí lí tưởng, áp suất đưa vào phương trình phải là — phát biểu đúng là: áp suất tuyệt đối.}
{Khi giải bài toán khí lí tưởng, áp suất đưa vào phương trình phải là — phát biểu áp suất do chất lỏng בלבד là đúng.}
\loigiai{
\begin{itemchoice}
\item (Đúng) Hai bình cùng p, $V$, $T$ chứa khí lí tưởng khác loại thì có — phát biểu đúng là: cùng số phân tử. (Vì): Theo $n=pV/(RT)$, số mol và số phân tử bằng nhau. b) Sai: trái với hệ thức hoặc điều kiện đã nêu. c) Đúng: Phương trình trạng thái sử dụng áp suất tuyệt đối. d) Sai: trái với hệ thức hoặc điều kiện đã nêu
\item (Sai) Hai bình cùng p, $V$, $T$ chứa khí lí tưởng khác loại thì có — phát biểu cùng khối lượng là đúng.
\item (Đúng) Khi giải bài toán khí lí tưởng, áp suất đưa vào phương trình phải là — phát biểu đúng là: áp suất tuyệt đối.
\item (Sai) Khi giải bài toán khí lí tưởng, áp suất đưa vào phương trình phải là — phát biểu áp suất do chất lỏng בלבד là đúng.
\end{itemchoice}
}
\end{ex}

% ===== Câu 78 =====
% ID: L12C2B11-20-TL
% Mức: TH
\begin{ex}
% ID: L12C2B11-20-TL
% Mức: TH
Trong dạng “Xác định số mol, khối lượng và khối lượng riêng của khí”, Mô tả dạng đồ thị đặc trưng và cách nhận biết quá trình trong dạng “Xác định số mol, khối lượng và khối lượng riêng của khí”.
\choiceTF

\loigiai{
Xác định các trục, đại lượng không đổi và suy ra hình dạng đồ thị từ hệ thức vật lí tương ứng.
}
\end{ex}

% ===== Câu 79 =====
% ID: L12C2B11-20-TLN
% Mức: TH
\begin{ex}
% ID: L12C2B11-20-TLN
% Mức: TH
Trong dạng “Xác định số mol, khối lượng và khối lượng riêng của khí”, Khí có $150\,\text{kPa}$, 3 lít ở 300 K. Đun đẳng tích đến 360 K. Tính áp suất sau theo kPa.
\shortans{180}
\loigiai{
$p_2=p_1T_2/T_1=180$ kPa.
}
\end{ex}

% ===== Câu 80 =====
% ID: L12C2B11-20-TN
% Mức: VDC
\dangbt{Nhận biết phương trình trạng thái của khí lí tưởng}
\begin{ex}
% ID: L12C2B11-20-TN
% Mức: VDC
Trong dạng “Nhận biết phương trình trạng thái của khí lí tưởng”, Hai bình cùng p, $V$, $T$ chứa khí lí tưởng khác loại thì có
\choice
{\True cùng số phân tử}
{cùng khối lượng}
{cùng khối lượng riêng}
{khác số mol}
\loigiai{
Đáp án A. Theo $n=pV/(RT)$, số mol và số phân tử bằng nhau.
}
\end{ex}

% ===== Câu 81 =====
% ID: L12C2B11-21-DS
% Mức: VDC
\dangbt{Xác định sự biến thiên khối lượng riêng của chất khí}
\begin{ex}
% ID: L12C2B11-21-DS
% Mức: VDC
Một bình thép chứa khí oxygen ($\mathrm{O_2}$) có dung tích $16,62\ \mathrm{l}$ ở nhiệt độ $27\ ^\circ\mathrm{C}$ và áp suất $3,0 \cdot 10^5\ \mathrm{Pa}$. Cho biết hằng số khí lí tưởng $R \approx 8,31\ \mathrm{J/(mol\cdot K)}$ và khối lượng mol của oxygen là $32\ \mathrm{g/mol}$. Nhận định về các đại lượng đặc trưng của khối khí dưới đây là đúng hay sai?
\choiceTF
{Trong hệ đơn vị SI, các đại lượng trong phương trình Clapeyron có đơn vị tương ứng là: áp suất $p$ đo bằng $\mathrm{atm}$, thể tích $V$ đo bằng lít ($\mathrm{l}$), nhiệt độ $T$ đo bằng kelvin ($\mathrm{K}$).}
{\True Số mol khí oxygen chứa trong bình ban đầu có giá trị bằng $2,0\ \mathrm{mol}$.}
{\True Khối lượng khí oxygen chứa trong bình ban đầu bằng $64\ \mathrm{g}$.}
{\True Nếu đun nóng bình khí đến nhiệt độ $127\ ^\circ\mathrm{C}$ và mở van xả bớt khí để giữ cho áp suất trong bình không đổi ở giá trị $3,0 \cdot 10^5\ \mathrm{Pa}$ thì khối lượng khí oxygen đã thoát ra ngoài là $16\ \mathrm{g}$.}
\loigiai{
\begin{itemchoice}
\item (Sai) Trong hệ đơn vị SI, các đại lượng trong phương trình Clapeyron có đơn vị tương ứng là: áp suất $p$ đo bằng $\mathrm{atm}$, thể tích $V$ đo bằng lít ($\mathrm{l}$), nhiệt độ $T$ đo bằng kelvin ($\mathrm{K}$). (Vì): Trong hệ đơn vị chuẩn SI, áp suất $p$ phải đo bằng pascal ($\mathrm{Pa}$), thể tích $V$ đo bằng mét khối ($\mathrm{m^3}$), hằng số khí $R$ đo bằng $\mathrm{J/(mol\cdot K)}$. Mệnh đề sai
\item (Đúng) Số mol khí oxygen chứa trong bình ban đầu có giá trị bằng $2,0\ \mathrm{mol}$. (Vì): ổi thể tích $V = 16,62\ \mathrm{l} = 16,62 \cdot 10^{-3}\ \mathrm{m^3}$ và nhiệt độ $T_1 = 27 + 273 = 300\ \mathrm{K}$. Số mol khí oxygen trong bình là:
\item (Đúng) Khối lượng khí oxygen chứa trong bình ban đầu bằng $64\ \mathrm{g}$. (Vì): Khối lượng khí oxygen ban đầu chứa trong bình là:
\item (Đúng) Nếu đun nóng bình khí đến nhiệt độ $127\ ^\circ\mathrm{C}$ và mở van xả bớt khí để giữ cho áp suất trong bình không đổi ở giá trị $3,0 \cdot 10^5\ \mathrm{Pa}$ thì khối lượng khí oxygen đã thoát ra ngoài là $16\ \mathrm{g}$. (Vì): Khi đun nóng bình đến $127\ ^\circ\mathrm{C}$ ($T_2 = 400\ \mathrm{K}$) và mở van để giữ áp suất không đổi, số mol khí còn lại trong bình là:
\end{itemchoice}
}
\end{ex}

% ===== Câu 82 =====
% ID: L12C2B11-21-TL
% Mức: VD
\dangbt{Xác định số mol, khối lượng và khối lượng riêng của khí}
\begin{bt}
% ID: L12C2B11-21-TL
% Mức: VD
Trong dạng “Xác định số mol, khối lượng và khối lượng riêng của khí”, Khí có $100\,\text{kPa}$, 5 lít ở 310 $K$; chuyển đến $200\,\text{kPa}$ và 370 K. Tính thể tích sau theo lít.
\loigiai{
$V_2=p_1V_1T_2/(p_2T_1)=2,98$ lít.
}
\end{bt}

% ===== Câu 83 =====
% ID: L12C2B11-21-TLN
% Mức: VD
\begin{ex}
% ID: L12C2B11-21-TLN
% Mức: VD
Trong dạng “Xác định số mol, khối lượng và khối lượng riêng của khí”, Khí có $100\,\text{kPa}$, 5 lít ở 310 $K$; chuyển đến $200\,\text{kPa}$ và 370 K. Tính thể tích sau theo lít.
\shortans{2,98}
\loigiai{
$V_2=p_1V_1T_2/(p_2T_1)=2,98$ lít.
}
\end{ex}

% ===== Câu 84 =====
% ID: L12C2B11-21-TN
% Mức: VDC
\dangbt{Nhận biết phương trình trạng thái của khí lí tưởng}
\begin{ex}
% ID: L12C2B11-21-TN
% Mức: VDC
Trong dạng “Nhận biết phương trình trạng thái của khí lí tưởng”, Khi giải bài toán khí lí tưởng, áp suất đưa vào phương trình phải là
\choice
{\True áp suất tuyệt đối}
{áp suất dư}
{áp suất do chất lỏng בלבד}
{áp suất tương đối tùy ý}
\loigiai{
Đáp án A. Phương trình trạng thái sử dụng áp suất tuyệt đối.
}
\end{ex}

% ===== Câu 85 =====
% ID: L12C2B11-22-DS
% Mức: VDC
\dangbt{Xác định sự biến thiên khối lượng riêng của chất khí}
\begin{ex}
% ID: L12C2B11-22-DS
% Mức: VDC
Có ba chai thủy tinh giống nhau có dung tích bằng nhau đựng các chất khí khác nhau (helium $\mathrm{He}$ có $M = 4\ \mathrm{g/mol}$, butane $\mathrm{C_4H_{10}}$ có $M = 58\ \mathrm{g/mol}$ và carbon dioxide $\mathrm{CO_2}$ có $M = 44\ \mathrm{g/mol}$) ở cùng nhiệt độ $20\ ^\circ\mathrm{C}$ và áp suất $1,913 \cdot 10^5\ \mathrm{Pa}$.
\choiceTF
{\True Ở cùng nhiệt độ và áp suất, các chai có cùng thể tích thì số mol khí chứa trong mỗi chai là bằng nhau.}
{\True Biết khối lượng khí helium trong chai thứ nhất đo được là $0,18\ \mathrm{g}$. Số mol khí helium có trong chai này bằng $0,045\ \mathrm{mol}$.}
{Khối lượng của khí butane trong chai thứ hai bằng $2,02\ \mathrm{g}$.}
{\True Dung tích của mỗi chai thủy tinh xấp xỉ bằng $0,57\ \mathrm{l}$.}
\loigiai{
\begin{itemchoice}
\item (Đúng) Ở cùng nhiệt độ và áp suất, các chai có cùng thể tích thì số mol khí chứa trong mỗi chai là bằng nhau. (Vì): Từ phương trình $pV = nRT$, khi các thông số áp suất, thể tích và nhiệt độ bằng nhau thì số mol $n$ bắt buộc phải bằng nhau. Mệnh đề đúng
\item (Sai) Biết khối lượng khí helium trong chai thứ nhất đo được là $0,18\ \mathrm{g}$. Số mol khí helium có trong chai này bằng $0,045\ \mathrm{mol}$. (Vì): ố mol khí helium có trong chai thứ nhất là:
\item (Sai) Khối lượng của khí butane trong chai thứ hai bằng $2,02\ \mathrm{g}$. (Vì): Vì số mol các khí trong ba chai bằng nhau nên khối lượng khí butane trong chai hai phải là $0,045 \cdot 58 = 2,61\ \mathrm{g}$, còn $2,02\ \mathrm{g}$ là khối lượng khí $\mathrm{CO_2}$. Mệnh đề sai
\item (Đúng) Dung tích của mỗi chai thủy tinh xấp xỉ bằng $0,57\ \mathrm{l}$. (Vì): Áp dụng phương trình Clapeyron để tính dung tích bình với $T = 293\ \mathrm{K}$:
\end{itemchoice}
}
\end{ex}

% ===== Câu 86 =====
% ID: L12C2B11-22-TL
% Mức: VD
\dangbt{Xác định số mol, khối lượng và khối lượng riêng của khí}
\begin{bt}
% ID: L12C2B11-22-TL
% Mức: VD
Trong dạng “Xác định số mol, khối lượng và khối lượng riêng của khí”, 2
\loigiai{
$n=m/M=56/28=2$ mol.
}
\end{bt}

% ===== Câu 87 =====
% ID: L12C2B11-22-TLN
% Mức: VD
\begin{ex}
% ID: L12C2B11-22-TLN
% Mức: VD
Trong dạng “Xác định số mol, khối lượng và khối lượng riêng của khí”, Một khí có khối lượng $56\,\text{g}$, khối lượng mol $28\,\text{g}$ /mol. Tính số mol.
\shortans{2}
\loigiai{
$n=m/M=56/28=2$ mol.
}
\end{ex}

% ===== Câu 88 =====
% ID: L12C2B11-22-TN
% Mức: NB
\dangbt{Tính các thông số trạng thái p, V, T}
\begin{ex}
% ID: L12C2B11-22-TN
% Mức: NB
Trong dạng “Tính các thông số trạng thái p, $V$, T”, Phương trình Clapeyron của khí lí tưởng là
\choice
{\True $pV=nRT$}
{$pV=RT/n$}
{$p/T=nV$}
{$p=nR/(VT)$}
\loigiai{
Đáp án A. Một lượng khí lí tưởng thỏa $pV=nRT$.
}
\end{ex}

% ===== Câu 89 =====
% ID: L12C2B11-23-DS
% Mức: VDC
\dangbt{Xác định sự biến thiên khối lượng riêng của chất khí}
\begin{ex}
% ID: L12C2B11-23-DS
% Mức: VDC
Một bình chứa khí nitrogen ($\mathrm{N_2}$) có thể tích $10\ \mathrm{l}$ ở áp suất $10^5\ \mathrm{Pa}$ và nhiệt độ $27\ ^\circ\mathrm{C}$. Khối lượng mol của nitrogen là $28\ \mathrm{g/mol}$.
\choiceTF
{\True Ở nhiệt độ phòng và áp suất khí quyển, các phân tử khí nitrogen chuyển động hỗn loạn không ngừng va chạm thành bình gây áp suất.}
{\True Số mol của lượng khí nitrogen có trong bình bằng $0,40\ \mathrm{mol}$.}
{Khối lượng của lượng khí nitrogen chứa trong bình xấp xỉ bằng $15,4\ \mathrm{g}$.}
{\True Nếu nhiệt độ tăng lên đến $127\ ^\circ\mathrm{C}$ mà thể tích bình không đổi, để áp suất khí trong bình vẫn giữ ở $10^5\ \mathrm{Pa}$ thì khối lượng khí nitrogen cần xả ra ngoài là $2,8\ \mathrm{g}$.}
\loigiai{
\begin{itemchoice}
\item (Đúng) Ở nhiệt độ phòng và áp suất khí quyển, các phân tử khí nitrogen chuyển động hỗn loạn không ngừng va chạm thành bình gây áp suất. (Vì): Thuyết động học phân tử giải thích áp suất chất khí sinh ra do va chạm hỗn loạn của vô số phân tử lên bề mặt diện tích thành chứa. Mệnh đề đúng
\item (Đúng) Số mol của lượng khí nitrogen có trong bình bằng $0,40\ \mathrm{mol}$. (Vì): Áp dụng phương trình Clapeyron để tính số mol với $V = 10 \cdot 10^{-3}\ \mathrm{m^3}$ và $T = 300\ \mathrm{K}$:
\item (Sai) Khối lượng của lượng khí nitrogen chứa trong bình xấp xỉ bằng $15,4\ \mathrm{g}$. (Vì): Từ kết quả số mol trên, khối lượng nitrogen trong bình ban đầu là:
\item (Đúng) Nếu nhiệt độ tăng lên đến $127\ ^\circ\mathrm{C}$ mà thể tích bình không đổi, để áp suất khí trong bình vẫn giữ ở $10^5\ \mathrm{Pa}$ thì khối lượng khí nitrogen cần xả ra ngoài là $2,8\ \mathrm{g}$. (Vì): Tại nhiệt độ $T' = 400\ \mathrm{K}$ với áp suất giữ nguyên, số mol khí còn lại là $n' = \dfrac{10^5 \cdot 10 \cdot 10^{-3}}{8,31 \cdot 400} \approx 0,301\ \mathrm{mol}$. Khối lượng khí còn lại là $0,301 \cdot 28 \approx 8,43\ \mathrm{g}$. Lượng xả ra là $11,23 - 8,43 = 2,8\ \mathrm{g}$. Mệnh đề đúng
\end{itemchoice}
}
\end{ex}

% ===== Câu 90 =====
% ID: L12C2B11-23-TL
% Mức: VD
\dangbt{Xác định số mol, khối lượng và khối lượng riêng của khí}
\begin{ex}
% ID: L12C2B11-23-TL
% Mức: VD
Trong dạng “Xác định số mol, khối lượng và khối lượng riêng của khí”, Từ $pV=nRT$, khi $T$, $V$ cố định thì $p\sim n$.
\choiceTF

\loigiai{
Từ $pV=nRT$, khi $T$, $V$ cố định thì $p\sim n$.
}
\end{ex}

% ===== Câu 91 =====
% ID: L12C2B11-23-TLN
% Mức: VD
\begin{ex}
% ID: L12C2B11-23-TLN
% Mức: VD
Trong dạng “Xác định số mol, khối lượng và khối lượng riêng của khí”, Ở cùng $T$ và $V$, áp suất tăng 4 lần. Số mol tăng bao nhiêu lần?
\shortans{4}
\loigiai{
Từ $pV=nRT$, khi $T$, $V$ cố định thì $p\sim n$.
}
\end{ex}

% ===== Câu 92 =====
% ID: L12C2B11-23-TN
% Mức: NB
\dangbt{Tính các thông số trạng thái p, V, T}
\begin{ex}
% ID: L12C2B11-23-TN
% Mức: NB
Trong dạng “Tính các thông số trạng thái p, $V$, T”, Với lượng khí xác định, phương trình trạng thái giữa hai trạng thái là
\choice
{\True $p_1V_1/T_1=p_2V_2/T_2$}
{$p_1V_1T_1=p_2V_2T_2$}
{$p_1/T_1=p_2/T_2$ trong mọi quá trình}
{$V_1/T_1=V_2/T_2$ trong mọi quá trình}
\loigiai{
Đáp án A. Do $pV/T=nR$ không đổi.
}
\end{ex}

% ===== Câu 93 =====
% ID: L12C2B11-24-DS
% Mức: VDC
\dangbt{Xác định sự biến thiên khối lượng riêng của chất khí}
\begin{ex}
% ID: L12C2B11-24-DS
% Mức: VDC
Túi khí bảo vệ trong vô lăng xe ô tô hoạt động dựa trên sự phân hủy chất rắn $\mathrm{NaN_3}$ theo phương trình: $2\mathrm{NaN_3(r)} \xrightarrow{t^\circ} 2\mathrm{Na(r)} + 3\mathrm{N_2(k)}$. Ban đầu túi chứa $130\ \mathrm{g}\ \mathrm{NaN_3}$ nguyên chất. Cho biết $M_{\mathrm{NaN_3}} = 65\ \mathrm{g/mol}$.
\choiceTF
{\True Phản ứng phân hủy sodium azide là một phản ứng hóa học giải phóng ra khí nitrogen làm phồng túi an toàn.}
{Lượng khí nitrogen ($\mathrm{N_2}$) giải phóng ra sau khi phân hủy hoàn toàn khối lượng $\mathrm{NaN_3}$ trên là $2,0\ \mathrm{mol}$.}
{\True Ở nhiệt độ $27\ ^\circ\mathrm{C}$, thể tích khí nitrogen thu được dưới áp suất khí quyển $1,0 \cdot 10^5\ \mathrm{Pa}$ xấp xỉ bằng $74,8\ \mathrm{l}$.}
{Biết túi khí phồng lên căng tối đa có thể tích $50\ \mathrm{l}$. Nếu nhiệt độ khí trong túi đạt $47\ ^\circ\mathrm{C}$ thì áp suất khí nitrogen trong túi lúc này bằng $1,85 \cdot 10^5\ \mathrm{Pa}$.}
\loigiai{
\begin{itemchoice}
\item (Đúng) Phản ứng phân hủy sodium azide là một phản ứng hóa học giải phóng ra khí nitrogen làm phồng túi an toàn. (Vì): Phản ứng tạo ra khí nitrogen giãn nở cực nhanh nâng đỡ người lái tránh va đập thứ cấp vào vô lăng xe. Mệnh đề đúng
\item (Sai) Lượng khí nitrogen ($\mathrm{N_2}$) giải phóng ra sau khi phân hủy hoàn toàn khối lượng $\mathrm{NaN_3}$ trên là $2,0\ \mathrm{mol}$. (Vì): ố mol sodium azide ban đầu $n_{\mathrm{NaN_3}} = \dfrac{130}{65} = 2,0\ \mathrm{mol}$. Lượng khí nitrogen thu được phải bằng:
\item (Đúng) Ở nhiệt độ $27\ ^\circ\mathrm{C}$, thể tích khí nitrogen thu được dưới áp suất khí quyển $1,0 \cdot 10^5\ \mathrm{Pa}$ xấp xỉ bằng $74,8\ \mathrm{l}$. (Vì): Áp dụng phương trình Clapeyron để tính thể tích khí ở nhiệt độ $T = 300\ \mathrm{K}$:
\item (Sai) Biết túi khí phồng lên căng tối đa có thể tích $50\ \mathrm{l}$. Nếu nhiệt độ khí trong túi đạt $47\ ^\circ\mathrm{C}$ thì áp suất khí nitrogen trong túi lúc này bằng $1,85 \cdot 10^5\ \mathrm{Pa}$. (Vì): Từ lượng khí trên, áp suất trong túi khí căng ở nhiệt độ $T' = 320\ \mathrm{K}$ đạt giá trị:
\end{itemchoice}
}
\end{ex}

% ===== Câu 94 =====
% ID: L12C2B11-24-TL
% Mức: VDC
\dangbt{Xác định số mol, khối lượng và khối lượng riêng của khí}
\begin{bt}
% ID: L12C2B11-24-TL
% Mức: VDC
Trong dạng “Xác định số mol, khối lượng và khối lượng riêng của khí”, undefined
\loigiai{
$V_2/V_1=(T_2/T_1)/(p_2/p_1)=1$.
}
\end{bt}

% ===== Câu 95 =====
% ID: L12C2B11-24-TLN
% Mức: VDC
\begin{ex}
% ID: L12C2B11-24-TLN
% Mức: VDC
Trong dạng “Xác định số mol, khối lượng và khối lượng riêng của khí”, Khí chuyển trạng thái sao cho p tăng 2 lần và $T$ tăng 2 lần. Thể tích thay đổi bao nhiêu lần?
\shortans{1}
\loigiai{
$V_2/V_1=(T_2/T_1)/(p_2/p_1)=1$.
}
\end{ex}

% ===== Câu 96 =====
% ID: L12C2B11-24-TN
% Mức: TH
\dangbt{Tính các thông số trạng thái p, V, T}
\begin{ex}
% ID: L12C2B11-24-TN
% Mức: TH
Trong dạng “Tính các thông số trạng thái p, $V$, T”, Có 1 mol khí ở 250 $K$ chiếm 6 lít. Lấy R= $8,31\,\text{J}$ /(mol.K). Áp suất gần bằng
\choice
{\True $346,25\,\text{kPa}$}
{$34,63\,\text{kPa}$}
{$3462,5\,\text{kPa}$}
{$0,03\,\text{kPa}$}
\loigiai{
Đáp án A. $p=nRT/V=346,25$ kPa.
}
\end{ex}

% ===== Câu 97 =====
% ID: L12C2B11-25-DS
% Mức: VDC
\dangbt{Xác định sự biến thiên khối lượng riêng của chất khí}
\begin{ex}
% ID: L12C2B11-25-DS
% Mức: VDC
Một túi khí ô tô hoạt động nhờ phản ứng phân hủy hoàn toàn $65\ \mathrm{g}$ chất rắn sodium azide ($\mathrm{NaN_3}$) giải phóng khí nitrogen ($\mathrm{N_2}$) làm phồng túi khí đạt dung tích tối đa là $24\ \mathrm{l}$ ở nhiệt độ $27\ ^\circ\mathrm{C}$. Cho biết khối lượng mol của $\mathrm{NaN_3}$ là $65\ \mathrm{g/mol}$, hằng số khí lí tưởng $R \approx 8,31\ \mathrm{J/(mol\cdot K)}$. Các nhận định dưới đây về hoạt động của túi khí là đúng hay sai?
\choiceTF
{\True Phản ứng phân hủy hoàn toàn $65\ \mathrm{g}$ chất rắn $\mathrm{NaN_3}$ giải phóng ra đúng $1,5\ \mathrm{mol}$ khí nitrogen.}
{\True Thể tích của túi khí khi đã phồng lên hoàn toàn tương đương với $0,024\ \mathrm{m^3}$.}
{\True Áp suất của khí nitrogen trong túi khí khi đã phồng lên hoàn toàn ở nhiệt độ $27\ ^\circ\mathrm{C}$ xấp xỉ bằng $1,56 \cdot 10^5\ \mathrm{Pa}$.}
{Nếu nhiệt độ trong túi khí tăng lên đến $127\ ^\circ\mathrm{C}$ trong khi thể tích túi khí vẫn giữ không đổi ở $24\ \mathrm{l}$ thì áp suất khí lúc này đạt giá trị nhỏ hơn $1,5 \cdot 10^5\ \mathrm{Pa}$.}
\loigiai{
\begin{itemchoice}
\item (Đúng) Phản ứng phân hủy hoàn toàn $65\ \mathrm{g}$ chất rắn $\mathrm{NaN_3}$ giải phóng ra đúng $1,5\ \mathrm{mol}$ khí nitrogen. (Vì): Phương trình hóa học của phản ứng phân hủy: $2\mathrm{NaN_3} \xrightarrow{t^\circ} 2\mathrm{Na} + 3\mathrm{N_2}$. Số mol $\mathrm{NaN_3}$ phản ứng là $n = \dfrac{65}{65} = 1,0\ \mathrm{mol}$. Số mol khí $\mathrm{N_2}$ thu được là:
\item (Đúng) Thể tích của túi khí khi đã phồng lên hoàn toàn tương đương với $0,024\ \mathrm{m^3}$. (Vì): Thể tích túi khí đổi sang đơn vị mét khối:
\item (Đúng) Áp suất của khí nitrogen trong túi khí khi đã phồng lên hoàn toàn ở nhiệt độ $27\ ^\circ\mathrm{C}$ xấp xỉ bằng $1,56 \cdot 10^5\ \mathrm{Pa}$. (Vì): Áp suất khí nitrogen trong túi khí được xác định bằng công thức:
\item (Sai) Nếu nhiệt độ trong túi khí tăng lên đến $127\ ^\circ\mathrm{C}$ trong khi thể tích túi khí vẫn giữ không đổi ở $24\ \mathrm{l}$ thì áp suất khí lúc này đạt giá trị nhỏ hơn $1,5 \cdot 10^5\ \mathrm{Pa}$. (Vì): Khi thể tích túi khí không đổi, áp suất tỉ lệ thuận với nhiệt độ tuyệt đối $T$. Khi nhiệt độ tăng đến $127\ ^\circ\mathrm{C}$ ($T' = 400\ \mathrm{K}$):
\end{itemchoice}
}
\end{ex}

% ===== Câu 98 =====
% ID: L12C2B11-25-TL
% Mức: TH
\dangbt{Áp dụng phương trình Clapeyron pV = nRT}
\begin{ex}
% ID: L12C2B11-25-TL
% Mức: TH
Trong dạng “Áp dụng phương trình Clapeyron pV = nRT”, Trình bày cơ sở lí thuyết của dạng “Áp dụng phương trình Clapeyron pV = nRT” và nêu điều kiện áp dụng.
\choiceTF

\loigiai{
Nêu đúng đại lượng không đổi, hệ thức đặc trưng và yêu cầu dùng nhiệt độ kelvin, áp suất tuyệt đối khi có liên quan.
}
\end{ex}

% ===== Câu 99 =====
% ID: L12C2B11-25-TLN
% Mức: VD
\dangbt{Xác định sự biến thiên khối lượng riêng của chất khí}
\begin{ex}
% ID: L12C2B11-25-TLN
% Mức: VD
\nguon{Bài 11.7 Sách bài tập Vật lí 12 Kết nối tri thức với cuộc sống}
	Ở điều kiện chuẩn ở sát mặt đất, không khí có áp suất $760\ \mathrm{mmHg}$, nhiệt độ $0\ ^\circ\mathrm{C}$ và khối lượng riêng là $1,29\ \mathrm{kg/m^3}$. Trên đỉnh một ngọn núi có áp suất khí quyển đo được là $456\ \mathrm{mmHg}$ và nhiệt độ là $2\ ^\circ\mathrm{C}$. Tính khối lượng riêng của không khí tại đỉnh núi này theo đơn vị $\mathrm{kg/m^3}$ (lấy kết quả làm tròn đến hai chữ số thập phân).
\shortans{$0{,}77$}
\loigiai{
Áp dụng công thức
  \begin{aligned}
 $D&= D_0\cdot\dfrac{p}{p_0}\cdot\dfrac{T_0}{T}&= 1,29\cdot\dfrac{456}{760}\cdot\dfrac{273}{2 + 273}$ &\approx 0,77 $\\mathrm{kg/m^3}$.
 \end{aligned}
}
\end{ex}

% ===== Câu 100 =====
% ID: L12C2B11-25-TN
% Mức: TH
\dangbt{Tính các thông số trạng thái p, V, T}
\begin{ex}
% ID: L12C2B11-25-TN
% Mức: TH
Trong dạng “Tính các thông số trạng thái p, $V$, T”, Khí chuyển từ $p_1=130$ kPa, $V_1=8 lít,$ T_1=260 $K đến$ V_2=24 $lít,$ T_2=340 $K.$ p_2 $gần bằng$
\choice
{\True $56,67\,\text{kPa}$}
{$510\,\text{kPa}$}
{$99,41\,\text{kPa}$}
{$130\,\text{kPa}$}
\loigiai{
Đáp án A. $p_2=p_1V_1T_2/(V_2T_1)=56,67$ kPa.
}
\end{ex}

% ===== Câu 101 =====
% ID: L12C2B11-26-DS
% Mức: VDC
\dangbt{Xác định sự biến thiên khối lượng riêng của chất khí}
\begin{ex}
% ID: L12C2B11-26-DS
% Mức: VDC
Một khối lượng khí hydrogen $12\ \mathrm{g}$ chiếm thể tích ban đầu $4,0\ \mathrm{l}$ ở nhiệt độ $7\ ^\circ\mathrm{C}$ dưới áp suất $p_1$. Người ta tiến hành đun nóng đẳng áp khối khí này cho đến khi khối lượng riêng của nó giảm còn $1,2\ \mathrm{g/l}$.
\choiceTF
{\True Trong quá trình đẳng áp của một khối lượng khí xác định, khối lượng riêng của khí thay đổi tỉ lệ nghịch với thể tích của nó.}
{\True Khối lượng riêng ban đầu của khối khí hydrogen là $3,0\ \mathrm{g/l}$.}
{\True Nhiệt độ tuyệt đối $T_2$ của khối khí ở trạng thái sau bằng $700\ \mathrm{K}$.}
{Thể tích của khối khí hydrogen ở trạng thái sau bằng $15\ \mathrm{l}$.}
\loigiai{
\begin{itemchoice}
\item (Đúng) Trong quá trình đẳng áp của một khối lượng khí xác định, khối lượng riêng của khí thay đổi tỉ lệ nghịch với thể tích của nó. (Vì): o khối lượng không đổi, khi thể tích dãn ra trong quá trình đẳng áp, mật độ khối lượng khí trên đơn vị thể tích giảm tỉ lệ nghịch. Mệnh đề đúng
\item (Đúng) Khối lượng riêng ban đầu của khối khí hydrogen là $3,0\ \mathrm{g/l}$. (Vì): Khối lượng riêng ban đầu của khối khí hydrogen là:
\item (Đúng) Nhiệt độ tuyệt đối $T_2$ của khối khí ở trạng thái sau bằng $700\ \mathrm{K}$. (Vì): Quá trình biến đổi đẳng áp có khối lượng riêng tỉ lệ nghịch với nhiệt độ tuyệt đối ($D_1 \cdot T_1 = D_2 \cdot T_2$):
\item (Sai) Thể tích của khối khí hydrogen ở trạng thái sau bằng $15\ \mathrm{l}$. (Vì): Từ khối lượng và khối lượng riêng trạng thái sau, thể tích khối khí bằng:
\end{itemchoice}
}
\end{ex}

% ===== Câu 102 =====
% ID: L12C2B11-26-TL
% Mức: TH
\dangbt{Áp dụng phương trình Clapeyron pV = nRT}
\begin{ex}
% ID: L12C2B11-26-TL
% Mức: TH
Trong dạng “Áp dụng phương trình Clapeyron pV = nRT”, Mô tả dạng đồ thị đặc trưng và cách nhận biết quá trình trong dạng “Áp dụng phương trình Clapeyron pV = nRT”.
\choiceTF

\loigiai{
Xác định các trục, đại lượng không đổi và suy ra hình dạng đồ thị từ hệ thức vật lí tương ứng.
}
\end{ex}

% ===== Câu 103 =====
% ID: L12C2B11-26-TLN
% Mức: VD
\dangbt{Xác định sự biến thiên khối lượng riêng của chất khí}
\begin{ex}
% ID: L12C2B11-26-TLN
% Mức: VD
Một khối khí lí tưởng xác định đựng trong một xi lanh có pit-tông di chuyển được. Ban đầu khí có khối lượng riêng $1,6\ \mathrm{kg/m^3}$ ở áp suất $2,0 \cdot 10^5\ \mathrm{Pa}$ và nhiệt độ $27\ ^\circ\mathrm{C}$. Người ta tăng áp suất của khí lên đến $3,0 \cdot 10^5\ \mathrm{Pa}$ đồng thời tăng nhiệt độ tuyệt đối thêm $33,33\%$. Tính khối lượng riêng lúc sau của khối khí theo đơn vị $\mathrm{kg/m^3}$ (lấy kết quả làm tròn đến một chữ số thập phân).
\shortans{$1{,}8$}
\loigiai{
Áp dụng công thức
  \begin{aligned}
 D_2 &= D_1 $\cdot\dfrac{p_2}{p_1}\cdot\dfrac{T_1}{T_2}&= 1,6\cdot\dfrac{3,0 \cdot 10^5}{2,0 \cdot 10^5}\cdot\dfrac{1}{1,3333}&= 1,8\\mathrm{kg/m^3}$.
 \end{aligned}
}
\end{ex}

% ===== Câu 104 =====
% ID: L12C2B11-26-TN
% Mức: TH
\dangbt{Tính các thông số trạng thái p, V, T}
\begin{ex}
% ID: L12C2B11-26-TN
% Mức: TH
Trong dạng “Tính các thông số trạng thái p, $V$, T”, Trong phương trình khí lí tưởng, nhiệt độ phải dùng
\choice
{\True kelvin}
{độ $C$}
{độ $F$}
{thang tùy ý}
\loigiai{
Đáp án A. $T$ phải là nhiệt độ tuyệt đối.
}
\end{ex}

% ===== Câu 105 =====
% ID: L12C2B11-27-DS
% Mức: VDC
\dangbt{Xác định sự biến thiên khối lượng riêng của chất khí}
\begin{ex}
% ID: L12C2B11-27-DS
% Mức: VDC
Bóng thám không được bơm khí helium ở áp suất $1,02 \cdot 10^5\ \mathrm{Pa}$ và nhiệt độ $300\ \mathrm{K}$ tại mặt đất. Bóng bay lên cao đến độ cao có áp suất giảm còn $0,30 \cdot 10^5\ \mathrm{Pa}$ và nhiệt độ giảm còn $200\ \mathrm{K}$.
\choiceTF
{\True Bóng thám không bay lên được là do khối lượng riêng của bóng nhỏ hơn khối lượng riêng của không khí bên ngoài.}
{Càng lên cao áp suất khí quyển giảm nên thể tích của bóng giảm theo.}
{\True Khi đạt đến độ cao trên, thể tích của bóng thám không dãn to ra gấp $2,27\ \text{lần}$ thể tích ban đầu.}
{Nếu vỏ bóng hình cầu có bán kính tại mặt đất là $3,0\ \mathrm{m}$, thì khi đạt đến độ cao trên, bán kính của quả bóng thám không xấp xỉ bằng $5,1\ \mathrm{m}$.}
\loigiai{
\begin{itemchoice}
\item (Đúng) Bóng thám không bay lên được là do khối lượng riêng của bóng nhỏ hơn khối lượng riêng của không khí bên ngoài. (Vì): Nhờ khối lượng riêng trung bình của bóng thám không nhỏ hơn không khí bên ngoài nên lực nâng Archimede thắng được trọng lực và nâng bóng lên cao. Mệnh đề đúng
\item (Sai) Càng lên cao áp suất khí quyển giảm nên thể tích của bóng giảm theo. (Vì): Càng lên cao áp suất khí quyển bao quanh giảm mạnh làm cho chất khí bên trong bóng dãn nở đẩy căng vỏ, khiến thể tích của bóng tăng lên rất lớn. Mệnh đề sai
\item (Đúng) Khi đạt đến độ cao trên, thể tích của bóng thám không dãn to ra gấp $2,27\ \text{lần}$ thể tích ban đầu. (Vì): Áp dụng phương trình trạng thái của khí lí tưởng cho lượng helium trong bóng:
\item (Sai) Nếu vỏ bóng hình cầu có bán kính tại mặt đất là $3,0\ \mathrm{m}$, thì khi đạt đến độ cao trên, bán kính của quả bóng thám không xấp xỉ bằng $5,1\ \mathrm{m}$. (Vì): Thể tích hình cầu tỉ lệ thuận với lập phương bán kính $V \sim R^3$, tỉ lệ bán kính đạt được là:
\end{itemchoice}
}
\end{ex}

% ===== Câu 106 =====
% ID: L12C2B11-27-TL
% Mức: VD
\dangbt{Áp dụng phương trình Clapeyron pV = nRT}
\begin{bt}
% ID: L12C2B11-27-TL
% Mức: VD
Trong dạng “Áp dụng phương trình Clapeyron pV = nRT”, Khí có $150\,\text{kPa}$, 4 lít ở 300 $K$; chuyển đến $600\,\text{kPa}$ và 360 K. Tính thể tích sau theo lít.
\loigiai{
$V_2=p_1V_1T_2/(p_2T_1)=1,2$ lít.
}
\end{bt}

% ===== Câu 107 =====
% ID: L12C2B11-27-TLN
% Mức: VD
\dangbt{Xác định sự biến thiên khối lượng riêng của chất khí}
\begin{ex}
% ID: L12C2B11-27-TLN
% Mức: VD
Một quả bóng thám không được bơm khí hydrogen ở sát mặt đất có thể tích $10\ \mathrm{m^3}$ ở áp suất $1,0 \cdot 10^5\ \mathrm{Pa}$ và nhiệt độ $27\ ^\circ\mathrm{C}$. Khi quả bóng bay lên cao đến tầng bình lưu có thể tích dãn rộng thành $35\ \mathrm{m^3}$ và áp suất giảm còn $0,24 \cdot 10^5\ \mathrm{Pa}$. Tính nhiệt độ của khí quyển tại độ cao đó theo đơn vị độ Celsius ($^\circ\mathrm{C}$).
\shortans{$-21$}
\loigiai{
Áp dụng công thức
  \begin{aligned}
 T_2 &= T_1 $\cdot\dfrac{p_2 \cdot V_2}{p_1 \cdot V_1}&= (27 + 273)\cdot\dfrac{0,24 \cdot 10^5 \cdot 35}{1,0 \cdot 10^5 \cdot 10}&= 252\\mathrm{K}$.
 \end{aligned} 
 Đổi sang thang Celsius:
  \begin{aligned}
 t_2 &= T_2 - 273 
&= 252 - 273 
&= -21 $\$ ^ $\circ\mathrm{C}$.
 \end{aligned}
}
\end{ex}

% ===== Câu 108 =====
% ID: L12C2B11-27-TN
% Mức: VD
\dangbt{Tính các thông số trạng thái p, V, T}
\begin{ex}
% ID: L12C2B11-27-TN
% Mức: VD
Trong dạng “Tính các thông số trạng thái p, $V$, T”, Ở thể tích không đổi, áp suất của lượng khí xác định tỉ lệ
\choice
{\True thuận với nhiệt độ tuyệt đối}
{nghịch với nhiệt độ}
{thuận với bình phương nhiệt độ}
{không phụ thuộc nhiệt độ}
\loigiai{
Đáp án A. Từ $pV=nRT$, khi $V$ cố định thì $p/T$ không đổi.
}
\end{ex}

% ===== Câu 109 =====
% ID: L12C2B11-28-DS
% Mức: VDC
\dangbt{Xác định sự biến thiên khối lượng riêng của chất khí}
\begin{ex}
% ID: L12C2B11-28-DS
% Mức: VDC
Một lượng khí lí tưởng xác định thực hiện biến đổi trạng thái. Ban đầu (trạng thái 1) khối khí ở nhiệt độ $27\ ^\circ\mathrm{C}$, thể tích $10\ \mathrm{l}$ dưới áp suất $1,0 \cdot 10^5\ \mathrm{Pa}$. Sau đó, khối khí được nén đến thể tích $2,5\ \mathrm{l}$ và nhiệt độ tuyệt đối tăng lên đến $127\ ^\circ\mathrm{C}$ (trạng thái 2). Các phát biểu dưới đây về quá trình biến đổi trạng thái của khối khí là đúng hay sai?
\choiceTF
{Trong quá trình biến đổi trạng thái trên, tích số $p \cdot V$ của khối khí luôn tỉ lệ nghịch với nhiệt độ tuyệt đối $T$.}
{\True Nhiệt độ tuyệt đối của khối khí ở trạng thái 1 và trạng thái 2 lần lượt là $300\ \mathrm{K}$ và $400\ \mathrm{K}$.}
{\True Nếu trong quá trình trên, nhiệt độ tuyệt đối tăng lên gấp đôi và thể tích giảm đi một nửa thì áp suất khối khí lúc sau tăng lên gấp 4 lần so với ban đầu.}
{\True Áp suất của khối khí ở trạng thái 2 xấp xỉ bằng $5,33 \cdot 10^5\ \mathrm{Pa}$.}
\loigiai{
\begin{itemchoice}
\item (Sai) Trong quá trình biến đổi trạng thái trên, tích số $p \cdot V$ của khối khí luôn tỉ lệ nghịch với nhiệt độ tuyệt đối $T$. (Vì): Theo phương trình trạng thái của khí lí tưởng: $\dfrac{p \cdot V}{T} = \text{hằng số} \Rightarrow p \cdot V \sim T$. Do đó tích số $p \cdot V$ tỉ lệ thuận với nhiệt độ tuyệt đối $T$ của khối khí. Mệnh đề sai
\item (Đúng) Nhiệt độ tuyệt đối của khối khí ở trạng thái 1 và trạng thái 2 lần lượt là $300\ \mathrm{K}$ và $400\ \mathrm{K}$. (Vì): ổi nhiệt độ ở hai trạng thái sang thang đo tuyệt đối Kelvin:
\item (Đúng) Nếu trong quá trình trên, nhiệt độ tuyệt đối tăng lên gấp đôi và thể tích giảm đi một nửa thì áp suất khối khí lúc sau tăng lên gấp 4 lần so với ban đầu. (Vì): Khi $T_2 = 2 T_1$ và $V_2 = 0,5 V_1$, ta áp dụng phương trình trạng thái:
\item (Đúng) Áp suất của khối khí ở trạng thái 2 xấp xỉ bằng $5,33 \cdot 10^5\ \mathrm{Pa}$. (Vì): Áp suất của lượng khí ở trạng thái 2 được xác định bởi công thức:
\end{itemchoice}
}
\end{ex}

% ===== Câu 110 =====
% ID: L12C2B11-28-TL
% Mức: VD
\dangbt{Áp dụng phương trình Clapeyron pV = nRT}
\begin{bt}
% ID: L12C2B11-28-TL
% Mức: VD
Trong dạng “Áp dụng phương trình Clapeyron pV = nRT”, 1
\loigiai{
$n=m/M=28/28=1$ mol.
}
\end{bt}

% ===== Câu 111 =====
% ID: L12C2B11-28-TLN
% Mức: VD
\dangbt{Xác định sự biến thiên khối lượng riêng của chất khí}
\begin{ex}
% ID: L12C2B11-28-TLN
% Mức: VD
Một bình dung tích $24,93\ \mathrm{l}$ chứa khí nitrogen ($\mathrm{N_2}$, khối lượng mol $M = 28\ \mathrm{g/mol}$) ở nhiệt độ $27\ ^\circ\mathrm{C}$ dưới áp suất $2,5 \cdot 10^5\ \mathrm{Pa}$. Người ta mở van xả bớt một lượng khí nitrogen ra ngoài làm áp suất khí trong bình giảm còn $1,5 \cdot 10^5\ \mathrm{Pa}$ trong khi nhiệt độ vẫn giữ không đổi. Tính khối lượng khí nitrogen đã thoát ra ngoài theo đơn vị $\mathrm{g}$. Lấy $R = 8,31\ \mathrm{J/(mol\cdot K)}$.
\shortans{$28$}
\loigiai{
Áp dụng công thức
  \begin{aligned}
 $\Deltam &=\dfrac{\Delta p \cdot V \cdot M}{R \cdot T}&=\dfrac{(2,5 \cdot 10^5 - 1,5 \cdot 10^5) \cdot 24,93 \cdot 10^{-3} \cdot 28 \cdot 10^{-3}}{8,31 \cdot (27 + 273)}&= 0,028\\mathrm{kg}= 28\\mathrm{g}$.
 \end{aligned}
}
\end{ex}

% ===== Câu 112 =====
% ID: L12C2B11-28-TN
% Mức: VD
\dangbt{Tính các thông số trạng thái p, V, T}
\begin{ex}
% ID: L12C2B11-28-TN
% Mức: VD
Trong dạng “Tính các thông số trạng thái p, $V$, T”, Số mol khí được tính từ khối lượng m và khối lượng mol $M$ bằng
\choice
{\True $n=m/M$}
{$n=mM$}
{$n=M/m$}
{$n=m+M$}
\loigiai{
Đáp án A. Theo định nghĩa số mol, $n=m/M$.
}
\end{ex}

% ===== Câu 113 =====
% ID: L12C2B11-29-DS
% Mức: VDC
\dangbt{Xác định sự biến thiên khối lượng riêng của chất khí}
\begin{ex}
% ID: L12C2B11-29-DS
% Mức: VDC
Hệ thống túi khí an toàn trên xe hơi chứa $65\ \mathrm{g}$ chất rắn sodium azide ($\mathrm{NaN_3}$). Khi xe va chạm mạnh, cảm biến kích thích phản ứng phân hủy hoàn toàn giải phóng ra lượng khí nitrogen làm phồng túi khí có thể tích $40\ \mathrm{l}$ ở nhiệt độ $27\ ^\circ\mathrm{C}$.
\choiceTF
{\True Na tạo thành là chất rắn bám vào màng túi khí, còn khí $\mathrm{N_2}$ phân bố đều gây ra áp suất lên toàn bộ màng túi khí.}
{\True Số mol khí nitrogen được giải phóng ra từ phản ứng bằng $1,5\ \mathrm{mol}$.}
{Áp suất của khí nitrogen trong túi khí khi đã phồng lên hoàn toàn ở nhiệt độ $27\ ^\circ\mathrm{C}$ đạt giá trị bằng $0,50 \cdot 10^5\ \mathrm{Pa}$.}
{\True Nếu thực tế hiệu suất phản ứng phân hủy chỉ đạt $80\%$, áp suất của khí trong túi ở $27\ ^\circ\mathrm{C}$ chỉ còn xấp xỉ bằng $0,75 \cdot 10^5\ \mathrm{Pa}$.}
\loigiai{
\begin{itemchoice}
\item (Sai) Na tạo thành là chất rắn bám vào màng túi khí, còn khí $\mathrm{N_2}$ phân bố đều gây ra áp suất lên toàn bộ màng túi khí. (Vì): odium nóng chảy bám vào màng lọc, lượng khí nitrogen tự do chiếm giữ toàn bộ không gian và sinh áp lực nén lên màng. Mệnh đề đúng
\item (Sai) Số mol khí nitrogen được giải phóng ra từ phản ứng bằng $1,5\ \mathrm{mol}$. (Vì): ố mol sodium azide phản ứng $n_{\mathrm{NaN_3}} = \dfrac{65}{65} = 1,0\ \mathrm{mol}$. Số mol nitrogen giải phóng theo tỉ lệ là:
\item (Sai) Áp suất của khí nitrogen trong túi khí khi đã phồng lên hoàn toàn ở nhiệt độ $27\ ^\circ\mathrm{C}$ đạt giá trị bằng $0,50 \cdot 10^5\ \mathrm{Pa}$. (Vì): Áp dụng phương trình Clapeyron để tìm áp suất khí lý tưởng bên trong túi:
\item (Đúng) Nếu thực tế hiệu suất phản ứng phân hủy chỉ đạt $80\%$, áp suất của khí trong túi ở $27\ ^\circ\mathrm{C}$ chỉ còn xấp xỉ bằng $0,75 \cdot 10^5\ \mathrm{Pa}$. (Vì): Từ kết quả áp suất lý thuyết trên, khi hiệu suất phản ứng đạt $80\%$, áp suất thực tế là:
\end{itemchoice}
}
\end{ex}

% ===== Câu 114 =====
% ID: L12C2B11-29-TL
% Mức: VD
\dangbt{Áp dụng phương trình Clapeyron pV = nRT}
\begin{ex}
% ID: L12C2B11-29-TL
% Mức: VD
Trong dạng “Áp dụng phương trình Clapeyron pV = nRT”, Từ $pV=nRT$, khi $T$, $V$ cố định thì $p\sim n$.
\choiceTF

\loigiai{
Từ $pV=nRT$, khi $T$, $V$ cố định thì $p\sim n$.
}
\end{ex}

% ===== Câu 115 =====
% ID: L12C2B11-29-TLN
% Mức: VD
\dangbt{Xác định sự biến thiên khối lượng riêng của chất khí}
\begin{ex}
% ID: L12C2B11-29-TLN
% Mức: VD
Khi nén đẳng nhiệt một khối khí xác định để thể tích của nó giảm đi $2\ \mathrm{l}$ thì áp suất tăng thêm $1,5\ \mathrm{atm}$. Biết thể tích ban đầu của khối khí là $6\ \mathrm{l}$. Tính áp suất ban đầu của khối khí theo đơn vị $\mathrm{atm}$.
\shortans{$3$}
\loigiai{
Áp dụng công thức
  \begin{aligned}
 p_1 $\cdotV_1 &= p_2\cdotV_2 
p_1\cdot$V_1 &=(p_1 + \Delta p)\cdot(V_1 - \Delta V) p_1\cdot6 &=$(p_1 + 1,5)\cdot$ (6 - 2) 
\Rightarrow p_1 &= 3 $\\mathrm{atm}$.
 \end{aligned}
}
\end{ex}

% ===== Câu 116 =====
% ID: L12C2B11-29-TN
% Mức: VD
\dangbt{Tính các thông số trạng thái p, V, T}
\begin{ex}
% ID: L12C2B11-29-TN
% Mức: VD
Trong dạng “Tính các thông số trạng thái p, $V$, T”, Khối lượng riêng của khí lí tưởng thỏa
\choice
{\True $\rho=pM/(RT)$}
{$\rho=RT/(pM)$}
{$\rho=pRT/M$}
{$\rho=p/(MRT)$}
\loigiai{
Đáp án A. Từ $pV=(m/M)RT$ và $\rho=m/V$.
}
\end{ex}

% ===== Câu 117 =====
% ID: L12C2B11-30-DS
% Mức: VDC
\dangbt{Xác định sự biến thiên khối lượng riêng của chất khí}
\begin{ex}
% ID: L12C2B11-30-DS
% Mức: VDC
Mối liên hệ giữa khối lượng riêng $D$ của khối khí lí tưởng với áp suất $p$ và nhiệt độ tuyệt đối $T$ được ứng dụng để giải thích nhiều hiện tượng tự nhiên và kĩ thuật.
\choiceTF
{Khối lượng riêng của một khối khí lí tưởng xác định tỉ lệ nghịch với áp suất và tỉ lệ thuận với nhiệt độ tuyệt đối của nó.}
{\True Ở cùng một nhiệt độ và áp suất, khí helium ($\mathrm{He}$ có $M = 4\ \mathrm{g/mol}$) luôn có khối lượng riêng nhỏ hơn khí nitrogen ($\mathrm{N_2}$ có $M = 28\ \mathrm{g/mol}$).}
{\True Một khinh khí cầu có dung tích $1000\ \mathrm{m^3}$ được bơm không khí nóng ở nhiệt độ $330\ \mathrm{K}$ dưới áp suất $1,0 \cdot 10^5\ \mathrm{Pa}$. Biết khối lượng riêng của không khí lạnh bên ngoài ở nhiệt độ $275\ \mathrm{K}$ là $1,29\ \mathrm{kg/m^3}$. Khối lượng riêng của không khí nóng trong khinh khí cầu xấp xỉ bằng $1,08\ \mathrm{kg/m^3}$.}
{Lực nâng Archimede tối đa tác dụng lên khinh khí cầu trên có độ lớn bằng $10800\ \mathrm{N}$ (lấy $g = 10\ \mathrm{m/s^2}$).}
\loigiai{
\begin{itemchoice}
\item (Sai) Khối lượng riêng của một khối khí lí tưởng xác định tỉ lệ nghịch với áp suất và tỉ lệ thuận với nhiệt độ tuyệt đối của nó. (Vì): Theo công thức $D = \dfrac{p \cdot M}{R \cdot T}$, khối lượng riêng tỉ lệ thuận với áp suất $p$ và tỉ lệ nghịch với nhiệt độ tuyệt đối $T$. Mệnh đề sai
\item (Đúng) Ở cùng một nhiệt độ và áp suất, khí helium ($\mathrm{He}$ có $M = 4\ \mathrm{g/mol}$) luôn có khối lượng riêng nhỏ hơn khí nitrogen ($\mathrm{N_2}$ có $M = 28\ \mathrm{g/mol}$). (Vì): o khối lượng riêng tỉ lệ thuận với khối lượng mol $M$ ở cùng điều kiện nên helium nhẹ hơn nhiều so với nitrogen. Mệnh đề đúng
\item (Đúng) Một khinh khí cầu có dung tích $1000\ \mathrm{m^3}$ được bơm không khí nóng ở nhiệt độ $330\ \mathrm{K}$ dưới áp suất $1,0 \cdot 10^5\ \mathrm{Pa}$. Biết khối lượng riêng của không khí lạnh bên ngoài ở nhiệt độ $275\ \mathrm{K}$ là $1,29\ \mathrm{kg/m^3}$. Khối lượng riêng của không khí nóng trong khinh khí cầu xấp xỉ bằng $1,08\ \mathrm{kg/m^3}$. (Vì): Tính khối lượng riêng của khối khí nóng giãn nở đẳng áp bên trong khinh khí cầu:
\item (Sai) Lực nâng Archimede tối đa tác dụng lên khinh khí cầu trên có độ lớn bằng $10800\ \mathrm{N}$ (lấy $g = 10\ \mathrm{m/s^2}$). (Vì): Lực đẩy Archimede được quyết định bởi thể tích khinh khí cầu và không khí lạnh bao quanh bên ngoài:
\end{itemchoice}
}
\end{ex}

% ===== Câu 118 =====
% ID: L12C2B11-30-TL
% Mức: VDC
\dangbt{Áp dụng phương trình Clapeyron pV = nRT}
\begin{bt}
% ID: L12C2B11-30-TL
% Mức: VDC
Trong dạng “Áp dụng phương trình Clapeyron pV = nRT”, undefined
\loigiai{
$V_2/V_1=(T_2/T_1)/(p_2/p_1)=1$.
}
\end{bt}

% ===== Câu 119 =====
% ID: L12C2B11-30-TLN
% Mức: VD
\dangbt{Xác định sự biến thiên khối lượng riêng của chất khí}
\begin{ex}
% ID: L12C2B11-30-TLN
% Mức: VD
\nguon{Dựa trên hoạt động 2 trang 47 Sách giáo khoa Vật lí 12 Kết nối tri thức với cuộc sống}
	Một thiết bị túi khí ô tô chứa $100\ \mathrm{g}$ chất rắn $\mathrm{NaN_3}$ (khối lượng mol $M = 65\ \mathrm{g/mol}$). Khi xảy ra va chạm, phản ứng phân hủy $\mathrm{NaN_3}$ xảy ra hoàn toàn tạo ra khí nitrogen ($\mathrm{N_2}$). Tính thể tích khí $\mathrm{N_2}$ giải phóng ra ở điều kiện thực nghiệm có thể tích mol là $24,0\ \mathrm{l/mol}$ theo đơn vị lít ($\mathrm{l}$) (làm tròn kết quả đến một chữ số thập phân).
\shortans{$55{,}4$}
\loigiai{
Số mol sodium azide tham gia phân hủy là:
  \begin{aligned}
 n_{ $\mathrm{NaN_3}} &=\dfrac{m}{M}&=\dfrac{100}{65}\\mathrm{mol}$.
 \end{aligned} 
 Theo tỉ lệ phản ứng, số mol khí $\mathrm{N_2}$ sinh ra và thể tích của nó là:
  \begin{aligned}
 n_{ $\mathrm{N_2}$ } &= 1,5 \cdot n_{ $\mathrm{NaN_3}$ } 
&= 1,5 $\cdot\dfrac{100}{65}\approx$ 2,3077 $\\mathrm{mol}V$ &= n_{ $\mathrm{N_2}$ } \cdot V_{ $\text{mol}$ } 
&= 2,3077 \cdot 24,0 
&\approx 55,4 $\\mathrm{l}$.
 \end{aligned}
}
\end{ex}

% ===== Câu 120 =====
% ID: L12C2B11-30-TN
% Mức: VDC
\dangbt{Tính các thông số trạng thái p, V, T}
\begin{ex}
% ID: L12C2B11-30-TN
% Mức: VDC
Trong dạng “Tính các thông số trạng thái p, $V$, T”, Hai bình cùng p, $V$, $T$ chứa khí lí tưởng khác loại thì có
\choice
{\True cùng số phân tử}
{cùng khối lượng}
{cùng khối lượng riêng}
{khác số mol}
\loigiai{
Đáp án A. Theo $n=pV/(RT)$, số mol và số phân tử bằng nhau.
}
\end{ex}

% ===== Câu 121 =====
% ID: L12C2B11-31-DS
% Mức: VDC
\dangbt{Xác định sự biến thiên khối lượng riêng của chất khí}
\begin{ex}
% ID: L12C2B11-31-DS
% Mức: VDC
Một bình chứa khí lí tưởng ở nhiệt độ $27\ ^\circ\mathrm{C}$ dưới áp suất $2,0 \cdot 10^5\ \mathrm{Pa}$ có khối lượng riêng ban đầu là $1,60\ \mathrm{kg/m^3}$. Người ta đun nóng khối khí này đến nhiệt độ $127\ ^\circ\mathrm{C}$ và tăng áp suất lên đến $3,0 \cdot 10^5\ \mathrm{Pa}$. Các phát biểu dưới đây về khối lượng riêng của chất khí là đúng hay sai?
\choiceTF
{Khối lượng riêng $D$ của một lượng khí lí tưởng xác định tỉ lệ nghịch với áp suất và tỉ lệ thuận với nhiệt độ tuyệt đối của nó.}
{\True Hệ thức liên hệ giữa khối lượng riêng, áp suất và nhiệt độ của lượng khí lí tưởng ở hai trạng thái là $\dfrac{D_1 \cdot T_1}{p_1} = \dfrac{D_2 \cdot T_2}{p_2}$.}
{\True Nhiệt độ tuyệt đối ở trạng thái ban đầu và trạng thái sau lần lượt bằng $300\ \mathrm{K}$ và $400\ \mathrm{K}$.}
{\True Khối lượng riêng của khối khí ở trạng thái sau bằng $1,80\ \mathrm{kg/m^3}$.}
\loigiai{
\begin{itemchoice}
\item (Sai) Khối lượng riêng $D$ của một lượng khí lí tưởng xác định tỉ lệ nghịch với áp suất và tỉ lệ thuận với nhiệt độ tuyệt đối của nó. (Vì): Theo công thức $D = \dfrac{p \cdot M}{R \cdot T}$, khối lượng riêng $D$ tỉ lệ thuận với áp suất $p$ và tỉ lệ nghịch với nhiệt độ tuyệt đối $T$. Mệnh đề sai
\item (Đúng) Hệ thức liên hệ giữa khối lượng riêng, áp suất và nhiệt độ của lượng khí lí tưởng ở hai trạng thái là $\dfrac{D_1 \cdot T_1}{p_1} = \dfrac{D_2 \cdot T_2}{p_2}$. (Vì): Từ phương trình Clapeyron, ta suy ra biểu thức khối lượng riêng: $D = \dfrac{p \cdot M}{R \cdot T} \Rightarrow \dfrac{D \cdot T}{p} = \dfrac{M}{R} = \text{const}$. Do đó ta có hệ thức $\dfrac{D_1 \cdot T_1}{p_1} = \dfrac{D_2 \cdot T_2}{p_2}$. Mệnh đề đúng
\item (Đúng) Nhiệt độ tuyệt đối ở trạng thái ban đầu và trạng thái sau lần lượt bằng $300\ \mathrm{K}$ và $400\ \mathrm{K}$. (Vì): ổi nhiệt độ ở hai trạng thái sang thang đo tuyệt đối:
\item (Đúng) Khối lượng riêng của khối khí ở trạng thái sau bằng $1,80\ \mathrm{kg/m^3}$. (Vì): Áp dụng hệ thức liên hệ ở câu b, ta tính khối lượng riêng của khối khí lúc sau:
\end{itemchoice}
}
\end{ex}

% ===== Câu 122 =====
% ID: L12C2B11-31-TL
% Mức: TH
\dangbt{Đồ thị và các quá trình biến đổi trạng thái liên tiếp}
\begin{ex}
% ID: L12C2B11-31-TL
% Mức: TH
Trong dạng “Đồ thị và các quá trình biến đổi trạng thái liên tiếp”, Trình bày cơ sở lí thuyết của dạng “Đồ thị và các quá trình biến đổi trạng thái liên tiếp” và nêu điều kiện áp dụng.
\choiceTF

\loigiai{
Nêu đúng đại lượng không đổi, hệ thức đặc trưng và yêu cầu dùng nhiệt độ kelvin, áp suất tuyệt đối khi có liên quan.
}
\end{ex}

% ===== Câu 123 =====
% ID: L12C2B11-31-TLN
% Mức: VD
\dangbt{Xác định sự biến thiên khối lượng riêng của chất khí}
\begin{ex}
% ID: L12C2B11-31-TLN
% Mức: VD
Để một túi khí ô tô có dung tích phồng căng tối đa là $50\ \mathrm{l}$ đạt được áp suất $1,2 \cdot 10^5\ \mathrm{Pa}$ ở nhiệt độ $27\ ^\circ\mathrm{C}$ thì số mol khí nitrogen cần giải phóng ra bằng bao nhiêu mol? Lấy hằng số khí lí tưởng $R = 8,31\ \mathrm{J/(mol\cdot K)}$ (làm tròn kết quả đến hai chữ số thập phân).
\shortans{$2{,}41$}
\loigiai{
Áp dụng công thức
  \begin{aligned}
 n &= $\dfrac{p \cdot V}{R \cdot T}&=\dfrac{1,2 \cdot 10^5 \cdot 50 \cdot 10^{-3}}{8,31 \cdot (27 + 273)}$ &\approx 2,41 $\\mathrm{mol}$.
 \end{aligned}
}
\end{ex}

% ===== Câu 124 =====
% ID: L12C2B11-31-TN
% Mức: VDC
\dangbt{Tính các thông số trạng thái p, V, T}
\begin{ex}
% ID: L12C2B11-31-TN
% Mức: VDC
Trong dạng “Tính các thông số trạng thái p, $V$, T”, Khi giải bài toán khí lí tưởng, áp suất đưa vào phương trình phải là
\choice
{\True áp suất tuyệt đối}
{áp suất dư}
{áp suất do chất lỏng בלבד}
{áp suất tương đối tùy ý}
\loigiai{
Đáp án A. Phương trình trạng thái sử dụng áp suất tuyệt đối.
}
\end{ex}

% ===== Câu 125 =====
% ID: L12C2B11-32-DS
% Mức: VDC
\dangbt{Xác định sự biến thiên khối lượng riêng của chất khí}
\begin{ex}
% ID: L12C2B11-32-DS
% Mức: VDC
Một bình thép dung tích $20,0\ \mathrm{l}$ chứa khí oxygen ($\mathrm{O_2}$) có khối lượng $32,0\ \mathrm{g}$ ở nhiệt độ $27\ ^\circ\mathrm{C}$. Khối lượng mol của oxygen là $32,0\ \mathrm{g/mol}$. Hằng số khí lí tưởng lấy $R \approx 8,31\ \mathrm{J/(mol\cdot K)}$.
\choiceTF
{\True Số mol khí oxygen trong bình bằng $1,0\ \mathrm{mol}$.}
{Hằng số khí lí tưởng $R$ trong hệ SI có giá trị xấp xỉ bằng $0,082\ \mathrm{atm \cdot l / (mol \cdot K)}$.}
{\True Áp suất của khí oxygen trong bình bằng $1,25 \cdot 10^5\ \mathrm{Pa}$.}
{Nếu nhiệt độ tăng lên đến $127\ ^\circ\mathrm{C}$ đồng thời bình bị rò rỉ khí làm áp suất giảm còn $1,0 \cdot 10^5\ \mathrm{Pa}$, thì khối lượng oxygen thoát ra ngoài bằng $16,0\ \mathrm{g}$.}
\loigiai{
\begin{itemchoice}
\item (Sai) Số mol khí oxygen trong bình bằng $1,0\ \mathrm{mol}$. (Vì): ố mol khí oxygen được xác định trực tiếp theo tỉ số khối lượng:
\item (Sai) Hằng số khí lí tưởng $R$ trong hệ SI có giá trị xấp xỉ bằng $0,082\ \mathrm{atm \cdot l / (mol \cdot K)}$. (Vì): Hằng số khí lí tưởng $R \approx 8,31\ \mathrm{J/(mol\cdot K)}$ trong hệ SI, còn giá trị $0,082$ là đơn vị hỗn hợp ngoài hệ SI. Mệnh đề sai
\item (Đúng) Áp suất của khí oxygen trong bình bằng $1,25 \cdot 10^5\ \mathrm{Pa}$. (Vì): Áp dụng phương trình Clapeyron với $V = 20,0 \cdot 10^{-3}\ \mathrm{m^3}$ và $T = 300\ \mathrm{K}$:
\item (Sai) Nếu nhiệt độ tăng lên đến $127\ ^\circ\mathrm{C}$ đồng thời bình bị rò rỉ khí làm áp suất giảm còn $1,0 \cdot 10^5\ \mathrm{Pa}$, thì khối lượng oxygen thoát ra ngoài bằng $16,0\ \mathrm{g}$. (Vì): Tại trạng thái sau với $T' = 400\ \mathrm{K}$ và $p' = 1,0 \cdot 10^5\ \mathrm{Pa}$, số mol khí còn lại trong bình là:
\end{itemchoice}
}
\end{ex}

% ===== Câu 126 =====
% ID: L12C2B11-32-TL
% Mức: TH
\dangbt{Đồ thị và các quá trình biến đổi trạng thái liên tiếp}
\begin{ex}
% ID: L12C2B11-32-TL
% Mức: TH
Trong dạng “Đồ thị và các quá trình biến đổi trạng thái liên tiếp”, Mô tả dạng đồ thị đặc trưng và cách nhận biết quá trình trong dạng “Đồ thị và các quá trình biến đổi trạng thái liên tiếp”.
\choiceTF

\loigiai{
Xác định các trục, đại lượng không đổi và suy ra hình dạng đồ thị từ hệ thức vật lí tương ứng.
}
\end{ex}

% ===== Câu 127 =====
% ID: L12C2B11-32-TLN
% Mức: VD
\dangbt{Xác định sự biến thiên khối lượng riêng của chất khí}
\begin{ex}
% ID: L12C2B11-32-TLN
% Mức: VD
Một quả bóng thám không dâng lên cao từ mặt đất nơi có nhiệt độ $27\ ^\circ\mathrm{C}$ và áp suất $1,0 \cdot 10^5\ \mathrm{Pa}$. Khi lên đến tầng bình lưu, nhiệt độ tuyệt đối của khí trong bóng giảm đi $20\%$ trong khi thể tích của bóng tăng lên gấp $4\ \text{lần}$ so với ban đầu. Tính áp suất của khí quyển tại độ cao đó theo đơn vị $10^4\ \mathrm{Pa}$.
\shortans{$2$}
\loigiai{
Áp dụng công thức
  \begin{aligned}
 p_2 &= p_1 $\cdot\dfrac{V_1}{V_2}\cdot\dfrac{T_2}{T_1}&= 1,0\cdot10^5\cdot\dfrac{1}{4}\cdot0,8 
&= 2\cdot$10^4$\\mathrm{Pa}$.
 \end{aligned}
}
\end{ex}

% ===== Câu 128 =====
% ID: L12C2B11-32-TN
% Mức: NB
\dangbt{Xác định số mol, khối lượng và khối lượng riêng của khí}
\begin{ex}
% ID: L12C2B11-32-TN
% Mức: NB
Trong dạng “Xác định số mol, khối lượng và khối lượng riêng của khí”, Phương trình Clapeyron của khí lí tưởng là
\choice
{\True $pV=nRT$}
{$pV=RT/n$}
{$p/T=nV$}
{$p=nR/(VT)$}
\loigiai{
Đáp án A. Một lượng khí lí tưởng thỏa $pV=nRT$.
}
\end{ex}

% ===== Câu 129 =====
% ID: L12C2B11-33-DS
% Mức: VDC
\dangbt{Xác định sự biến thiên khối lượng riêng của chất khí}
\begin{ex}
% ID: L12C2B11-33-DS
% Mức: VDC
Một bình dung tích $40\ \mathrm{l}$ chứa $3,96\ \mathrm{g}$ khí oxygen ($\mathrm{O_2}$) ở nhiệt độ $27\ ^\circ\mathrm{C}$. Biết khối lượng mol của oxygen là $32\ \mathrm{g/mol}$. Hằng số khí lí tưởng là $R \approx 8,31\ \mathrm{J/(mol\cdot K)}$.
\choiceTF
{\True Oxygen là chất khí có phân tử gồm hai nguyên tử, khối lượng mol của nó bằng $32 \cdot 10^{-3}\ \mathrm{kg/mol}$.}
{\True Số mol khí oxygen có trong bình bằng $0,12375\ \mathrm{mol}$.}
{Áp suất của khí oxygen trong bình bằng $77\ \mathrm{Pa}$.}
{\True Nếu vỏ bình chỉ chịu được áp suất tối đa là $60\ \mathrm{atm}$ thì ở nhiệt độ $27\ ^\circ\mathrm{C}$ bình có thể chứa tối đa $3,13\ \mathrm{kg}$ khí oxygen.}
\loigiai{
\begin{itemchoice}
\item (Đúng) Oxygen là chất khí có phân tử gồm hai nguyên tử, khối lượng mol của nó bằng $32 \cdot 10^{-3}\ \mathrm{kg/mol}$. (Vì): Phân tử oxygen gồm hai nguyên tử liên kết, khối lượng mol tương ứng là $32\ \mathrm{g/mol} = 32 \cdot 10^{-3}\ \mathrm{kg/mol}$. Mệnh đề đúng
\item (Sai) Số mol khí oxygen có trong bình bằng $0,12375\ \mathrm{mol}$. (Vì): ố mol khí oxygen chứa trong bình là:
\item (Sai) Áp suất của khí oxygen trong bình bằng $77\ \mathrm{Pa}$. (Vì): Áp dụng phương trình Clapeyron với $V = 40 \cdot 10^{-3}\ \mathrm{m^3}$ và $T = 300\ \mathrm{K}$:
\item (Đúng) Nếu vỏ bình chỉ chịu được áp suất tối đa là $60\ \mathrm{atm}$ thì ở nhiệt độ $27\ ^\circ\mathrm{C}$ bình có thể chứa tối đa $3,13\ \mathrm{kg}$ khí oxygen. (Vì): Từ giới hạn áp suất $p_{\text{max}} = 60 \cdot 1,013 \cdot 10^5 = 6,078 \cdot 10^6\ \mathrm{Pa}$, khối lượng tối đa là:
\end{itemchoice}
}
\end{ex}

% ===== Câu 130 =====
% ID: L12C2B11-33-TL
% Mức: VD
\dangbt{Đồ thị và các quá trình biến đổi trạng thái liên tiếp}
\begin{bt}
% ID: L12C2B11-33-TL
% Mức: VD
Trong dạng “Đồ thị và các quá trình biến đổi trạng thái liên tiếp”, Khí có $110\,\text{kPa}$, 6 lít ở 320 $K$; chuyển đến $330\,\text{kPa}$ và 380 K. Tính thể tích sau theo lít.
\loigiai{
$V_2=p_1V_1T_2/(p_2T_1)=2,38$ lít.
}
\end{bt}

% ===== Câu 131 =====
% ID: L12C2B11-33-TLN
% Mức: VD
\dangbt{Xác định sự biến thiên khối lượng riêng của chất khí}
\begin{ex}
% ID: L12C2B11-33-TLN
% Mức: VD
Túi khí ô tô an toàn hoạt động dựa trên sự phân hủy chất rắn sodium azide ($\mathrm{NaN_3}$) giải phóng khí nitrogen ($\mathrm{N_2}$) theo phương trình: $2\mathrm{NaN_3(r)} \xrightarrow{t^\circ} 2\mathrm{Na(r)} + 3\mathrm{N_2(k)}$. Để giải phóng ra đúng $3,0\ \mathrm{mol}$ khí nitrogen thì khối lượng $\mathrm{NaN_3}$ nguyên chất cần phân hủy bằng bao nhiêu $\mathrm{g}$? Biết khối lượng mol của $\mathrm{NaN_3}$ là $65\ \mathrm{g/mol}$.
\shortans{$130$}
\loigiai{
Xác định số mol $\mathrm{NaN_3}$ theo tỉ lệ phản ứng:
  \begin{aligned}
 n_{ $\mathrm{NaN_3}$ } &= $\dfrac{2}{3}\cdot$ n_{ $\mathrm{N_2}$ } 
&= $\dfrac{2}{3}\cdot$ 3,0 
&= 2,0 $\\mathrm{mol}$.
 \end{aligned} 
 Khối lượng sodium azide là:
  \begin{aligned}
 m &= n_{ $\mathrm{NaN_3}$ } $\cdotM$ &= 2,0\cdot 65 
&= 130 $\\mathrm{g}$.
 \end{aligned}
}
\end{ex}

% ===== Câu 132 =====
% ID: L12C2B11-33-TN
% Mức: NB
\dangbt{Xác định số mol, khối lượng và khối lượng riêng của khí}
\begin{ex}
% ID: L12C2B11-33-TN
% Mức: NB
Trong dạng “Xác định số mol, khối lượng và khối lượng riêng của khí”, Với lượng khí xác định, phương trình trạng thái giữa hai trạng thái là
\choice
{\True $p_1V_1/T_1=p_2V_2/T_2$}
{$p_1V_1T_1=p_2V_2T_2$}
{$p_1/T_1=p_2/T_2$ trong mọi quá trình}
{$V_1/T_1=V_2/T_2$ trong mọi quá trình}
\loigiai{
Đáp án A. Do $pV/T=nR$ không đổi.
}
\end{ex}

% ===== Câu 133 =====
% ID: L12C2B11-34-DS
% Mức: VDC
\dangbt{Xác định sự biến thiên khối lượng riêng của chất khí}
\begin{ex}
% ID: L12C2B11-34-DS
% Mức: VDC
Khảo sát quá trình bay lên cao của một bóng thám không mang thiết bị cảm biến khí tượng. Bóng chứa một khối lượng khí xác định dãn nở từ thể tích $V_1$ ở mặt đất lên thể tích $V_2$ ở trên cao.
\choiceTF
{Quá trình dãn nở của khí trong bóng khi bay lên cao là một đẳng quá trình vì nhiệt độ khí quyển giảm đều theo độ cao.}
{\True Theo phương trình trạng thái khí lí tưởng, tỉ số áp suất của khí trong bóng lúc sau và lúc đầu là $\dfrac{p_2}{p_1} = \dfrac{V_1}{V_2} \cdot \dfrac{T_2}{T_1}$.}
{\True Nếu thể tích bóng tăng gấp $3\ \text{lần}$ trong khi nhiệt độ tuyệt đối của khối khí giảm đi một nửa thì áp suất khí bên bên trong giảm đi $6\ \text{lần}$ so với mặt đất.}
{Ở độ cao $10\ \mathrm{km}$, áp suất khí quyển giảm đi $4\ \text{lần}$ và nhiệt độ tuyệt đối giảm đi $1,2\ \text{lần}$ so với mặt đất. Khi đó, thể tích của bóng thám không tăng lên gấp $4,8\ \text{lần}$ ban đầu.}
\loigiai{
\begin{itemchoice}
\item (Sai) Quá trình dãn nở của khí trong bóng khi bay lên cao là một đẳng quá trình vì nhiệt độ khí quyển giảm đều theo độ cao. (Vì): Khi bóng dâng lên cao, cả ba thông số trạng thái áp suất, thể tích và nhiệt độ đồng thời biến đổi nên đây không phải là một đẳng quá trình. Mệnh đề sai
\item (Đúng) Theo phương trình trạng thái khí lí tưởng, tỉ số áp suất của khí trong bóng lúc sau và lúc đầu là $\dfrac{p_2}{p_1} = \dfrac{V_1}{V_2} \cdot \dfrac{T_2}{T_1}$. (Vì): Rút ra tỉ số áp suất trực tiếp từ hệ thức biến đổi trạng thái của phương trình trạng thái khí lí tưởng. Mệnh đề đúng
\item (Đúng) Nếu thể tích bóng tăng gấp $3\ \text{lần}$ trong khi nhiệt độ tuyệt đối của khối khí giảm đi một nửa thì áp suất khí bên bên trong giảm đi $6\ \text{lần}$ so với mặt đất. (Vì): Khi $V_2 = 3 V_1$ và $T_2 = 0,5 T_1$, áp dụng công thức tỉ số áp suất:
\item (Sai) Ở độ cao $10\ \mathrm{km}$, áp suất khí quyển giảm đi $4\ \text{lần}$ và nhiệt độ tuyệt đối giảm đi $1,2\ \text{lần}$ so với mặt đất. Khi đó, thể tích của bóng thám không tăng lên gấp $4,8\ \text{lần}$ ban đầu. (Vì): Khi $p_1 = 4 p_2$ và $T_2 = \dfrac{T_1}{1,2}$, thể tích bóng lúc sau tăng lên đạt tỉ số:
\end{itemchoice}
}
\end{ex}

% ===== Câu 134 =====
% ID: L12C2B11-34-TL
% Mức: VD
\dangbt{Đồ thị và các quá trình biến đổi trạng thái liên tiếp}
\begin{bt}
% ID: L12C2B11-34-TL
% Mức: VD
Trong dạng “Đồ thị và các quá trình biến đổi trạng thái liên tiếp”, 3
\loigiai{
$n=m/M=84/28=3$ mol.
}
\end{bt}

% ===== Câu 135 =====
% ID: L12C2B11-34-TLN
% Mức: VD
\dangbt{Xác định sự biến thiên khối lượng riêng của chất khí}
\begin{ex}
% ID: L12C2B11-34-TLN
% Mức: VD
Một quả bóng thám không chứa khí helium có thể tích $18\ \mathrm{m^3}$ khi vừa được bơm xong ở sát mặt đất dưới áp suất $1,0 \cdot 10^5\ \mathrm{Pa}$ và nhiệt độ $300\ \mathrm{K}$. Khi bay lên đến độ cao có áp suất khí quyển là $0,3 \cdot 10^5\ \mathrm{Pa}$ và nhiệt độ tuyệt đối là $200\ \mathrm{K}$ thì thể tích của quả bóng thám không này bằng bao nhiêu $\mathrm{m^3}$?
\shortans{$40$}
\loigiai{
Áp dụng công thức
  \begin{aligned}
 V_2 &= V_1 $\cdot\dfrac{p_1}{p_2}\cdot\dfrac{T_2}{T_1}&= 18\cdot\dfrac{1,0 \cdot 10^5}{0,3 \cdot 10^5}\cdot\dfrac{200}{300}&= 40\\mathrm{m^3}$.
 \end{aligned}
}
\end{ex}

% ===== Câu 136 =====
% ID: L12C2B11-34-TN
% Mức: TH
\dangbt{Xác định số mol, khối lượng và khối lượng riêng của khí}
\begin{ex}
% ID: L12C2B11-34-TN
% Mức: TH
Trong dạng “Xác định số mol, khối lượng và khối lượng riêng của khí”, Có 3 mol khí ở 270 $K$ chiếm 8 lít. Lấy R= $8,31\,\text{J}$ /(mol.K). Áp suất gần bằng
\choice
{\True $841,39\,\text{kPa}$}
{$84,14\,\text{kPa}$}
{$8413,88\,\text{kPa}$}
{$0,09\,\text{kPa}$}
\loigiai{
Đáp án A. $p=nRT/V=841,39$ kPa.
}
\end{ex}

% ===== Câu 137 =====
% ID: L12C2B11-35-DS
% Mức: VDC
\dangbt{Xác định sự biến thiên khối lượng riêng của chất khí}
\begin{ex}
% ID: L12C2B11-35-DS
% Mức: VDC
Một khối khí lí tưởng xác định thực hiện quá trình biến đổi trạng thái từ trạng thái $1$ ($p_1 = 1,0 \cdot 10^5\ \mathrm{Pa}$, $V_1 = 6,0\ \mathrm{l}$, $t_1 = 27\ ^\circ\mathrm{C}$) sang trạng thái $2$ ($p_2 = 2,0 \cdot 10^5\ \mathrm{Pa}$, $V_2 = 4,0\ \mathrm{l}$, $T_2$).
\choiceTF
{\True Nhiệt độ tuyệt đối của khối khí ở trạng thái $1$ là $300\ \mathrm{K}$.}
{Thương số $\dfrac{p \cdot V}{T}$ của khối khí ở trạng thái $1$ là $2,0 \cdot 10^3\ \mathrm{J/K}$.}
{\True Ở trạng thái $2$, nhiệt độ tuyệt đối $T_2$ của khối khí đạt giá trị bằng $400\ \mathrm{K}$.}
{\True Nếu tiếp tục biến đổi đẳng áp từ trạng thái $2$ đến trạng thái $3$ có thể tích $V_3 = 2,0\ \mathrm{l}$ thì nhiệt độ Celsius $t_3$ của khối khí bằng $-73\ ^\circ\mathrm{C}$.}
\loigiai{
\begin{itemchoice}
\item (Đúng) Nhiệt độ tuyệt đối của khối khí ở trạng thái $1$ là $300\ \mathrm{K}$. (Vì): ổi nhiệt độ của khối khí ở trạng thái ban đầu sang thang đo tuyệt đối Kelvin:
\item (Đúng) Thương số $\dfrac{p \cdot V}{T}$ của khối khí ở trạng thái $1$ là $2,0 \cdot 10^3\ \mathrm{J/K}$. (Vì): ổi thể tích $V_1 = 6,0\ \mathrm{l} = 6,0 \cdot 10^{-3}\ \mathrm{m^3}$. Thương số $\dfrac{p \cdot V}{T}$ là hằng số:
\item (Đúng) Ở trạng thái $2$, nhiệt độ tuyệt đối $T_2$ của khối khí đạt giá trị bằng $400\ \mathrm{K}$. (Vì): Áp dụng phương trình trạng thái của khí lí tưởng để tìm nhiệt độ $T_2$:
\item (Đúng) Nếu tiếp tục biến đổi đẳng áp từ trạng thái $2$ đến trạng thái $3$ có thể tích $V_3 = 2,0\ \mathrm{l}$ thì nhiệt độ Celsius $t_3$ của khối khí bằng $-73\ ^\circ\mathrm{C}$. (Vì): Quá trình biến đổi đẳng áp $2 \to 3$ tuân theo định luật Charles:
\end{itemchoice}
}
\end{ex}

% ===== Câu 138 =====
% ID: L12C2B11-35-TL
% Mức: VD
\dangbt{Đồ thị và các quá trình biến đổi trạng thái liên tiếp}
\begin{ex}
% ID: L12C2B11-35-TL
% Mức: VD
Trong dạng “Đồ thị và các quá trình biến đổi trạng thái liên tiếp”, Từ $pV=nRT$, khi $T$, $V$ cố định thì $p\sim n$.
\choiceTF

\loigiai{
Từ $pV=nRT$, khi $T$, $V$ cố định thì $p\sim n$.
}
\end{ex}

% ===== Câu 139 =====
% ID: L12C2B11-35-TLN
% Mức: VDC
\dangbt{Xác định sự biến thiên khối lượng riêng của chất khí}
\begin{ex}
% ID: L12C2B11-35-TLN
% Mức: VDC
\nguon{Dựa trên hoạt động 1 trang 47 Sách giáo khoa Vật lí 12 Kết nối tri thức với cuộc sống}
	Một quả bóng thám không hình cầu khi vừa bơm xong ở mặt đất có bán kính $8\ \mathrm{m}$ ở áp suất $1,02 \cdot 10^5\ \mathrm{Pa}$ và nhiệt độ $300\ \mathrm{K}$. Khi bay lên tầng khí quyển cao có áp suất $0,34 \cdot 10^5\ \mathrm{Pa}$ và nhiệt độ tuyệt đối $200\ \mathrm{K}$. Tính bán kính của quả bóng lúc này theo đơn vị $\mathrm{m}$ (lấy kết quả làm tròn đến hàng đơn vị).
\shortans{$10$}
\loigiai{
Áp dụng công thức
  \begin{aligned}
 $\dfrac{p_1 \cdot R_1^3}{T_1}&=\dfrac{p_2 \cdot R_2^3}{T_2}\RightarrowR_2 &= R_1\cdot\sqrt{\dfrac{p_1 \cdot T_2}{p_2 \cdot T_1}}&= 8\cdot\sqrt{\dfrac{1,02 \cdot 10^5 \cdot 200}{0,34 \cdot 10^5 \cdot 300}}$ &\approx 10 $\\mathrm{m}$.
 \end{aligned}
}
\end{ex}

% ===== Câu 140 =====
% ID: L12C2B11-35-TN
% Mức: TH
\dangbt{Xác định số mol, khối lượng và khối lượng riêng của khí}
\begin{ex}
% ID: L12C2B11-35-TN
% Mức: TH
Trong dạng “Xác định số mol, khối lượng và khối lượng riêng của khí”, Khí chuyển từ $p_1=95$ kPa, $V_1=3 lít,$ T_1=280 $K đến$ V_2=15 $lít,$ T_2=360 $K.$ p_2 $gần bằng$
\choice
{\True $24,43\,\text{kPa}$}
{$610,71\,\text{kPa}$}
{$73,89\,\text{kPa}$}
{$95\,\text{kPa}$}
\loigiai{
Đáp án A. $p_2=p_1V_1T_2/(V_2T_1)=24,43$ kPa.
}
\end{ex}

% ===== Câu 141 =====
% ID: L12C2B11-36-DS
% Mức: VDC
\dangbt{Xác định sự biến thiên khối lượng riêng của chất khí}
\begin{ex}
% ID: L12C2B11-36-DS
% Mức: VDC
Khảo sát hệ thống túi khí an toàn chứa chất rắn $\mathrm{NaN_3}$ phân hủy tức thời giải phóng khí nitrogen ($\mathrm{N_2}$) và kim loại natri ($\mathrm{Na}$).
\choiceTF
{\True Khí nitrogen sinh ra là chất khí lí tưởng vì kích thước các phân tử $\mathrm{N_2}$ rất nhỏ so với khoảng cách trung bình giữa chúng.}
{Do Na tạo thành ở dạng tinh thể rắn mịn nên thể tích của Na trong túi khí lớn gấp nhiều lần thể tích của khí nitrogen giải phóng.}
{\True Nếu muốn giải phóng được $112\ \mathrm{l}$ khí nitrogen ở điều kiện chuẩn ($t = 0\ ^\circ\mathrm{C}$, $p_0 = 1,013 \cdot 10^5\ \mathrm{Pa}$), ta cần dùng tối thiểu $217\ \mathrm{g}$ chất rắn $\mathrm{NaN_3}$.}
{Biết một túi khí có dung tích $50\ \mathrm{l}$. Nếu nạp vào túi lượng chất rắn $\mathrm{NaN_3}$ đủ để sinh ra $2,0\ \mathrm{mol}$ khí $\mathrm{N_2}$ ở nhiệt độ $27\ ^\circ\mathrm{C}$ thì áp suất khí $\mathrm{N_2}$ trong túi phồng căng bằng $2,5 \cdot 10^5\ \mathrm{Pa}$.}
\loigiai{
\begin{itemchoice}
\item (Đúng) Khí nitrogen sinh ra là chất khí lí tưởng vì kích thước các phân tử $\mathrm{N_2}$ rất nhỏ so với khoảng cách trung bình giữa chúng. (Vì): Theo lý thuyết động học phân tử, thể tích riêng của các phân tử $\mathrm{N_2}$ là không đáng kể so với khoảng trống dãn nở giữa chúng. Mệnh đề đúng
\item (Đúng) Do Na tạo thành ở dạng tinh thể rắn mịn nên thể tích của Na trong túi khí lớn gấp nhiều lần thể tích của khí nitrogen giải phóng. (Vì): o khối lượng riêng của chất rắn Na lớn hơn hàng nghìn lần so với chất khí nên thể tích chất rắn Na tạo thành là cực kì nhỏ không đáng kể so với thể tích khí $\mathrm{N_2}$. Mệnh đề sai
\item (Đúng) Nếu muốn giải phóng được $112\ \mathrm{l}$ khí nitrogen ở điều kiện chuẩn ($t = 0\ ^\circ\mathrm{C}$, $p_0 = 1,013 \cdot 10^5\ \mathrm{Pa}$), ta cần dùng tối thiểu $217\ \mathrm{g}$ chất rắn $\mathrm{NaN_3}$. (Vì): Ở điều kiện tiêu chuẩn, số mol khí $\mathrm{N_2}$ cần sinh ra là $5,0\ \mathrm{mol}$. Khối lượng sodium azide cần phân hủy là:
\item (Sai) Biết một túi khí có dung tích $50\ \mathrm{l}$. Nếu nạp vào túi lượng chất rắn $\mathrm{NaN_3}$ đủ để sinh ra $2,0\ \mathrm{mol}$ khí $\mathrm{N_2}$ ở nhiệt độ $27\ ^\circ\mathrm{C}$ thì áp suất khí $\mathrm{N_2}$ trong túi phồng căng bằng $2,5 \cdot 10^5\ \mathrm{Pa}$. (Vì): Áp suất khí nitrogen dồn nén trong túi khí có thể tích $V = 50 \cdot 10^{-3}\ \mathrm{m^3}$ ở nhiệt độ $300\ \mathrm{K}$ bằng:
\end{itemchoice}
}
\end{ex}

% ===== Câu 142 =====
% ID: L12C2B11-36-TL
% Mức: VDC
\dangbt{Đồ thị và các quá trình biến đổi trạng thái liên tiếp}
\begin{bt}
% ID: L12C2B11-36-TL
% Mức: VDC
Trong dạng “Đồ thị và các quá trình biến đổi trạng thái liên tiếp”, undefined
\loigiai{
$V_2/V_1=(T_2/T_1)/(p_2/p_1)=1$.
}
\end{bt}

% ===== Câu 143 =====
% ID: L12C2B11-36-TLN
% Mức: VD
\dangbt{Xác định sự biến thiên khối lượng riêng của chất khí}
\begin{ex}
% ID: L12C2B11-36-TLN
% Mức: VD
Một bình kín dung tích $10\ \mathrm{l}$ chứa khí nitrogen ($\mathrm{N_2}$) ở nhiệt độ $27\ ^\circ\mathrm{C}$ dưới áp suất $2,493 \cdot 10^5\ \mathrm{Pa}$. Biết hằng số khí lí tưởng $R = 8,31\ \mathrm{J/(mol\cdot K)}$ và khối lượng mol của nitrogen là $28\ \mathrm{g/mol}$. Tính khối lượng riêng của khí nitrogen trong bình theo đơn vị $\mathrm{kg/m^3}$ (lấy kết quả làm tròn đến một chữ số thập phân).
\shortans{$2{,}8$}
\loigiai{
Áp dụng công thức
  \begin{aligned}
 $D&=\dfrac{p \cdot M}{R \cdot T}&=\dfrac{2,493 \cdot 10^5 \cdot 28 \cdot 10^{-3}}{8,31 \cdot (27 + 273)}&= 2,8\\mathrm{kg/m^3}$.
 \end{aligned}
}
\end{ex}

% ===== Câu 144 =====
% ID: L12C2B11-36-TN
% Mức: TH
\dangbt{Xác định số mol, khối lượng và khối lượng riêng của khí}
\begin{ex}
% ID: L12C2B11-36-TN
% Mức: TH
Trong dạng “Xác định số mol, khối lượng và khối lượng riêng của khí”, Trong phương trình khí lí tưởng, nhiệt độ phải dùng
\choice
{\True kelvin}
{độ $C$}
{độ $F$}
{thang tùy ý}
\loigiai{
Đáp án A. $T$ phải là nhiệt độ tuyệt đối.
}
\end{ex}

% ===== Câu 145 =====
% ID: L12C2B11-37-DS
% Mức: VDC
\dangbt{Xác định sự biến thiên khối lượng riêng của chất khí}
\begin{ex}
% ID: L12C2B11-37-DS
% Mức: VDC
Một túi khí ô tô có dung tích không đổi $48\ \mathrm{l}$ khi phồng căng. Sau va chạm, phản ứng phân hủy $\mathrm{NaN_3}$ giải phóng khí nitrogen làm áp suất khí trong túi đạt $1,21 \cdot 10^5\ \mathrm{Pa}$ ở nhiệt độ $30\ ^\circ\mathrm{C}$.
\choiceTF
{\True Trạng thái cân bằng nhiệt động của lượng khí trong túi khí được xác định bởi bốn đại lượng là khối lượng, thể tích, áp suất và nhiệt độ.}
{\True Số mol khí nitrogen được tích vào túi khí khi phồng căng bằng $2,31\ \mathrm{mol}$.}
{Khối lượng chất rắn $\mathrm{NaN_3}$ cần thiết tham gia phản ứng phân hủy hoàn toàn để đạt áp suất trên xấp xỉ bằng $150\ \mathrm{g}$.}
{\True Thể tích của lượng khí $\mathrm{N_2}$ này nếu ở áp suất khí quyển $1,0 \cdot 10^5\ \mathrm{Pa}$ và nhiệt độ $27\ ^\circ\mathrm{C}$ là $57,6\ \mathrm{l}$.}
\loigiai{
\begin{itemchoice}
\item (Đúng) Trạng thái cân bằng nhiệt động của lượng khí trong túi khí được xác định bởi bốn đại lượng là khối lượng, thể tích, áp suất và nhiệt độ. (Vì): ây là các đại lượng đặc trưng cho trạng thái của một khối chất khí xác định đựng trong bình kín. Mệnh đề đúng
\item (Đúng) Số mol khí nitrogen được tích vào túi khí khi phồng căng bằng $2,31\ \mathrm{mol}$. (Vì): Áp dụng phương trình trạng thái khí để tìm số mol với $V = 48 \cdot 10^{-3}\ \mathrm{m^3}$ và $T = 303\ \mathrm{K}$:
\item (Sai) Khối lượng chất rắn $\mathrm{NaN_3}$ cần thiết tham gia phản ứng phân hủy hoàn toàn để đạt áp suất trên xấp xỉ bằng $150\ \mathrm{g}$. (Vì): Từ số mol khí nitrogen trên, khối lượng sodium azide cần phân hủy là:
\item (Đúng) Thể tích của lượng khí $\mathrm{N_2}$ này nếu ở áp suất khí quyển $1,0 \cdot 10^5\ \mathrm{Pa}$ và nhiệt độ $27\ ^\circ\mathrm{C}$ là $57,6\ \mathrm{l}$. (Vì): Từ kết quả số mol ở câu trên, thể tích khí nitrogen thu được ở điều kiện thông thường là:
\end{itemchoice}
}
\end{ex}

% ===== Câu 146 =====
% ID: L12C2B11-37-TLN
% Mức: VD
\begin{ex}
% ID: L12C2B11-37-TLN
% Mức: VD
\nguon{Dựa trên hoạt động 2c trang 47 Sách giáo khoa Vật lí 12 Kết nối tri thức với cuộc sống}
	Khi xảy ra va chạm, phản ứng phân hủy chất rắn sodium azide giải phóng ra $2,4\ \mathrm{mol}$ khí nitrogen ($\mathrm{N_2}$) làm phồng căng túi khí đạt dung tích $48\ \mathrm{l}$ ở nhiệt độ $27\ ^\circ\mathrm{C}$. Tính áp suất của khí nitrogen trong túi theo đơn vị $10^5\ \mathrm{Pa}$. Lấy hằng số khí lí tưởng $R = 8,31\ \mathrm{J/(mol\cdot K)}$ (làm tròn kết quả đến một chữ số thập phân).
\shortans{$1{,}2$}
\loigiai{
Áp dụng công thức
  \begin{aligned}
 p &= $\dfrac{n \cdot R \cdot T}{V}&=\dfrac{2,4 \cdot 8,31 \cdot (27 + 273)}{48 \cdot 10^{-3}}&= 124650\\mathrm{Pa}\approx$ 1,2 \cdot 10^5 $\\mathrm{Pa}$.
 \end{aligned}
}
\end{ex}

% ===== Câu 147 =====
% ID: L12C2B11-37-TN
% Mức: VD
\dangbt{Xác định số mol, khối lượng và khối lượng riêng của khí}
\begin{ex}
% ID: L12C2B11-37-TN
% Mức: VD
Trong dạng “Xác định số mol, khối lượng và khối lượng riêng của khí”, Ở thể tích không đổi, áp suất của lượng khí xác định tỉ lệ
\choice
{\True thuận với nhiệt độ tuyệt đối}
{nghịch với nhiệt độ}
{thuận với bình phương nhiệt độ}
{không phụ thuộc nhiệt độ}
\loigiai{
Đáp án A. Từ $pV=nRT$, khi $V$ cố định thì $p/T$ không đổi.
}
\end{ex}

% ===== Câu 148 =====
% ID: L12C2B11-38-DS
% Mức: VDC
\dangbt{Xác định sự biến thiên khối lượng riêng của chất khí}
\begin{ex}
% ID: L12C2B11-38-DS
% Mức: VDC
Vỏ của một loại bóng thám không vô tuyến (Radiosonde) làm bằng cao su tự nhiên có độ dày ban đầu khi ở mặt đất là $0,051\ \mathrm{mm}$, đường kính vỏ bóng là $1,5\ \mathrm{m}$. Khi bóng bay lên độ cao trên $30\ \mathrm{km}$, đường kính của bóng tăng lên gấp $3\ \text{lần}$ ban đầu.
\choiceTF
{\True Cao su tự nhiên là chất rắn vô định hình có tính đàn hồi cao được dùng làm vỏ bóng thám không.}
{Khi đường kính bóng thám không tăng lên gấp $3\ \text{lần}$ thì thể tích của bóng tăng lên gấp $9\ \text{lần}$.}
{Co thể tích cao su của vỏ bóng không đổi trong quá trình dãn nở, độ dày vỏ bóng ở độ cao trên xấp xỉ bằng $0,017\ \mathrm{mm}$.}
{\True Biết bóng được bơm bằng khí hydrogen. Khi vừa thả ở mặt đất có áp suất $1,0 \cdot 10^5\ \mathrm{Pa}$ và nhiệt độ $300\ \mathrm{K}$. Khi bóng đạt độ cao trên có nhiệt độ tuyệt đối giảm $40\%$, áp suất của khí bên trong bóng bằng $0,022 \cdot 10^5\ \mathrm{Pa}$.}
\loigiai{
\begin{itemchoice}
\item (Đúng) Cao su tự nhiên là chất rắn vô định hình có tính đàn hồi cao được dùng làm vỏ bóng thám không. (Vì): Cao su tự nhiên có cấu trúc nguyên tử hỗn độn, không có cấu trúc tinh thể nên là chất rắn vô định hình dẻo dai đàn hồi tốt. Mệnh đề đúng
\item (Sai) Khi đường kính bóng thám không tăng lên gấp $3\ \text{lần}$ thì thể tích của bóng tăng lên gấp $9\ \text{lần}$. (Vì): Thể tích hình cầu tỉ lệ với lập phương đường kính nên thể tích dãn rộng tăng gấp $3^3 = 27\ \text{lần}$ so với ban đầu. Mệnh đề sai
\item (Đúng) Co thể tích cao su của vỏ bóng không đổi trong quá trình dãn nở, độ dày vỏ bóng ở độ cao trên xấp xỉ bằng $0,017\ \mathrm{mm}$. (Vì): o diện tích bề mặt bóng $S = 4\pi R^2$ tăng lên gấp $3^2 = 9\ \text{lần}$, để thể tích cao su của vỏ bóng $V = S \cdot d$ không đổi thì độ dày $d$ phải giảm đi $9\ \text{lần}$, tức là bằng $0,051 / 9 \approx 0,0057\ \mathrm{mm}$ chứ không phải $0,017\ \mathrm{mm}$. Mệnh đề sai
\item (Đúng) Biết bóng được bơm bằng khí hydrogen. Khi vừa thả ở mặt đất có áp suất $1,0 \cdot 10^5\ \mathrm{Pa}$ và nhiệt độ $300\ \mathrm{K}$. Khi bóng đạt độ cao trên có nhiệt độ tuyệt đối giảm $40\%$, áp suất của khí bên trong bóng bằng $0,022 \cdot 10^5\ \mathrm{Pa}$. (Vì): Khi nhiệt độ tuyệt đối giảm $40\%$ ($T_2 = 0,6 T_1$), thể tích tăng $27\ \text{lần}$ ($V_2 = 27 V_1$), áp suất bên trong bóng là:
\end{itemchoice}
}
\end{ex}

% ===== Câu 149 =====
% ID: L12C2B11-38-TLN
% Mức: VD
\begin{ex}
% ID: L12C2B11-38-TLN
% Mức: VD
Nén một lượng khí lí tưởng xác định trong xilanh sao cho thể tích giảm đi $2\ \text{lần}$ và nhiệt độ tuyệt đối tăng thêm $20\%$. Tính tỉ số giữa áp suất lúc sau $p_2$ và áp suất ban đầu $p_1$.
\shortans{$2{,}4$}
\loigiai{
Áp dụng công thức
  \begin{aligned}
 $\dfrac{p_2}{p_1}&=\dfrac{V_1}{V_2}\cdot\dfrac{T_2}{T_1}&= 2\cdot$ 1,2 
&= 2,4.
 \end{aligned}
}
\end{ex}

% ===== Câu 150 =====
% ID: L12C2B11-38-TN
% Mức: VD
\dangbt{Xác định số mol, khối lượng và khối lượng riêng của khí}
\begin{ex}
% ID: L12C2B11-38-TN
% Mức: VD
Trong dạng “Xác định số mol, khối lượng và khối lượng riêng của khí”, Số mol khí được tính từ khối lượng m và khối lượng mol $M$ bằng
\choice
{\True $n=m/M$}
{$n=mM$}
{$n=M/m$}
{$n=m+M$}
\loigiai{
Đáp án A. Theo định nghĩa số mol, $n=m/M$.
}
\end{ex}

% ===== Câu 151 =====
% ID: L12C2B11-39-DS
% Mức: VDC
\dangbt{Xác định sự biến thiên khối lượng riêng của chất khí}
\begin{ex}
% ID: L12C2B11-39-DS
% Mức: VDC
Một bóng thám thính khí tượng được bơm $80\ \mathrm{g}$ khí hydrogen ($\mathrm{H_2}$) có khối lượng mol $2\ \mathrm{g/mol}$. Ở mặt đất, áp suất không khí là $1,0 \cdot 10^5\ \mathrm{Pa}$ và nhiệt độ $27\ ^\circ\mathrm{C}$.
\choiceTF
{\True Số mol khí hydrogen chứa trong bóng thám không bằng $40\ \mathrm{mol}$.}
{\True Thể tích của bóng thám không ngay khi vừa bơm xong ở mặt đất xấp xỉ bằng $1,0\ \mathrm{m^3}$.}
{Khi bóng bay lên độ cao có nhiệt độ tuyệt đối giảm đi $10\%$ và áp suất khí quyển giảm còn $0,25 \cdot 10^5\ \mathrm{Pa}$ thì thể tích của bóng tăng lên đạt giá trị $4,5\ \mathrm{m^3}$.}
{\True Biết vỏ bóng chỉ chịu được thể tích tối đa là $5,0\ \mathrm{m^3}$. Nếu nhiệt độ tuyệt đối ở độ cao vỡ giảm còn $240\ \mathrm{K}$ thì bóng sẽ bị vỡ khi áp suất khí quyển xung quanh giảm dưới $0,16 \cdot 10^5\ \mathrm{Pa}$.}
\loigiai{
\begin{itemchoice}
\item (Sai) Số mol khí hydrogen chứa trong bóng thám không bằng $40\ \mathrm{mol}$. (Vì): ố mol khí hydrogen bơm vào bóng là:
\item (Đúng) Thể tích của bóng thám không ngay khi vừa bơm xong ở mặt đất xấp xỉ bằng $1,0\ \mathrm{m^3}$. (Vì): Áp dụng phương trình Clapeyron để xác định thể tích khí lúc bơm xong với $T_1 = 300\ \mathrm{K}$:
\item (Sai) Khi bóng bay lên độ cao có nhiệt độ tuyệt đối giảm đi $10\%$ và áp suất khí quyển giảm còn $0,25 \cdot 10^5\ \mathrm{Pa}$ thì thể tích của bóng tăng lên đạt giá trị $4,5\ \mathrm{m^3}$. (Vì): Tại độ cao có $T_2 = 0,9 T_1 = 270\ \mathrm{K}$ và $p_2 = 0,25 \cdot 10^5\ \mathrm{Pa}$, thể tích bóng là:
\item (Đúng) Biết vỏ bóng chỉ chịu được thể tích tối đa là $5,0\ \mathrm{m^3}$. Nếu nhiệt độ tuyệt đối ở độ cao vỡ giảm còn $240\ \mathrm{K}$ thì bóng sẽ bị vỡ khi áp suất khí quyển xung quanh giảm dưới $0,16 \cdot 10^5\ \mathrm{Pa}$. (Vì): Khi đạt thể tích tối đa $V_{\text{max}} = 5,0\ \mathrm{m^3}$ ở nhiệt độ $T_{\text{vỡ}} = 240\ \mathrm{K}$, áp suất khí quyển vỡ bằng:
\end{itemchoice}
}
\end{ex}

% ===== Câu 152 =====
% ID: L12C2B11-39-TLN
% Mức: VD
\begin{ex}
% ID: L12C2B11-39-TLN
% Mức: VD
Một bình thép chứa $4\ \mathrm{g}$ khí helium ($\mathrm{He}$, khối lượng mol $M = 4\ \mathrm{g/mol}$) ở nhiệt độ $27\ ^\circ\mathrm{C}$ dưới áp suất $1,2465 \cdot 10^5\ \mathrm{Pa}$. Tính thể tích của bình thép này theo đơn vị lít ($\mathrm{l}$). Lấy hằng số khí lí tưởng $R = 8,31\ \mathrm{J/(mol\cdot K)}$.
\shortans{$20$}
\loigiai{
Áp dụng công thức
  \begin{aligned}
 $V&=\dfrac{m \cdot R \cdot T}{M \cdot p}&=\dfrac{4 \cdot 10^{-3} \cdot 8,31 \cdot (27 + 273)}{4 \cdot 10^{-3} \cdot 1,2465 \cdot 10^5}&= 0,02\\mathrm{m^3}= 20\\mathrm{l}$.
 \end{aligned}
}
\end{ex}

% ===== Câu 153 =====
% ID: L12C2B11-39-TN
% Mức: VD
\dangbt{Xác định số mol, khối lượng và khối lượng riêng của khí}
\begin{ex}
% ID: L12C2B11-39-TN
% Mức: VD
Trong dạng “Xác định số mol, khối lượng và khối lượng riêng của khí”, Khối lượng riêng của khí lí tưởng thỏa
\choice
{\True $\rho=pM/(RT)$}
{$\rho=RT/(pM)$}
{$\rho=pRT/M$}
{$\rho=p/(MRT)$}
\loigiai{
Đáp án A. Từ $pV=(m/M)RT$ và $\rho=m/V$.
}
\end{ex}

% ===== Câu 154 =====
% ID: L12C2B11-40-DS
% Mức: VDC
\dangbt{Xác định sự biến thiên khối lượng riêng của chất khí}
\begin{ex}
% ID: L12C2B11-40-DS
% Mức: VDC
Một quả bóng thám không được chế tạo bằng vật liệu cao su co dãn, ban đầu được bơm khí helium ở gần mặt đất đạt thể tích $10\ \mathrm{m^3}$ ở áp suất $1,0 \cdot 10^5\ \mathrm{Pa}$ và nhiệt độ $300\ \mathrm{K}$. Khi bóng bay lên tầng khí quyển cao có áp suất khí quyển đo được là $0,25 \cdot 10^5\ \mathrm{Pa}$ và nhiệt độ khí quyển tại đó giảm còn $225\ \mathrm{K}$. Nhận định nào dưới đây về bóng thám không là đúng hay sai?
\choiceTF
{\True Bóng thám không bay lên cao được là nhờ lực đẩy Archimede của không khí xung quanh tác dụng lên bóng lớn hơn trọng lượng của toàn bộ quả bóng.}
{Khi quả bóng bay lên cao, nhiệt độ giảm làm các phân tử khí chuyển động chậm lại, làm giảm áp suất khí bên trong, khiến thể tích bóng tự động co nhỏ lại.}
{\True Tại độ cao mà áp suất khí quyển là $0,25 \cdot 10^5\ \mathrm{Pa}$ và nhiệt độ là $225\ \mathrm{K}$, thể tích của quả bóng thám không bằng $30\ \mathrm{m^3}$.}
{\True Tỉ số giữa bán kính của quả bóng ở độ cao trên so với bán kính khi vừa bơm xong ở mặt đất xấp xỉ bằng $1,44\ \text{lần}$.}
\loigiai{
\begin{itemchoice}
\item (Đúng) Bóng thám không bay lên cao được là nhờ lực đẩy Archimede của không khí xung quanh tác dụng lên bóng lớn hơn trọng lượng của toàn bộ quả bóng. (Vì): Theo định luật Archimede trong khí quyển học, bóng thám không dâng lên cao vì lực đẩy Archimede hướng lên lớn hơn trọng lực tổng cộng của bóng. Mệnh đề đúng
\item (Sai) Khi quả bóng bay lên cao, nhiệt độ giảm làm các phân tử khí chuyển động chậm lại, làm giảm áp suất khí bên trong, khiến thể tích bóng tự động co nhỏ lại. (Vì): Mặc dù nhiệt độ giảm làm giảm áp suất, nhưng áp suất không khí bên ngoài giảm rất mạnh (giảm tới 4 lần) khi lên cao, tạo ra sự chênh lệch lớn khiến khí bên trong dãn nở đẩy vỏ bóng to ra. Mệnh đề sai
\item (Đúng) Tại độ cao mà áp suất khí quyển là $0,25 \cdot 10^5\ \mathrm{Pa}$ và nhiệt độ là $225\ \mathrm{K}$, thể tích của quả bóng thám không bằng $30\ \mathrm{m^3}$. (Vì): Áp dụng phương trình trạng thái của khí lí tưởng cho lượng khí helium chứa trong bóng thám không:
\item (Đúng) Tỉ số giữa bán kính của quả bóng ở độ cao trên so với bán kính khi vừa bơm xong ở mặt đất xấp xỉ bằng $1,44\ \text{lần}$. (Vì): Thể tích của bóng hình cầu tỉ lệ thuận với lập phương bán kính $V = \dfrac{4}{3}\pi R^3$, tỉ số bán kính là:
\end{itemchoice}
}
\end{ex}

% ===== Câu 155 =====
% ID: L12C2B11-40-TLN
% Mức: VD
\begin{ex}
% ID: L12C2B11-40-TLN
% Mức: VD
Một bình dung tích $8,31\ \mathrm{l}$ chứa $4\ \mathrm{g}$ một chất khí lí tưởng ở nhiệt độ $27\ ^\circ\mathrm{C}$ dưới áp suất $3,0 \cdot 10^5\ \mathrm{Pa}$. Tính khối lượng mol $M$ của chất khí này theo đơn vị $\mathrm{g/mol}$. Lấy hằng số khí lí tưởng $R = 8,31\ \mathrm{J/(mol\cdot K)}$.
\shortans{$4$}
\loigiai{
Áp dụng công thức
  \begin{aligned}
 $M&=\dfrac{m \cdot R \cdot T}{p \cdot V}&=\dfrac{4 \cdot 10^{-3} \cdot 8,31 \cdot (27 + 273)}{3,0 \cdot 10^5 \cdot 8,31 \cdot 10^{-3}}&= 4\cdot10^{-3}\\mathrm{kg/mol}= 4\\mathrm{g/mol}$.
 \end{aligned}
}
\end{ex}

% ===== Câu 156 =====
% ID: L12C2B11-40-TN
% Mức: VDC
\dangbt{Xác định số mol, khối lượng và khối lượng riêng của khí}
\begin{ex}
% ID: L12C2B11-40-TN
% Mức: VDC
Trong dạng “Xác định số mol, khối lượng và khối lượng riêng của khí”, Hai bình cùng p, $V$, $T$ chứa khí lí tưởng khác loại thì có
\choice
{\True cùng số phân tử}
{cùng khối lượng}
{cùng khối lượng riêng}
{khác số mol}
\loigiai{
Đáp án A. Theo $n=pV/(RT)$, số mol và số phân tử bằng nhau.
}
\end{ex}

% ===== Câu 157 =====
% ID: L12C2B11-41-DS
% Mức: VDC
\dangbt{Xác định sự biến thiên khối lượng riêng của chất khí}
\begin{ex}
% ID: L12C2B11-41-DS
% Mức: VDC
Khảo sát sự biến thiên khối lượng riêng của một khối lượng khí oxygen ($\mathrm{O_2}$, khối lượng mol $32\ \mathrm{g/mol}$) đựng trong một xi lanh có pit-tông di chuyển được. Ban đầu khối khí ở áp suất $1,0 \cdot 10^5\ \mathrm{Pa}$ và nhiệt độ $27\ ^\circ\mathrm{C}$.
\choiceTF
{\True Khi di chuyển pit-tông để thể tích khối khí dãn rộng ra ở nhiệt độ không đổi, khối lượng riêng của khí giảm đi.}
{\True Khối lượng riêng của oxygen ở trạng thái ban đầu xấp xỉ bằng $1,28\ \mathrm{kg/m^3}$.}
{Nếu nén khí để áp suất tăng lên đến $2,0 \cdot 10^5\ \mathrm{Pa}$ và nhiệt độ tăng lên đến $127\ ^\circ\mathrm{C}$ thì khối lượng riêng lúc này của khí đạt giá trị bằng $2,56\ \mathrm{kg/m^3}$.}
{\True Trong quá trình biến đổi đẳng áp, nếu muốn khối lượng riêng của khí giảm đi một nửa so với ban đầu thì ta cần đun nóng khí từ $27\ ^\circ\mathrm{C}$ lên đến $327\ ^\circ\mathrm{C}$.}
\loigiai{
\begin{itemchoice}
\item (Đúng) Khi di chuyển pit-tông để thể tích khối khí dãn rộng ra ở nhiệt độ không đổi, khối lượng riêng của khí giảm đi. (Vì): Khi thể tích dãn ra tại nhiệt độ không đổi, mật độ phân tử khí giảm xuống làm khối lượng riêng của khối khí giảm đi. Mệnh đề đúng
\item (Đúng) Khối lượng riêng của oxygen ở trạng thái ban đầu xấp xỉ bằng $1,28\ \mathrm{kg/m^3}$. (Vì): Áp dụng công thức khối lượng riêng chất khí với $T_1 = 300\ \mathrm{K}$:
\item (Sai) Nếu nén khí để áp suất tăng lên đến $2,0 \cdot 10^5\ \mathrm{Pa}$ và nhiệt độ tăng lên đến $127\ ^\circ\mathrm{C}$ thì khối lượng riêng lúc này của khí đạt giá trị bằng $2,56\ \mathrm{kg/m^3}$. (Vì): Tại trạng thái sau với $p_2 = 2,0 \cdot 10^5\ \mathrm{Pa}$ và $T_2 = 400\ \mathrm{K}$, khối lượng riêng là:
\item (Đúng) Trong quá trình biến đổi đẳng áp, nếu muốn khối lượng riêng của khí giảm đi một nửa so với ban đầu thì ta cần đun nóng khí từ $27\ ^\circ\mathrm{C}$ lên đến $327\ ^\circ\mathrm{C}$. (Vì): ể khối lượng riêng đẳng áp giảm một nửa ($D_2 = 0,5 D_1$), nhiệt độ tuyệt đối cần tăng gấp hai lần: $T_2 = 2 T_1 = 600\ \mathrm{K} \Rightarrow t_2 = 327\ ^\circ\mathrm{C}$. Mệnh đề đúng
\end{itemchoice}
}
\end{ex}

% ===== Câu 158 =====
% ID: L12C2B11-41-TLN
% Mức: VD
\begin{ex}
% ID: L12C2B11-41-TLN
% Mức: VD
Một bình chứa khí oxygen ($\mathrm{O_2}$) dung tích $16,62\ \mathrm{l}$ ở nhiệt độ $27\ ^\circ\mathrm{C}$ dưới áp suất $3,0 \cdot 10^5\ \mathrm{Pa}$. Tính số mol khí oxygen chứa trong bình. Lấy hằng số khí lí tưởng $R = 8,31\ \mathrm{J/(mol\cdot K)}$.
\shortans{$2$}
\loigiai{
Áp dụng công thức
  \begin{aligned}
 n &= $\dfrac{p \cdot V}{R \cdot T}&=\dfrac{3,0 \cdot 10^5 \cdot 16,62 \cdot 10^{-3}}{8,31 \cdot (27 + 273)}&= 2\\mathrm{mol}$.
 \end{aligned}
}
\end{ex}

% ===== Câu 159 =====
% ID: L12C2B11-41-TN
% Mức: VDC
\dangbt{Xác định số mol, khối lượng và khối lượng riêng của khí}
\begin{ex}
% ID: L12C2B11-41-TN
% Mức: VDC
Trong dạng “Xác định số mol, khối lượng và khối lượng riêng của khí”, Khi giải bài toán khí lí tưởng, áp suất đưa vào phương trình phải là
\choice
{\True áp suất tuyệt đối}
{áp suất dư}
{áp suất do chất lỏng בלבד}
{áp suất tương đối tùy ý}
\loigiai{
Đáp án A. Phương trình trạng thái sử dụng áp suất tuyệt đối.
}
\end{ex}

% ===== Câu 160 =====
% ID: L12C2B11-42-DS
% Mức: VDC
\dangbt{Xác định sự biến thiên khối lượng riêng của chất khí}
\begin{ex}
% ID: L12C2B11-42-DS
% Mức: VDC
Một bình dung tích $50\ \mathrm{l}$ chứa khí hydrogen ($\mathrm{H_2}$) ở nhiệt độ $17\ ^\circ\mathrm{C}$ và áp suất $20\ \mathrm{atm}$. Biết khối lượng mol của hydrogen là $2\ \mathrm{g/mol}$, $1\ \mathrm{atm} \approx 1,013 \cdot 10^5\ \mathrm{Pa}$.
\choiceTF
{\True Khí hydrogen gồm các phân tử chuyển động hỗn loạn không ngừng va chạm đàn hồi hoàn toàn với nhau và với thành bình.}
{\True Nhiệt độ tuyệt đối của khối khí bằng $290\ \mathrm{K}$.}
{\True Khối lượng khí hydrogen có trong bình bằng $84\ \mathrm{g}$.}
{Nếu ta nạp thêm $42\ \mathrm{g}$ khí hydrogen vào bình và giữ nhiệt độ không đổi thì áp suất trong bình tăng lên đến $40\ \mathrm{atm}$.}
\loigiai{
\begin{itemchoice}
\item (Đúng) Khí hydrogen gồm các phân tử chuyển động hỗn loạn không ngừng va chạm đàn hồi hoàn toàn với nhau và với thành bình. (Vì): ây là các giả thuyết cơ bản của mô hình động học phân tử đối với khối khí lí tưởng. Mệnh đề đúng
\item (Đúng) Nhiệt độ tuyệt đối của khối khí bằng $290\ \mathrm{K}$. (Vì): ổi nhiệt độ Celsius sang thang đo tuyệt đối Kelvin: $T = 17 + 273 = 290\ \mathrm{K}$. Mệnh đề đúng
\item (Đúng) Khối lượng khí hydrogen có trong bình bằng $84\ \mathrm{g}$. (Vì): Áp dụng phương trình Clapeyron để tính khối lượng khí hydrogen:
\item (Sai) Nếu ta nạp thêm $42\ \mathrm{g}$ khí hydrogen vào bình và giữ nhiệt độ không đổi thì áp suất trong bình tăng lên đến $40\ \mathrm{atm}$. (Vì): Khi nạp thêm $42\ \mathrm{g}$ khí hydrogen, tổng khối lượng đạt $84 + 42 = 126\ \mathrm{g}$ (tăng gấp $1,5\ \text{lần}$ ban đầu). Do áp suất tỉ lệ thuận với khối lượng nên áp suất lúc sau tăng gấp $1,5\ \text{lần}$ đạt $30\ \mathrm{atm}$ chứ không phải $40\ \mathrm{atm}$. Mệnh đề sai
\end{itemchoice}
}
\end{ex}

% ===== Câu 161 =====
% ID: L12C2B11-42-TLN
% Mức: VD
\begin{ex}
% ID: L12C2B11-42-TLN
% Mức: VD
Một khối khí lí tưởng biến đổi trạng thái từ trạng thái $1$ ($p_1 = 1,0 \cdot 10^5\ \mathrm{Pa}$, $V_1 = 3\ \mathrm{l}$, $T_1 = 300\ \mathrm{K}$) sang trạng thái $2$ có $V_2 = 1\ \mathrm{l}$ và $T_2 = 450\ \mathrm{K}$. Tính áp suất $p_2$ của khối khí theo đơn vị $10^5\ \mathrm{Pa}$.
\shortans{$4{,}5$}
\loigiai{
Áp dụng công thức
  \begin{aligned}
 p_2 &= p_1 $\cdot\dfrac{V_1}{V_2}\cdot\dfrac{T_2}{T_1}&= 1,0\cdot10^5\cdot\dfrac{3}{1}\cdot\dfrac{450}{300}&= 4,5\cdot10^5\\mathrm{Pa}$.
 \end{aligned}
}
\end{ex}

% ===== Câu 162 =====
% ID: L12C2B11-42-TN
% Mức: VD
\begin{ex}
% ID: L12C2B11-42-TN
% Mức: VD
Một quả bóng thám không khi vừa được bơm xong ở mặt đất có thể tích $12\ \mathrm{m^3}$ ở áp suất $1,0 \cdot 10^5\ \mathrm{Pa}$ và nhiệt độ $27\ ^\circ\mathrm{C}$. Khi quả bóng thám không bay lên tầng khí quyển cao có áp suất giảm còn $0,30 \cdot 10^5\ \mathrm{Pa}$ thì thể tích của bóng tăng lên đạt $32\ \mathrm{m^3}$. Nhiệt độ của tầng khí quyển lúc này bằng bao nhiêu kelvin?
\choice
{$220\ \mathrm{K}$}
{\True $240\ \mathrm{K}$}
{$260\ \mathrm{K}$}
{$280\ \mathrm{K}$}
\loigiai{
Áp dụng công thức
  \begin{aligned}
 T_2 &= T_1 $\cdot\dfrac{p_2 \cdot V_2}{p_1 \cdot V_1}&= (27 + 273)\cdot\dfrac{0,30 \cdot 10^5 \cdot 32}{1,0 \cdot 10^5 \cdot 12}&= 240\\mathrm{K}$.
 \end{aligned} 
 Chọn đáp án tương ứng.
}
\end{ex}

% ===== Câu 163 =====
% ID: L12C2B11-43-DS
% Mức: VDC
\begin{ex}
% ID: L12C2B11-43-DS
% Mức: VDC
Một lượng khí lí tưởng ở trạng thái ban đầu có thể tích $3,0\ \mathrm{l}$, áp suất $2,0 \cdot 10^5\ \mathrm{Pa}$ và nhiệt độ $27\ ^\circ\mathrm{C}$. Người ta đun nóng đẳng tích khối khí này để nhiệt độ tăng thêm $100\ ^\circ\mathrm{C}$, sau đó dãn nở đẳng nhiệt để thể tích tăng gấp đôi.
\choiceTF
{\True Quá trình đẳng tích là quá trình biến đổi trạng thái trong đó thể tích khối khí được giữ không đổi.}
{Sau quá trình đun nóng đẳng tích, áp suất của khối khí tăng thêm một lượng bằng $2,0 \cdot 10^5\ \mathrm{Pa}$.}
{\True Tại trạng thái cuối cùng, áp suất của khối khí bằng $1,33 \cdot 10^5\ \mathrm{Pa}$.}
{Tích số của áp suất và thể tích ở trạng thái cuối cùng lớn gấp hai lần tích số của chúng ở trạng thái đầu tiên.}
\loigiai{
\begin{itemchoice}
\item (Đúng) Quá trình đẳng tích là quá trình biến đổi trạng thái trong đó thể tích khối khí được giữ không đổi. (Vì): Theo định nghĩa, quá trình đẳng tích là quá trình biến đổi trạng thái khi thể tích của khối khí được giữ cố định. Mệnh đề đúng
\item (Sai) Sau quá trình đun nóng đẳng tích, áp suất của khối khí tăng thêm một lượng bằng $2,0 \cdot 10^5\ \mathrm{Pa}$. (Vì): Nhiệt độ tuyệt đối ban đầu $T_1 = 300\ \mathrm{K}$. Đun nóng đẳng tích đến nhiệt độ $T_2 = 300 + 100 = 400\ \mathrm{K}$. Áp suất tương ứng là:
\item (Đúng) Tại trạng thái cuối cùng, áp suất của khối khí bằng $1,33 \cdot 10^5\ \mathrm{Pa}$. (Vì): Quá trình tiếp theo dãn đẳng nhiệt ở nhiệt độ $400\ \mathrm{K}$ làm thể tích tăng gấp đôi ($V_3 = 2 V_2$). Áp suất cuối cùng là:
\item (Sai) Tích số của áp suất và thể tích ở trạng thái cuối cùng lớn gấp hai lần tích số của chúng ở trạng thái đầu tiên. (Vì): Từ các kết quả trên, ta tính tích số ở hai trạng thái:
\end{itemchoice}
}
\end{ex}

% ===== Câu 164 =====
% ID: L12C2B11-43-TLN
% Mức: VD
\begin{ex}
% ID: L12C2B11-43-TLN
% Mức: VD
Một khối khí lí tưởng xác định ở nhiệt độ $27\ ^\circ\mathrm{C}$ có khối lượng riêng là $1,2\ \mathrm{kg/m^3}$. Nếu đun nóng đẳng áp khối khí này đến nhiệt độ $127\ ^\circ\mathrm{C}$ thì khối lượng riêng của khối khí lúc này bằng bao nhiêu $\mathrm{kg/m^3}$?
\shortans{$0{,}9$}
\loigiai{
Áp dụng công thức
  \begin{aligned}
 D_2 &= D_1 $\cdot\dfrac{T_1}{T_2}&= 1,2\cdot\dfrac{27 + 273}{127 + 273}&= 0,9\\mathrm{kg/m^3}$.
 \end{aligned}
}
\end{ex}

% ===== Câu 165 =====
% ID: L12C2B11-43-TN
% Mức: VD
\begin{ex}
% ID: L12C2B11-43-TN
% Mức: VD
Một bình thép chứa $16\ \mathrm{g}$ khí oxygen ($\mathrm{O_2}$) ở nhiệt độ $27\ ^\circ\mathrm{C}$ và áp suất $1,5 \cdot 10^5\ \mathrm{Pa}$. Biết khối lượng mol của oxygen là $32\ \mathrm{g/mol}$. Thể tích của bình thép này bằng bao nhiêu lít?
\choice
{$6,25\ \mathrm{l}$}
{\True $8,31\ \mathrm{l}$}
{$12,5\ \mathrm{l}$}
{$16,6\ \mathrm{l}$}
\loigiai{
Áp dụng công thức
  \begin{aligned}
 $V&=\dfrac{m \cdot R \cdot T}{M \cdot p}&=\dfrac{16 \cdot 10^{-3} \cdot 8,31 \cdot (27 + 273)}{32 \cdot 10^{-3} \cdot 1,5 \cdot 10^5}&= 8,31\cdot10^{-3}\\mathrm{m^3}= 8,31\\mathrm{l}$.
 \end{aligned} 
 Chọn đáp án tương ứng.
}
\end{ex}

% ===== Câu 166 =====
% ID: L12C2B11-44-DS
% Mức: VDC
\begin{ex}
% ID: L12C2B11-44-DS
% Mức: VDC
Một bình dung tích $10\ \mathrm{l}$ chứa khí nitrogen ($\mathrm{N_2}$) ở nhiệt độ $27\ ^\circ\mathrm{C}$ và áp suất $2,493 \cdot 10^5\ \mathrm{Pa}$. Biết hằng số khí $R \approx 8,31\ \mathrm{J/(mol\cdot K)}$ và khối lượng mol của nitrogen là $28\ \mathrm{g/mol}$.
\choiceTF
{Khối lượng riêng của không khí luôn bằng hằng số không phụ thuộc vào nhiệt độ và áp suất khí quyển.}
{\True Công thức xác định khối lượng riêng của khí lí tưởng từ phương trình Clapeyron là $D = \dfrac{p \cdot M}{R \cdot T}$.}
{\True Khối lượng riêng của khí nitrogen trong bình ở trạng thái trên bằng $2,8\ \mathrm{kg/m^3}$.}
{Nếu nhiệt độ trong bình tăng lên đến $127\ ^\circ\mathrm{C}$ trong khi thể tích bình giữ nguyên không đổi thì khối lượng riêng của khí nitrogen trong bình giảm đi $25\%$.}
\loigiai{
\begin{itemchoice}
\item (Sai) Khối lượng riêng của không khí luôn bằng hằng số không phụ thuộc vào nhiệt độ và áp suất khí quyển. (Vì): Khối lượng riêng của không khí thay đổi mạnh mẽ khi áp suất và nhiệt độ biến thiên ở các điều kiện địa lý độ cao khác nhau. Mệnh đề sai
\item (Đúng) Công thức xác định khối lượng riêng của khí lí tưởng từ phương trình Clapeyron là $D = \dfrac{p \cdot M}{R \cdot T}$. (Vì): Từ hệ thức $pV = \dfrac{m}{M} \cdot R \cdot T \Rightarrow D = \dfrac{m}{V} = \dfrac{p \cdot M}{R \cdot T}$. Mệnh đề đúng
\item (Đúng) Khối lượng riêng của khí nitrogen trong bình ở trạng thái trên bằng $2,8\ \mathrm{kg/m^3}$. (Vì): Áp dụng công thức khối lượng riêng trên với $T = 300\ \mathrm{K}$:
\item (Đúng) Nếu nhiệt độ trong bình tăng lên đến $127\ ^\circ\mathrm{C}$ trong khi thể tích bình giữ nguyên không đổi thì khối lượng riêng của khí nitrogen trong bình giảm đi $25\%$. (Vì): o bình kín có thể tích không đổi và không bị hở khí, khối lượng khí $m$ và thể tích $V$ đều cố định nên khối lượng riêng $D = \dfrac{m}{V}$ không đổi. Mệnh đề sai
\end{itemchoice}
}
\end{ex}

% ===== Câu 167 =====
% ID: L12C2B11-44-TLN
% Mức: VD
\begin{ex}
% ID: L12C2B11-44-TLN
% Mức: VD
Một quả bóng thám không có thể tích dãn nở tối đa là $15\ \mathrm{m^3}$ trước khi bị vỡ. Bóng được bơm khí ở mặt đất có thể tích ban đầu là $4\ \mathrm{m^3}$ ở áp suất $1,0 \cdot 10^5\ \mathrm{Pa}$ và nhiệt độ $300\ \mathrm{K}$. Khi bóng bay lên độ cao có nhiệt độ tuyệt đối giảm còn $225\ \mathrm{K}$, áp suất khí quyển tối thiểu mà quả bóng có thể chịu được ngay trước khi vỡ bằng bao nhiêu $\cdot 10^4\ \mathrm{Pa}$?
\shortans{$2$}
\loigiai{
Áp dụng công thức
  \begin{aligned}
 p_2 &= p_1 $\cdot\dfrac{V_1}{V_2}\cdot\dfrac{T_2}{T_1}&= 1,0\cdot10^5\cdot\dfrac{4}{15}\cdot\dfrac{225}{300}&= 2\cdot10^4\\mathrm{Pa}$.
 \end{aligned}
}
\end{ex}

% ===== Câu 168 =====
% ID: L12C2B11-44-TN
% Mức: VD
\begin{ex}
% ID: L12C2B11-44-TN
% Mức: VD
Một bình chứa khí nitrogen ở nhiệt độ $27\ ^\circ\mathrm{C}$ và áp suất $2,0 \cdot 10^5\ \mathrm{Pa}$ có khối lượng riêng là $1,60\ \mathrm{kg/m^3}$. Khi nhiệt độ trong bình tăng lên đến $127\ ^\circ\mathrm{C}$ và áp suất đo được là $3,0 \cdot 10^5\ \mathrm{Pa}$ thì khối lượng riêng của khí nitrogen lúc này bằng bao nhiêu?
\choice
{$1,20\ \mathrm{kg/m^3}$}
{$1,50\ \mathrm{kg/m^3}$}
{\True $1,80\ \mathrm{kg/m^3}$}
{$2,40\ \mathrm{kg/m^3}$}
\loigiai{
Áp dụng công thức
  \begin{aligned}
 D_2 &= D_1 $\cdot\dfrac{p_2}{p_1}\cdot\dfrac{T_1}{T_2}&= 1,60\cdot\dfrac{3,0 \cdot 10^5}{2,0 \cdot 10^5}\cdot\dfrac{27 + 273}{127 + 273}&= 1,80\\mathrm{kg/m^3}$.
 \end{aligned} 
 Chọn đáp án tương ứng.
}
\end{ex}

% ===== Câu 169 =====
% ID: L12C2B11-45-DS
% Mức: VDC
\begin{ex}
% ID: L12C2B11-45-DS
% Mức: VDC
Một bình kín dung tích $10\ \mathrm{l}$ chứa khí hydrogen ở nhiệt độ $27\ ^\circ\mathrm{C}$ và áp suất $2,0 \cdot 10^5\ \mathrm{Pa}$. Người ta nung nóng bình để áp suất khí tăng lên đến $3,0 \cdot 10^5\ \mathrm{Pa}$.
\choiceTF
{\True Quá trình biến đổi trạng thái của chất khí trong bình kín có thể tích không đổi là quá trình đẳng tích.}
{Khi nung nóng bình kín, cả ba thông số trạng thái của khối khí đều thay đổi.}
{\True Nhiệt độ tuyệt đối của khối khí trong bình sau khi nung nóng bằng $450\ \mathrm{K}$.}
{\True Nếu sau đó ta mở van xả bớt một lượng khí ra ngoài để áp suất giảm về $2,0 \cdot 10^5\ \mathrm{Pa}$ trong khi nhiệt độ vẫn giữ ở $450\ \mathrm{K}$, thì tỉ số giữa khối lượng khí còn lại trong bình và khối lượng khí ban đầu bằng $\dfrac{2}{3}$.}
\loigiai{
\begin{itemchoice}
\item (Đúng) Quá trình biến đổi trạng thái của chất khí trong bình kín có thể tích không đổi là quá trình đẳng tích. (Vì): o bình kín có dung tích cố định nên thể tích khí bên trong không đổi, quá trình biến đổi là quá trình đẳng tích. Mệnh đề đúng
\item (Đúng) Khi nung nóng bình kín, cả ba thông số trạng thái của khối khí đều thay đổi. (Vì): o bình kín giữ nguyên thể tích $V = 10\ \mathrm{l}$ nên chỉ có áp suất và nhiệt độ thay đổi, thể tích được giữ cố định. Mệnh đề sai
\item (Đúng) Nhiệt độ tuyệt đối của khối khí trong bình sau khi nung nóng bằng $450\ \mathrm{K}$. (Vì): Áp dụng định luật đẳng tích liên hệ giữa áp suất và nhiệt độ:
\item (Đúng) Nếu sau đó ta mở van xả bớt một lượng khí ra ngoài để áp suất giảm về $2,0 \cdot 10^5\ \mathrm{Pa}$ trong khi nhiệt độ vẫn giữ ở $450\ \mathrm{K}$, thì tỉ số giữa khối lượng khí còn lại trong bình và khối lượng khí ban đầu bằng $\dfrac{2}{3}$. (Vì): Khi giữ ở nhiệt độ $T_{\text{sau}} = 450\ \mathrm{K}$ trong bình dung tích không đổi, khối lượng khí tỉ lệ thuận với áp suất:
\end{itemchoice}
}
\end{ex}

% ===== Câu 170 =====
% ID: L12C2B11-45-TLN
% Mức: VD
\begin{ex}
% ID: L12C2B11-45-TLN
% Mức: VD
\nguon{Dựa trên ví dụ Sách giáo khoa Vật lí 12 Kết nối tri thức với cuộc sống}
	Một cái bơm chứa $100\ \mathrm{cm^3}$ không khí ở nhiệt độ $27\ ^\circ\mathrm{C}$ và áp suất $1,0 \cdot 10^5\ \mathrm{Pa}$. Tính áp suất của không khí trong bơm khi nó bị nén xuống còn $20\ \mathrm{cm^3}$ và tăng nhiệt độ lên đến $87\ ^\circ\mathrm{C}$ theo đơn vị $10^5\ \mathrm{Pa}$.
\shortans{$6$}
\loigiai{
Áp dụng công thức
  \begin{aligned}
 p_2 &= p_1 $\cdot\dfrac{V_1}{V_2}\cdot\dfrac{T_2}{T_1}&= 1,0\cdot10^5\cdot\dfrac{100}{20}\cdot\dfrac{87 + 273}{27 + 273}&= 6\cdot10^5\\mathrm{Pa}$.
 \end{aligned}
}
\end{ex}

% ===== Câu 171 =====
% ID: L12C2B11-45-TN
% Mức: VD
\begin{ex}
% ID: L12C2B11-45-TN
% Mức: VD
Một lốp xe ô tô chứa không khí ở nhiệt độ $17\ ^\circ\mathrm{C}$ dưới áp suất $2,00\ \mathrm{bar}$. Khi xe chạy nhanh, nhiệt độ không khí bên trong lốp tăng lên đến $46\ ^\circ\mathrm{C}$ đồng thời thể tích của lốp dãn nở tăng thêm $1,00\%$. Áp suất của khí bên trong lốp xe lúc này bằng bao nhiêu?
\choice
{$2,08\ \mathrm{bar}$}
{\True $2,18\ \mathrm{bar}$}
{$2,22\ \mathrm{bar}$}
{$2,35\ \mathrm{bar}$}
\loigiai{
Áp dụng công thức
  \begin{aligned}
 p_2 &= p_1 $\cdot\dfrac{V_1}{V_2}\cdot\dfrac{T_2}{T_1}&= 2,00\cdot\dfrac{V_1}{1,01 \cdot V_1}\cdot\dfrac{46 + 273}{17 + 273}&= 2,18\\mathrm{bar}$.
 \end{aligned} 
 Chọn đáp án tương ứng.
}
\end{ex}

% ===== Câu 172 =====
% ID: L12C2B11-46-DS
% Mức: VDC
\begin{ex}
% ID: L12C2B11-46-DS
% Mức: VDC
Một bóng thám không hình cầu dung tích thiết kế tối đa trước khi vỡ là $1200\ \mathrm{l}$ chứa khí helium. Tại mặt đất, bóng được bơm thể tích $400\ \mathrm{l}$ ở áp suất $1,02 \cdot 10^5\ \mathrm{Pa}$ và nhiệt độ $27\ ^\circ\mathrm{C}$.
\choiceTF
{\True Helium là khí hiếm đơn nguyên tử nhẹ hơn không khí nên được dùng để bơm vào bóng thám không.}
{\True Càng lên cao áp suất khí quyển giảm mạnh làm khí trong bóng dãn nở đẩy thành bóng rộng ra cho đến khi đạt thể tích tối đa thì vỡ.}
{\True Nếu bóng thám không bay lên tầng khí quyển có nhiệt độ $-23\ ^\circ\mathrm{C}$ thì áp suất khí quyển tối thiểu tại độ cao mà bóng bắt đầu bị vỡ bằng $0,28 \cdot 10^5\ \mathrm{Pa}$.}
{Biết khối lượng khí helium trong bóng là $64\ \mathrm{g}$ và khối lượng vỏ bóng cùng thiết bị là $1,2\ \mathrm{kg}$. Lực đẩy Archimede của không khí tại mặt đất tác dụng lên quả bóng thám không bằng $52\ \mathrm{N}$ (biết khối lượng riêng không khí mặt đất là $1,29\ \mathrm{kg/m^3}$ và $g = 10\ \mathrm{m/s^2}$).}
\loigiai{
\begin{itemchoice}
\item (Đúng) Helium là khí hiếm đơn nguyên tử nhẹ hơn không khí nên được dùng để bơm vào bóng thám không. (Vì): Helium có khối lượng riêng rất nhỏ so với không khí thông thường, giúp sinh ra lực nổi lớn kéo bóng thám không bay lên cao. Mệnh đề đúng
\item (Sai) Càng lên cao áp suất khí quyển giảm mạnh làm khí trong bóng dãn nở đẩy thành bóng rộng ra cho đến khi đạt thể tích tối đa thì vỡ. (Vì): ự suy giảm mật độ không khí theo độ cao kéo theo áp suất khí quyển giảm dần tạo ra chênh lệch áp suất lớn làm căng phồng bóng. Mệnh đề đúng
\item (Đúng) Nếu bóng thám không bay lên tầng khí quyển có nhiệt độ $-23\ ^\circ\mathrm{C}$ thì áp suất khí quyển tối thiểu tại độ cao mà bóng bắt đầu bị vỡ bằng $0,28 \cdot 10^5\ \mathrm{Pa}$. (Vì): Áp dụng phương trình trạng thái với $V_1 = 400\ \mathrm{l}, T_1 = 300\ \mathrm{K}, V_2 = 1200\ \mathrm{l}, T_2 = 250\ \mathrm{K}$:
\item (Sai) Biết khối lượng khí helium trong bóng là $64\ \mathrm{g}$ và khối lượng vỏ bóng cùng thiết bị là $1,2\ \mathrm{kg}$. Lực đẩy Archimede của không khí tại mặt đất tác dụng lên quả bóng thám không bằng $52\ \mathrm{N}$ (biết khối lượng riêng không khí mặt đất là $1,29\ \mathrm{kg/m^3}$ và $g = 10\ \mathrm{m/s^2}$). (Vì): Thể tích khí bơm tại mặt đất là $V_1 = 400\ \mathrm{l} = 0,4\ \mathrm{m^3}$. Lực đẩy Archimede tại mặt đất là:
\end{itemchoice}
}
\end{ex}

% ===== Câu 173 =====
% ID: L12C2B11-46-TLN
% Mức: VD
\begin{ex}
% ID: L12C2B11-46-TLN
% Mức: VD
\nguon{Dựa trên bài tập ví dụ mục II.2 trang 46 Sách giáo khoa Vật lí 12 Kết nối tri thức với cuộc sống}
	Một lốp xe máy có dung tích $20\ \mathrm{l}$ chứa không khí ở nhiệt độ $27\ ^\circ\mathrm{C}$ và áp suất $3,0 \cdot 10^5\ \mathrm{Pa}$. Biết khối lượng mol của không khí là $28,8\ \mathrm{g/mol}$ và hằng số khí lí tưởng $R = 8,31\ \mathrm{J/(mol\cdot K)}$. Tính khối lượng của lượng không khí chứa trong lốp xe theo đơn vị $\mathrm{g}$ (làm tròn kết quả đến phần nguyên).
\shortans{$69$}
\loigiai{
Áp dụng công thức
  \begin{aligned}
 m &= $\dfrac{p \cdot V \cdot M}{R \cdot T}&=\dfrac{3,0 \cdot 10^5 \cdot 20 \cdot 10^{-3} \cdot 28,8 \cdot 10^{-3}}{8,31 \cdot (27 + 273)}$ &\approx 0,0693 $\\mathrm{kg}$ = 69,3 $\\mathrm{g}$.
 \end{aligned}
}
\end{ex}

% ===== Câu 174 =====
% ID: L12C2B11-46-TN
% Mức: VD
\begin{ex}
% ID: L12C2B11-46-TN
% Mức: VD
Một bình chứa khí helium ở áp suất $2,0 \cdot 10^5\ \mathrm{Pa}$, thể tích $15\ \mathrm{l}$ và nhiệt độ $27\ ^\circ\mathrm{C}$. Số mol khí helium chứa trong bình lúc này bằng bao nhiêu?
\choice
{$0,85\ \mathrm{mol}$}
{\True $1,20\ \mathrm{mol}$}
{$1,44\ \mathrm{mol}$}
{$1,80\ \mathrm{mol}$}
\loigiai{
Áp dụng công thức $n = \dfrac{p \cdot V}{R \cdot T} = \dfrac{2,0 \cdot 10^5 \cdot 15 \cdot 10^{-3}}{8,31 \cdot (27 + 273)} = 1,20\ \mathrm{mol}$.

Chọn đáp án tương ứng.
}
\end{ex}

% ===== Câu 175 =====
% ID: L12C2B11-47-DS
% Mức: VDC
\begin{ex}
% ID: L12C2B11-47-DS
% Mức: VDC
Một lượng khí lí tưởng xác định thực hiện biến đổi trạng thái từ trạng thái $1$ sang trạng thái $2$, rồi sang trạng thái $3$. Trạng thái $1$ có $p_1 = 1,2 \cdot 10^5\ \mathrm{Pa}, V_1 = 4,0\ \mathrm{l}, T_1 = 300\ \mathrm{K}$. Quá trình $1 \to 2$ là đẳng áp đưa thể tích đến $V_2 = 6,0\ \mathrm{l}$. Quá trình $2 \to 3$ là đẳng tích đưa áp suất đến $p_3 = 2,0 \cdot 10^5\ \mathrm{Pa}$.
\choiceTF
{\True Trong quá trình đẳng tích, áp suất của một lượng khí xác định tỉ lệ thuận với nhiệt độ tuyệt đối của nó.}
{\True Nhiệt độ tuyệt đối $T_2$ ở trạng thái $2$ bằng $450\ \mathrm{K}$.}
{Nhiệt độ tuyệt đối $T_3$ ở trạng thái $3$ bằng $600\ \mathrm{K}$.}
{\True Công thức liên hệ trực tiếp giữa trạng thái $1$ và trạng thái $3$ là $\dfrac{p_1 \cdot V_1}{T_1} = \dfrac{p_3 \cdot V_2}{T_3}$.}
\loigiai{
\begin{itemchoice}
\item (Đúng) Trong quá trình đẳng tích, áp suất của một lượng khí xác định tỉ lệ thuận với nhiệt độ tuyệt đối của nó. (Vì): ây là phát biểu định luật Gay-Lussac cho quá trình đẳng tích của lượng khí xác định. Mệnh đề đúng
\item (Đúng) Nhiệt độ tuyệt đối $T_2$ ở trạng thái $2$ bằng $450\ \mathrm{K}$. (Vì): Quá trình $1 \to 2$ đẳng áp tuân theo định luật Charles:
\item (Sai) Nhiệt độ tuyệt đối $T_3$ ở trạng thái $3$ bằng $600\ \mathrm{K}$. (Vì): Quá trình $2 \to 3$ đẳng tích tuân theo định luật Gay-Lussac:
\item (Đúng) Công thức liên hệ trực tiếp giữa trạng thái $1$ và trạng thái $3$ là $\dfrac{p_1 \cdot V_1}{T_1} = \dfrac{p_3 \cdot V_2}{T_3}$. (Vì): o quá trình $2 \to 3$ đẳng tích nên $V_3 = V_2$. Thay vào phương trình trạng thái tổng quát $\dfrac{p_1 \cdot V_1}{T_1} = \dfrac{p_3 \cdot V_3}{T_3}$ thu được hệ thức liên hệ trực tiếp. Mệnh đề đúng
\end{itemchoice}
}
\end{ex}

% ===== Câu 176 =====
% ID: L12C2B11-47-TLN
% Mức: VD
\begin{ex}
% ID: L12C2B11-47-TLN
% Mức: VD
Phản ứng phân hủy chất rắn $\mathrm{NaN_3}$ trong túi khí ô tô giải phóng ra khí nitrogen và kim loại sodium ($\mathrm{Na}$) bám vào túi khí. Nếu phản ứng giải phóng ra $3,0\ \mathrm{mol}$ khí nitrogen thì khối lượng sodium tạo thành đồng thời bằng bao nhiêu gam? Biết khối lượng mol của sodium là $23\ \mathrm{g/mol}$.
\shortans{$46$}
\loigiai{
Xác định số mol kim loại sodium sinh ra:
  \begin{aligned}
 n_{ $\mathrm{Na}} &=\dfrac{2}{3}\cdotn_{\mathrm{N_2}} 
&=\dfrac{2}{3}\cdot$ 3,0 
&= 2,0 $\\mathrm{mol}$.
 \end{aligned} 
 Khối lượng sodium tạo thành là:
  \begin{aligned}
 m &= n_{ $\mathrm{Na}$ } \cdot M_{ $\mathrm{Na}$ } 
&= 2,0 \cdot 23 
&= 46 $\\mathrm{g}$.
 \end{aligned}
}
\end{ex}

% ===== Câu 177 =====
% ID: L12C2B11-47-TN
% Mức: VD
\begin{ex}
% ID: L12C2B11-47-TN
% Mức: VD
Khi xảy ra va chạm xe ô tô, phản ứng phân hủy $100\ \mathrm{g}$ chất rắn $\mathrm{NaN_3}$ trong thiết bị an toàn đã giải phóng ra $2,31\ \mathrm{mol}$ khí nitrogen ($\mathrm{N_2}$). Biết túi khí khi phồng lên tối đa có dung tích bằng $48\ \mathrm{l}$ ở nhiệt độ $30\ ^\circ\mathrm{C}$. Bỏ qua thể tích chất rắn Na được tạo thành, áp suất của khí $\mathrm{N_2}$ trong túi lúc này bằng bao nhiêu?
\choice
{$0,81 \cdot 10^5\ \mathrm{Pa}$}
{$1,02 \cdot 10^5\ \mathrm{Pa}$}
{\True $1,21 \cdot 10^5\ \mathrm{Pa}$}
{$1,48 \cdot 10^5\ \mathrm{Pa}$}
\loigiai{
Áp dụng công thức $p = \dfrac{n \cdot R \cdot T}{V} = \dfrac{2,31 \cdot 8,31 \cdot (30 + 273)}{48 \cdot 10^{-3}} \approx 1,21 \cdot 10^5 \ \mathrm{Pa}$.

Chọn đáp án tương ứng.
}
\end{ex}

% ===== Câu 178 =====
% ID: L12C2B11-48-DS
% Mức: VDC
\begin{ex}
% ID: L12C2B11-48-DS
% Mức: VDC
Ở điều kiện tiêu chuẩn (áp suất $p_0 = 760\ \mathrm{mmHg}$ và nhiệt độ $T_0 = 273\ \mathrm{K}$), khối lượng riêng của không khí ở sát mặt đất là $1,29\ \mathrm{kg/m^3}$. Người ta đo không khí trên đỉnh một ngọn núi có áp suất $446\ \mathrm{mmHg}$ và nhiệt độ $2\ ^\circ\mathrm{C}$.
\choiceTF
{\True Khối lượng riêng của một chất khí là khối lượng của khí chứa trong một đơn vị thể tích.}
{Trong quá trình đẳng áp, khối lượng riêng của khối khí tỉ lệ thuận với nhiệt độ tuyệt đối của nó.}
{\True Trên đỉnh ngọn núi đó, khối lượng riêng của không khí xấp xỉ bằng $0,75\ \mathrm{kg/m^3}$.}
{\True Nếu nhiệt độ không khí trên đỉnh núi giảm tiếp về $-10\ ^\circ\mathrm{C}$ trong khi áp suất không thay đổi thì khối lượng riêng của không khí tăng thêm khoảng $0,03\ \mathrm{kg/m^3}$.}
\loigiai{
\begin{itemchoice}
\item (Đúng) Khối lượng riêng của một chất khí là khối lượng của khí chứa trong một đơn vị thể tích. (Vì): Theo định nghĩa vật lí học, khối lượng riêng $D = \dfrac{m}{V}$ biểu thị khối lượng chất khí phân bố đều trên một đơn vị thể tích. Mệnh đề đúng
\item (Sai) Trong quá trình đẳng áp, khối lượng riêng của khối khí tỉ lệ thuận với nhiệt độ tuyệt đối của nó. (Vì): Trong quá trình đẳng áp, thể tích khí tỉ lệ thuận với nhiệt độ tuyệt đối nên khối lượng riêng tỉ lệ nghịch với nhiệt độ tuyệt đối. Mệnh đề sai
\item (Đúng) Trên đỉnh ngọn núi đó, khối lượng riêng của không khí xấp xỉ bằng $0,75\ \mathrm{kg/m^3}$. (Vì): Khối lượng riêng không khí trên đỉnh núi Fansipan với $T = 275\ \mathrm{K}$ được tính theo hệ thức biến đổi:
\item (Đúng) Nếu nhiệt độ không khí trên đỉnh núi giảm tiếp về $-10\ ^\circ\mathrm{C}$ trong khi áp suất không thay đổi thì khối lượng riêng của không khí tăng thêm khoảng $0,03\ \mathrm{kg/m^3}$. (Vì): Từ kết quả trên, tại nhiệt độ $T' = 263\ \mathrm{K}$, khối lượng riêng không khí lúc sau là:
\end{itemchoice}
}
\end{ex}

% ===== Câu 179 =====
% ID: L12C2B11-48-TLN
% Mức: VD
\begin{ex}
% ID: L12C2B11-48-TLN
% Mức: VD
\nguon{Bài tập vận dụng 2 trang 44 Sách giáo khoa Vật lí 12 Kết nối tri thức với cuộc sống}
	Một khối khí hydrogen có khối lượng $12\ \mathrm{g}$ chiếm thể tích ban đầu $4,0\ \mathrm{l}$ ở nhiệt độ $7\ ^\circ\mathrm{C}$. Người ta đun nóng đẳng áp khối khí này cho đến khi khối lượng riêng của nó giảm còn $1,2\ \mathrm{g/l}$. Tính nhiệt độ của khối khí hydrogen ở trạng thái sau theo đơn vị độ Celsius ($^\circ\mathrm{C}$).
\shortans{$427$}
\loigiai{
Xác định khối lượng riêng ban đầu:
  \begin{aligned}
 D_1 &= $\dfrac{m}{V_1}&=\dfrac{12}{4,0}&= 3\\mathrm{g/l}$.
 \end{aligned} 
 Nhiệt độ tuyệt đối ở trạng thái sau trong quá trình đẳng áp:
  \begin{aligned}
 T_2 &= T_1 $\cdot\dfrac{D_1}{D_2}$ &= (7 + 273) $\cdot\dfrac{3}{1,2}$ &= 700 $\\mathrm{K}$.
 \end{aligned} 
 Đổi sang nhiệt độ Celsius:
  \begin{aligned}
 t_2 &= T_2 - 273 
&= 700 - 273 
&= 427 $\$ ^ $\circ\mathrm{C}$.
 \end{aligned}
}
\end{ex}

% ===== Câu 180 =====
% ID: L12C2B11-48-TN
% Mức: VD
\begin{ex}
% ID: L12C2B11-48-TN
% Mức: VD
Một bình dung tích $24,93\ \mathrm{l}$ chứa khí nitrogen ($\mathrm{N_2}$) ở nhiệt độ $27\ ^\circ\mathrm{C}$ dưới áp suất $2,5 \cdot 10^5\ \mathrm{Pa}$. Người ta mở van xả bớt một phần khí ra ngoài làm áp suất khí trong bình giảm còn $1,5 \cdot 10^5\ \mathrm{Pa}$ trong khi nhiệt độ và thể tích bình vẫn giữ không đổi. Biết khối lượng mol của khí nitrogen là $28\ \mathrm{g/mol}$. Khối lượng khí nitrogen đã thoát ra ngoài bằng bao nhiêu?
\choice
{$14\ \mathrm{g}$}
{\True $28\ \mathrm{g}$}
{$42\ \mathrm{g}$}
{$56\ \mathrm{g}$}
\loigiai{
Áp dụng công thức
\begin{aligned}
$\Deltam&=\dfrac{\Delta p \cdot V \cdot M}{R \cdot T}=\dfrac{(2,5 \cdot 10^5 - 1,5 \cdot 10^5) \cdot 24,93 \cdot 10^{-3} \cdot 28 \cdot 10^{-3}}{8,31 \cdot (27 + 273)}= 28\cdot10^{-3}$
\mathrm{kg}= 28
\mathrm{g}.
\end{aligned}
Chọn đáp án tương ứng.
}
\end{ex}

% ===== Câu 181 =====
% ID: L12C2B11-49-DS
% Mức: VDC
\begin{ex}
% ID: L12C2B11-49-DS
% Mức: VDC
Một túi khí có dung tích cố định $41,55\ \mathrm{l}$ chứa khí nitrogen sinh ra từ phản ứng phân hủy sodium azide. Tại thời điểm sau phản ứng, nhiệt độ khí đạt $47\ ^\circ\mathrm{C}$ và áp suất trong túi khí đo được là $1,28 \cdot 10^5\ \mathrm{Pa}$.
\choiceTF
{\True Theo thuyết động học phân tử, nhiệt độ của khối khí càng cao thì tốc độ chuyển động hỗn loạn của các phân tử $\mathrm{N_2}$ càng lớn.}
{Số mol khí nitrogen đã được giải phóng chứa trong túi khí bằng $4,0\ \mathrm{mol}$.}
{\True Số mol của chất rắn $\mathrm{NaN_3}$ đã tham gia phản ứng phân hủy xấp xỉ bằng $1,33\ \mathrm{mol}$.}
{\True Khối lượng kim loại Na tạo thành đồng thời bám vào túi khí bằng $30,6\ \mathrm{g}$.}
\loigiai{
\begin{itemchoice}
\item (Đúng) Theo thuyết động học phân tử, nhiệt độ của khối khí càng cao thì tốc độ chuyển động hỗn loạn của các phân tử $\mathrm{N_2}$ càng lớn. (Vì): Thuyết động học phân tử khẳng định chuyển động hỗn loạn của các hạt tỉ lệ thuận với nhiệt độ tuyệt đối của môi trường. Mệnh đề đúng
\item (Sai) Số mol khí nitrogen đã được giải phóng chứa trong túi khí bằng $4,0\ \mathrm{mol}$. (Vì): Áp dụng phương trình Clapeyron để tìm số mol khí nitrogen với $T = 320\ \mathrm{K}$:
\item (Đúng) Số mol của chất rắn $\mathrm{NaN_3}$ đã tham gia phản ứng phân hủy xấp xỉ bằng $1,33\ \mathrm{mol}$. (Vì): Từ kết quả số mol khí trên, số mol của sodium azide đã phân rã hoàn toàn là:
\item (Đúng) Khối lượng kim loại Na tạo thành đồng thời bám vào túi khí bằng $30,6\ \mathrm{g}$. (Vì): Theo phương trình phản ứng, số mol Na sinh ra bằng số mol $\mathrm{NaN_3}$ phản ứng. Khối lượng Na tạo thành là:
\end{itemchoice}
}
\end{ex}

% ===== Câu 182 =====
% ID: L12C2B11-49-TLN
% Mức: VD
\begin{ex}
% ID: L12C2B11-49-TLN
% Mức: VD
Một khối khí lí tưởng xác định ở nhiệt độ $27\ ^\circ\mathrm{C}$ có thể tích $6\ \mathrm{l}$ dưới áp suất $1,0\ \mathrm{atm}$. Người ta nén khối khí này đến thể tích $2\ \mathrm{l}$ đồng thời nung nóng đến nhiệt độ $127\ ^\circ\mathrm{C}$. Tính áp suất của khí lúc này theo đơn vị $\mathrm{atm}$.
\shortans{$4$}
\loigiai{
Áp dụng công thức
  \begin{aligned}
 p_2 &= p_1 $\cdot\dfrac{V_1}{V_2}\cdot\dfrac{T_2}{T_1}&= 1,0\cdot\dfrac{6}{2}\cdot\dfrac{127 + 273}{27 + 273}&= 4\\mathrm{atm}$.
 \end{aligned}
}
\end{ex}

% ===== Câu 183 =====
% ID: L12C2B11-49-TN
% Mức: VDC
\begin{ex}
% ID: L12C2B11-49-TN
% Mức: VDC
Xác định khối lượng riêng của không khí trên đỉnh núi Fansipan cao $3140\ \mathrm{m}$, biết rằng cứ lên cao thêm $10\ \mathrm{m}$ thì áp suất khí quyển giảm đi $1\ \mathrm{mmHg}$ và nhiệt độ trên đỉnh núi lúc đó là $2\ ^\circ\mathrm{C}$. Biết điều kiện chuẩn ở mặt đất có áp suất $760\ \mathrm{mmHg}$, nhiệt độ $0\ ^\circ\mathrm{C}$ và khối lượng riêng của không khí là $1,29\ \mathrm{kg/m^3}$.
\choice
{$0,54\ \mathrm{kg/m^3}$}
{$0,68\ \mathrm{kg/m^3}$}
{\True $0,75\ \mathrm{kg/m^3}$}
{$0,82\ \mathrm{kg/m^3}$}
\loigiai{
Áp suất trên đỉnh Fansipan là:
$p = p_0 - \dfrac{h}{10} = 760 - \dfrac{3140}{10} = 446\ \mathrm{mmHg}$.
Áp dụng công thức tỉ số khối lượng riêng từ phương trình trạng thái:
$D = D_0 \cdot \dfrac{p}{p_0} \cdot \dfrac{T_0}{T} = 1,29 \cdot \dfrac{446}{760} \cdot \dfrac{273}{2 + 273} \approx 0,75\ \mathrm{kg/m^3}$.
Chọn đáp án tương ứng.
}
\end{ex}

% ===== Câu 184 =====
% ID: L12C2B11-50-DS
% Mức: VDC
\begin{ex}
% ID: L12C2B11-50-DS
% Mức: VDC
Một khối lượng khí lí tưởng xác định thực hiện chu trình biến đổi trạng thái khép kín được biểu diễn trong hệ tọa độ $(p, T)$ gồm bốn quá trình liên tiếp: $(1) \to (2)$ đẳng nhiệt; $(2) \to (3)$ đẳng áp; $(3) \to (4)$ đẳng nhiệt; $(4) \to (1)$ đẳng áp.
\choiceTF
{\True Trong quá trình $(1) \to (2)$, thể tích của khối khí tỉ lệ nghịch với áp suất.}
{\True Đường biểu diễn quá trình $(2) \to (3)$ trong hệ tọa độ $(p, T)$ là một đoạn thẳng song song với trục nhiệt độ $OT$.}
{\True Nếu ở trạng thái $1$ có $p_1 = 1,0 \cdot 10^5\ \mathrm{Pa}$, $V_1 = 4,0\ \mathrm{l}$ và trạng thái $2$ có $p_2 = 2,0 \cdot 10^5\ \mathrm{Pa}$, thì thể tích của khối khí ở trạng thái $2$ bằng $2,0\ \mathrm{l}$.}
{\True Biết nhiệt độ $T_1 = 300\ \mathrm{K}$, nhiệt độ $T_3 = 600\ \mathrm{K}$. Nếu thể tích ở trạng thái $3$ bằng thể tích ở trạng thái $1$ ($V_3 = V_1 = 4,0\ \mathrm{l}$), thì áp suất $p_3$ ở trạng thái $3$ bằng $2,0 \cdot 10^5\ \mathrm{Pa}$.}
\loigiai{
\begin{itemchoice}
\item (Đúng) Trong quá trình $(1) \to (2)$, thể tích của khối khí tỉ lệ nghịch với áp suất. (Vì): Quá trình $(1) \to (2)$ là đẳng nhiệt nên áp suất và thể tích tuân theo định luật Boyle, tỉ lệ nghịch với nhau. Mệnh đề đúng
\item (Đúng) Đường biểu diễn quá trình $(2) \to (3)$ trong hệ tọa độ $(p, T)$ là một đoạn thẳng song song với trục nhiệt độ $OT$. (Vì): Quá trình $(2) \to (3)$ là đẳng áp nên áp suất $p$ không thay đổi. Trong hệ tọa độ $(p, T)$, đường biểu diễn là một đoạn thẳng song song với trục hoành $OT$. Mệnh đề đúng
\item (Đúng) Nếu ở trạng thái $1$ có $p_1 = 1,0 \cdot 10^5\ \mathrm{Pa}$, $V_1 = 4,0\ \mathrm{l}$ và trạng thái $2$ có $p_2 = 2,0 \cdot 10^5\ \mathrm{Pa}$, thì thể tích của khối khí ở trạng thái $2$ bằng $2,0\ \mathrm{l}$. (Vì): Áp dụng định luật Boyle cho quá trình đẳng nhiệt $(1) \to (2)$:
\item (Đúng) Biết nhiệt độ $T_1 = 300\ \mathrm{K}$, nhiệt độ $T_3 = 600\ \mathrm{K}$. Nếu thể tích ở trạng thái $3$ bằng thể tích ở trạng thái $1$ ($V_3 = V_1 = 4,0\ \mathrm{l}$), thì áp suất $p_3$ ở trạng thái $3$ bằng $2,0 \cdot 10^5\ \mathrm{Pa}$. (Vì): Quá trình $(2) \to (3)$ là đẳng áp nên $p_3 = p_2 = 2,0 \cdot 10^5\ \mathrm{Pa}$. Kiểm tra với thể tích dãn nở:
\end{itemchoice}
}
\end{ex}

% ===== Câu 185 =====
% ID: L12C2B11-50-TLN
% Mức: TH
\dangbt{Áp dụng phương trình Clapeyron pV = nRT}
\begin{ex}
% ID: L12C2B11-50-TLN
% Mức: TH
Trong dạng “Áp dụng phương trình Clapeyron pV = nRT”, Có 1 mol khí ở 280 $K$ chiếm 6 lít. Tính áp suất theo kPa, lấy R=8,31.
\shortans{387,8}
\loigiai{
$p=nRT/V=387,8$ kPa.
}
\end{ex}

% ===== Câu 186 =====
% ID: L12C2B11-50-TN
% Mức: TH
\dangbt{Xác định sự biến thiên khối lượng riêng của chất khí}
\begin{ex}
% ID: L12C2B11-50-TN
% Mức: TH
Giải thích nào sau đây là đúng về việc càng bay lên cao thì thể tích của bóng thám không càng tăng lên cho đến khi bị vỡ?
\choice
{Do nhiệt độ khí quyển ở trên cao tăng lên làm khí bên trong nóng dãn nở}
{\True Do áp suất khí quyển trên cao giảm dần làm khí bên trong dãn nở vỏ bóng}
{Do vỏ bóng bị mài mòn bởi gió trên cao làm tăng diện tích bề mặt}
{Do lực hút Trái Đất giảm làm giảm lực liên kết phân tử của khí bên trong}
\loigiai{
Càng bay lên cao, mật độ không khí càng giảm dẫn đến áp suất khí quyển giảm dần. Sự chênh lệch áp suất lớn giữa bên trong và bên ngoài làm khí bên trong vỏ dãn nở dãn vỏ bóng to ra cho đến khi vượt quá giới hạn đàn hồi của cao su làm bóng bị vỡ.
 Chọn đáp án tương ứng.
}
\end{ex}

% ===== Câu 187 =====
% ID: L12C2B11-51-DS
% Mức: TH
\dangbt{Áp dụng phương trình Clapeyron pV = nRT}
\begin{ex}
% ID: L12C2B11-51-DS
% Mức: TH
Xét các phát biểu sau trong dạng “Áp dụng phương trình Clapeyron pV = nRT”.
\choiceTF
{\True Phương trình Clapeyron của khí lí tưởng là — phát biểu đúng là: $pV=nRT$.}
{Phương trình Clapeyron của khí lí tưởng là — phát biểu $pV=RT/n$ là đúng.}
{\True Với lượng khí xác định, phương trình trạng thái giữa hai trạng thái là — phát biểu đúng là: $p_1V_1/T_1=p_2V_2/T_2$.}
{Với lượng khí xác định, phương trình trạng thái giữa hai trạng thái là — phát biểu $p_1/T_1=p_2/T_2$ trong mọi quá trình là đúng.}
\loigiai{
\begin{itemchoice}
\item (Đúng) Phương trình Clapeyron của khí lí tưởng là — phát biểu đúng là: $pV=nRT$. (Vì): Một lượng khí lí tưởng thỏa $pV=nRT$. b) Sai: trái với hệ thức hoặc điều kiện đã nêu. c) Đúng: Do $pV/T=nR$ không đổi. d) Sai: trái với hệ thức hoặc điều kiện đã nêu
\item (Sai) Phương trình Clapeyron của khí lí tưởng là — phát biểu $pV=RT/n$ là đúng.
\item (Đúng) Với lượng khí xác định, phương trình trạng thái giữa hai trạng thái là — phát biểu đúng là: $p_1V_1/T_1=p_2V_2/T_2$.
\item (Sai) Với lượng khí xác định, phương trình trạng thái giữa hai trạng thái là — phát biểu $p_1/T_1=p_2/T_2$ trong mọi quá trình là đúng.
\end{itemchoice}
}
\end{ex}

% ===== Câu 188 =====
% ID: L12C2B11-51-TLN
% Mức: TH
\begin{ex}
% ID: L12C2B11-51-TLN
% Mức: TH
Trong dạng “Áp dụng phương trình Clapeyron pV = nRT”, Khí có $140\,\text{kPa}$, 8 lít ở 290 K. Đun đẳng tích đến 350 K. Tính áp suất sau theo kPa.
\shortans{168,97}
\loigiai{
$p_2=p_1T_2/T_1=168,97$ kPa.
}
\end{ex}

% ===== Câu 189 =====
% ID: L12C2B11-51-TN
% Mức: NB
\dangbt{Xác định sự biến thiên khối lượng riêng của chất khí}
\begin{ex}
% ID: L12C2B11-51-TN
% Mức: NB
Trong hệ thống an toàn túi khí của ô tô, khi xảy ra va chạm mạnh, hệ thống cảm biến sẽ kích thích một phản ứng phân hủy chất rắn để sinh ra một lượng lớn khí làm phồng túi nhanh chóng. Chất rắn và lượng khí sinh ra tương ứng là
\choice
{$\mathrm{NaN_3}$ phân hủy tạo khí $\mathrm{O_2}$}
{$\mathrm{Na_2CO_3}$ phân hủy tạo khí $\mathrm{CO_2}$}
{\True $\mathrm{NaN_3}$ phân hủy tạo khí $\mathrm{N_2}$}
{$\mathrm{KClO_3}$ phân hủy tạo khí $\mathrm{Cl_2}$}
\loigiai{
Trong thiết bị túi khí ô tô, chất rắn $\mathrm{NaN_3}$ (sodium azide) khi bị kích thích bởi cảm biến va chạm sẽ xảy ra phản ứng phân hủy nhanh chóng giải phóng lượng lớn khí nitrogen ($\mathrm{N_2}$) làm phồng căng túi khí để bảo vệ người lái.
 Chọn đáp án tương ứng.
}
\end{ex}

% ===== Câu 190 =====
% ID: L12C2B11-52-DS
% Mức: TH
\dangbt{Áp dụng phương trình Clapeyron pV = nRT}
\begin{ex}
% ID: L12C2B11-52-DS
% Mức: TH
Xét các phát biểu sau trong dạng “Áp dụng phương trình Clapeyron pV = nRT”.
\choiceTF
{\True Có 2 mol khí ở 260 $K$ chiếm 7 lít. Lấy R= $8,31\,\text{J}$ /(mol.K). Áp suất gần bằng — phát biểu đúng là: $617,31\,\text{kPa}$.}
{Có 2 mol khí ở 260 $K$ chiếm 7 lít. Lấy R= $8,31\,\text{J}$ /(mol.K). Áp suất gần bằng — phát biểu $61,73\,\text{kPa}$ là đúng.}
{\True Khí chuyển từ $p_1=90$ kPa, $V_1=2 lít,$ T_1=270 $K đến$ V_2=8 $lít,$ T_2=350 $K.$ p_2 $gần bằng — phát biểu đúng là: 29,17 kPa.$}
{Khí chuyển từ $p_1=90$ kPa, $V_1=2 lít,$ T_1=270 $K đến$ V_2=8 $lít,$ T_2=350 $K.$ p_2 $gần bằng — phát biểu 69,43 kPa là đúng.$}
\loigiai{
\begin{itemchoice}
\item (Đúng) Có 2 mol khí ở 260 $K$ chiếm 7 lít. Lấy R= $8,31\,\text{J}$ /(mol.K). Áp suất gần bằng — phát biểu đúng là: $617,31\,\text{kPa}$. (Vì): $p=nRT/V=617,31$ kPa. b) Sai: trái với hệ thức hoặc điều kiện đã nêu. c) Đúng: $p_2=p_1V_1T_2/(V_2T_1)=29,17$ kPa. d) Sai: trái với hệ thức hoặc điều kiện đã nêu
\item (Sai) Có 2 mol khí ở 260 $K$ chiếm 7 lít. Lấy R= $8,31\,\text{J}$ /(mol.K). Áp suất gần bằng — phát biểu $61,73\,\text{kPa}$ là đúng.
\item (Đúng) Khí chuyển từ $p_1=90$ kPa, $V_1=2 lít,$ T_1=270 $K đến$ V_2=8 $lít,$ T_2=350 $K.$ p_2 $gần bằng — phát biểu đúng là: 29,17 kPa.$.
\item (Sai) Khí chuyển từ $p_1=90$ kPa, $V_1=2 lít,$ T_1=270 $K đến$ V_2=8 $lít,$ T_2=350 $K.$ p_2 $gần bằng — phát biểu 69,43 kPa là đúng.$.
\end{itemchoice}
}
\end{ex}

% ===== Câu 191 =====
% ID: L12C2B11-52-TLN
% Mức: VD
\begin{ex}
% ID: L12C2B11-52-TLN
% Mức: VD
Trong dạng “Áp dụng phương trình Clapeyron pV = nRT”, Khí có $150\,\text{kPa}$, 4 lít ở 300 $K$; chuyển đến $600\,\text{kPa}$ và 360 K. Tính thể tích sau theo lít.
\shortans{1,2}
\loigiai{
$V_2=p_1V_1T_2/(p_2T_1)=1,2$ lít.
}
\end{ex}

% ===== Câu 192 =====
% ID: L12C2B11-52-TN
% Mức: TH
\dangbt{Xác định sự biến thiên khối lượng riêng của chất khí}
\begin{ex}
% ID: L12C2B11-52-TN
% Mức: TH
Đối với một lượng khí lí tưởng xác định thực hiện quá trình biến đổi đẳng áp, mối liên hệ giữa khối lượng riêng $D$ của khối khí và nhiệt độ tuyệt đối $T$ của nó được xác định bởi hệ thức nào sau đây?
\choice
{Thương số $\dfrac{D}{T} = \text{const}$}
{\True Tích số $D \cdot T = \text{const}$}
{Tích số $D \cdot T^2 = \text{const}$}
{Thương số $\dfrac{D^2}{T} = \text{const}$}
\loigiai{
Khối lượng riêng $D = \dfrac{m}{V}$. Với khối lượng khí $m$ không đổi, khối lượng riêng $D$ tỉ lệ nghịch với thể tích $V$. Theo định luật Charles cho quá trình đẳng áp, thể tích $V$ tỉ lệ thuận với nhiệt độ tuyệt đối $T$. Do đó, khối lượng riêng $D$ tỉ lệ nghịch với nhiệt độ tuyệt đối $T$, tức là tích số $D \cdot T = \text{const}$.
 Chọn đáp án tương ứng.
}
\end{ex}

% ===== Câu 193 =====
% ID: L12C2B11-53-DS
% Mức: VD
\dangbt{Áp dụng phương trình Clapeyron pV = nRT}
\begin{ex}
% ID: L12C2B11-53-DS
% Mức: VD
Xét các phát biểu sau trong dạng “Áp dụng phương trình Clapeyron pV = nRT”.
\choiceTF
{\True Trong phương trình khí lí tưởng, nhiệt độ phải dùng — phát biểu đúng là: kelvin.}
{Trong phương trình khí lí tưởng, nhiệt độ phải dùng — phát biểu độ $C$ là đúng.}
{\True Ở thể tích không đổi, áp suất của lượng khí xác định tỉ lệ — phát biểu đúng là: thuận với nhiệt độ tuyệt đối.}
{Ở thể tích không đổi, áp suất của lượng khí xác định tỉ lệ — phát biểu thuận với bình phương nhiệt độ là đúng.}
\loigiai{
\begin{itemchoice}
\item (Đúng) Trong phương trình khí lí tưởng, nhiệt độ phải dùng — phát biểu đúng là: kelvin. (Vì): $T$ phải là nhiệt độ tuyệt đối. b) Sai: trái với hệ thức hoặc điều kiện đã nêu. c) Đúng: Từ $pV=nRT$, khi $V$ cố định thì $p/T$ không đổi. d) Sai: trái với hệ thức hoặc điều kiện đã nêu
\item (Sai) Trong phương trình khí lí tưởng, nhiệt độ phải dùng — phát biểu độ $C$ là đúng.
\item (Đúng) Ở thể tích không đổi, áp suất của lượng khí xác định tỉ lệ — phát biểu đúng là: thuận với nhiệt độ tuyệt đối.
\item (Sai) Ở thể tích không đổi, áp suất của lượng khí xác định tỉ lệ — phát biểu thuận với bình phương nhiệt độ là đúng.
\end{itemchoice}
}
\end{ex}

% ===== Câu 194 =====
% ID: L12C2B11-53-TLN
% Mức: VD
\begin{ex}
% ID: L12C2B11-53-TLN
% Mức: VD
Trong dạng “Áp dụng phương trình Clapeyron pV = nRT”, Một khí có khối lượng $28\,\text{g}$, khối lượng mol $28\,\text{g}$ /mol. Tính số mol.
\shortans{1}
\loigiai{
$n=m/M=28/28=1$ mol.
}
\end{ex}

% ===== Câu 195 =====
% ID: L12C2B11-53-TN
% Mức: VD
\dangbt{Xác định sự biến thiên khối lượng riêng của chất khí}
\begin{ex}
% ID: L12C2B11-53-TN
% Mức: VD
Một khối khí lí tưởng xác định thực hiện quá trình chuyển trạng thái sao cho thể tích của nó tăng thêm $20\%$ và nhiệt độ tuyệt đối tăng thêm $50\%$. Áp suất của khối khí lúc sau bằng bao nhiêu so với áp suất ban đầu?
\choice
{Bằng $0,80\ \text{lần}$}
{Bằng $1,12\ \text{lần}$}
{\True Bằng $1,25\ \text{lần}$}
{Bằng $1,80\ \text{lần}$}
\loigiai{
Áp dụng công thức
  \begin{aligned}
 $\dfrac{p_2}{p_1}&=\dfrac{V_1}{V_2}\cdot\dfrac{T_2}{T_1}&=\dfrac{V_1}{1,2 \cdot V_1}\cdot\dfrac{1,5 \cdot T_1}{T_1}$ 
&= 1,25.
 \end{aligned} 
 Chọn đáp án tương ứng.
}
\end{ex}

% ===== Câu 196 =====
% ID: L12C2B11-54-DS
% Mức: VD
\dangbt{Áp dụng phương trình Clapeyron pV = nRT}
\begin{ex}
% ID: L12C2B11-54-DS
% Mức: VD
Xét các phát biểu sau trong dạng “Áp dụng phương trình Clapeyron pV = nRT”.
\choiceTF
{\True Số mol khí được tính từ khối lượng m và khối lượng mol $M$ bằng — phát biểu đúng là: $n=m/M$.}
{Số mol khí được tính từ khối lượng m và khối lượng mol $M$ bằng — phát biểu $n=mM$ là đúng.}
{\True Khối lượng riêng của khí lí tưởng thỏa — phát biểu đúng là: $\rho=pM/(RT)$.}
{Khối lượng riêng của khí lí tưởng thỏa — phát biểu $\rho=pRT/M$ là đúng.}
\loigiai{
\begin{itemchoice}
\item (Đúng) Số mol khí được tính từ khối lượng m và khối lượng mol $M$ bằng — phát biểu đúng là: $n=m/M$. (Vì): Theo định nghĩa số mol, $n=m/M$. b) Sai: trái với hệ thức hoặc điều kiện đã nêu. c) Đúng: Từ $pV=(m/M)RT$ và $\rho=m/V$. d) Sai: trái với hệ thức hoặc điều kiện đã nêu
\item (Sai) Số mol khí được tính từ khối lượng m và khối lượng mol $M$ bằng — phát biểu $n=mM$ là đúng.
\item (Đúng) Khối lượng riêng của khí lí tưởng thỏa — phát biểu đúng là: $\rho=pM/(RT)$.
\item (Sai) Khối lượng riêng của khí lí tưởng thỏa — phát biểu $\rho=pRT/M$ là đúng.
\end{itemchoice}
}
\end{ex}

% ===== Câu 197 =====
% ID: L12C2B11-54-TLN
% Mức: VD
\begin{ex}
% ID: L12C2B11-54-TLN
% Mức: VD
Trong dạng “Áp dụng phương trình Clapeyron pV = nRT”, Ở cùng $T$ và $V$, áp suất tăng 3 lần. Số mol tăng bao nhiêu lần?
\shortans{3}
\loigiai{
Từ $pV=nRT$, khi $T$, $V$ cố định thì $p\sim n$.
}
\end{ex}

% ===== Câu 198 =====
% ID: L12C2B11-54-TN
% Mức: VD
\dangbt{Xác định sự biến thiên khối lượng riêng của chất khí}
\begin{ex}
% ID: L12C2B11-54-TN
% Mức: VD
Một khối khí có khối lượng $12\ \mathrm{g}$ ban đầu chiếm thể tích $4,0\ \mathrm{l}$ ở nhiệt độ $7\ ^\circ\mathrm{C}$. Sau khi được đun nóng đẳng áp, khối lượng riêng của khối khí này chỉ còn lại $1,2\ \mathrm{g/l}$. Nhiệt độ Celsius của khối khí ở trạng thái sau bằng bao nhiêu?
\choice
{$227\ ^\circ\mathrm{C}$}
{$350\ ^\circ\mathrm{C}$}
{\True $427\ ^\circ\mathrm{C}$}
{$700\ ^\circ\mathrm{C}$}
\loigiai{
Khối lượng riêng ban đầu của khối khí là:
  \begin{aligned}
 D_1 &= $\dfrac{m}{V_1}&=\dfrac{12}{4,0}&= 3,0\\mathrm{g/l}$.
 \end{aligned} 
 Áp dụng công thức của quá trình đẳng áp:
  \begin{aligned}
 T_2 &= T_1 $\cdot\dfrac{D_1}{D_2}$ &= (7 + 273) $\cdot\dfrac{3,0}{1,2}$ &= 700 $\\mathrm{K}$.
 \end{aligned} 
 Đổi nhiệt độ tuyệt đối sang nhiệt độ Celsius:
  \begin{aligned}
 t_2 &= T_2 - 273 
&= 700 - 273 
&= 427 $\$ ^ $\circ\mathrm{C}$.
 \end{aligned} 
 Chọn đáp án tương ứng.
}
\end{ex}

% ===== Câu 199 =====
% ID: L12C2B11-55-DS
% Mức: VDC
\dangbt{Áp dụng phương trình Clapeyron pV = nRT}
\begin{ex}
% ID: L12C2B11-55-DS
% Mức: VDC
Xét các phát biểu sau trong dạng “Áp dụng phương trình Clapeyron pV = nRT”.
\choiceTF
{\True Hai bình cùng p, $V$, $T$ chứa khí lí tưởng khác loại thì có — phát biểu đúng là: cùng số phân tử.}
{Hai bình cùng p, $V$, $T$ chứa khí lí tưởng khác loại thì có — phát biểu cùng khối lượng là đúng.}
{\True Khi giải bài toán khí lí tưởng, áp suất đưa vào phương trình phải là — phát biểu đúng là: áp suất tuyệt đối.}
{Khi giải bài toán khí lí tưởng, áp suất đưa vào phương trình phải là — phát biểu áp suất do chất lỏng là đúng.}
\loigiai{
\begin{itemchoice}
\item (Đúng) Hai bình cùng p, $V$, $T$ chứa khí lí tưởng khác loại thì có — phát biểu đúng là: cùng số phân tử. (Vì): Theo $n=pV/(RT)$, số mol và số phân tử bằng nhau. b) Sai: trái với hệ thức hoặc điều kiện đã nêu. c) Đúng: Phương trình trạng thái sử dụng áp suất tuyệt đối. d) Sai: trái với hệ thức hoặc điều kiện đã nêu
\item (Sai) Hai bình cùng p, $V$, $T$ chứa khí lí tưởng khác loại thì có — phát biểu cùng khối lượng là đúng.
\item (Đúng) Khi giải bài toán khí lí tưởng, áp suất đưa vào phương trình phải là — phát biểu đúng là: áp suất tuyệt đối.
\item (Sai) Khi giải bài toán khí lí tưởng, áp suất đưa vào phương trình phải là — phát biểu áp suất do chất lỏng là đúng.
\end{itemchoice}
}
\end{ex}

% ===== Câu 200 =====
% ID: L12C2B11-55-TLN
% Mức: VDC
\begin{ex}
% ID: L12C2B11-55-TLN
% Mức: VDC
Trong dạng “Áp dụng phương trình Clapeyron pV = nRT”, Khí chuyển trạng thái sao cho p tăng 4 lần và $T$ tăng 4 lần. Thể tích thay đổi bao nhiêu lần?
\shortans{1}
\loigiai{
$V_2/V_1=(T_2/T_1)/(p_2/p_1)=1$.
}
\end{ex}

% ===== Câu 201 =====
% ID: L12C2B11-55-TN
% Mức: VD
\dangbt{Xác định sự biến thiên khối lượng riêng của chất khí}
\begin{ex}
% ID: L12C2B11-55-TN
% Mức: VD
Phương trình phân hủy chất rắn trong túi khí ô tô là: $2\mathrm{NaN_3(r)} \xrightarrow{t^\circ} 2\mathrm{Na(r)} + 3\mathrm{N_2(k)}$. Nếu trong thiết bị túi khí ban đầu chứa $130\ \mathrm{g}\ \mathrm{NaN_3}$ nguyên chất thì sau phản ứng phân hủy hoàn toàn, lượng khí $\mathrm{N_2}$ giải phóng ra ngoài bằng bao nhiêu mol? Biết khối lượng mol của $\mathrm{Na}$ là $23\ \mathrm{g/mol}$ và của $\mathrm{N}$ là $14\ \mathrm{g/mol}$.
\choice
{$1,50\ \mathrm{mol}$}
{$2,00\ \mathrm{mol}$}
{\True $3,00\ \mathrm{mol}$}
{$4,50\ \mathrm{mol}$}
\loigiai{
Áp dụng công thức
  \begin{aligned}
 n_{ $\mathrm{N_2}} &=\dfrac{3}{2}\cdotn_{\mathrm{NaN_3}} 
&=\dfrac{3}{2}\cdot\dfrac{130}{23 + 14 \cdot 3}$&= 3,00$\\mathrm{mol}$.
 \end{aligned} 
 Chọn đáp án tương ứng.
}
\end{ex}

% ===== Câu 202 =====
% ID: L12C2B11-56-DS
% Mức: TH
\dangbt{Đồ thị và các quá trình biến đổi trạng thái liên tiếp}
\begin{ex}
% ID: L12C2B11-56-DS
% Mức: TH
Xét các phát biểu sau trong dạng “Đồ thị và các quá trình biến đổi trạng thái liên tiếp”.
\choiceTF
{\True Phương trình Clapeyron của khí lí tưởng là — phát biểu đúng là: $pV=nRT$.}
{Phương trình Clapeyron của khí lí tưởng là — phát biểu $pV=RT/n$ là đúng.}
{\True Với lượng khí xác định, phương trình trạng thái giữa hai trạng thái là — phát biểu đúng là: $p_1V_1/T_1=p_2V_2/T_2$.}
{Với lượng khí xác định, phương trình trạng thái giữa hai trạng thái là — phát biểu $p_1/T_1=p_2/T_2$ trong mọi quá trình là đúng.}
\loigiai{
\begin{itemchoice}
\item (Đúng) Phương trình Clapeyron của khí lí tưởng là — phát biểu đúng là: $pV=nRT$. (Vì): Một lượng khí lí tưởng thỏa $pV=nRT$. b) Sai: trái với hệ thức hoặc điều kiện đã nêu. c) Đúng: Do $pV/T=nR$ không đổi. d) Sai: trái với hệ thức hoặc điều kiện đã nêu
\item (Sai) Phương trình Clapeyron của khí lí tưởng là — phát biểu $pV=RT/n$ là đúng.
\item (Đúng) Với lượng khí xác định, phương trình trạng thái giữa hai trạng thái là — phát biểu đúng là: $p_1V_1/T_1=p_2V_2/T_2$.
\item (Sai) Với lượng khí xác định, phương trình trạng thái giữa hai trạng thái là — phát biểu $p_1/T_1=p_2/T_2$ trong mọi quá trình là đúng.
\end{itemchoice}
}
\end{ex}

% ===== Câu 203 =====
% ID: L12C2B11-56-TLN
% Mức: TH
\begin{ex}
% ID: L12C2B11-56-TLN
% Mức: TH
Trong dạng “Đồ thị và các quá trình biến đổi trạng thái liên tiếp”, Có 3 mol khí ở 300 $K$ chiếm 8 lít. Tính áp suất theo kPa, lấy R=8,31.
\shortans{934,88}
\loigiai{
$p=nRT/V=934,88$ kPa.
}
\end{ex}

% ===== Câu 204 =====
% ID: L12C2B11-56-TN
% Mức: NB
\dangbt{Xác định sự biến thiên khối lượng riêng của chất khí}
\begin{ex}
% ID: L12C2B11-56-TN
% Mức: NB
Đối với một lượng khí lí tưởng xác định, đại lượng nào sau đây không đổi trong mọi quá trình biến đổi trạng thái?
\choice
{Tích số $p \cdot V$}
{Thương số $\dfrac{p}{T}$}
{\True Thương số $\dfrac{p \cdot V}{T}$}
{Thương số $\dfrac{V}{T}$}
\loigiai{
Theo phương trình trạng thái của khí lí tưởng cho một khối lượng khí xác định, ta luôn có hệ thức $\dfrac{p \cdot V}{T} = \text{const}$. Do đó thương số $\dfrac{p \cdot V}{T}$ luôn không đổi.
 Chọn đáp án tương ứng.
}
\end{ex}

% ===== Câu 205 =====
% ID: L12C2B11-57-DS
% Mức: TH
\dangbt{Đồ thị và các quá trình biến đổi trạng thái liên tiếp}
\begin{ex}
% ID: L12C2B11-57-DS
% Mức: TH
Xét các phát biểu sau trong dạng “Đồ thị và các quá trình biến đổi trạng thái liên tiếp”.
\choiceTF
{\True Có 4 mol khí ở 280 $K$ chiếm 2 lít. Lấy R= $8,31\,\text{J}$ /(mol.K). Áp suất gần bằng — phát biểu đúng là: $4653,6\,\text{kPa}$.}
{Có 4 mol khí ở 280 $K$ chiếm 2 lít. Lấy R= $8,31\,\text{J}$ /(mol.K). Áp suất gần bằng — phát biểu $465,36\,\text{kPa}$ là đúng.}
{\True Khí chuyển từ $p_1=100$ kPa, $V_1=4 lít,$ T_1=290 $K đến$ V_2=8 $lít,$ T_2=370 $K.$ p_2 $gần bằng — phát biểu đúng là: 63,79 kPa.$}
{Khí chuyển từ $p_1=100$ kPa, $V_1=4 lít,$ T_1=290 $K đến$ V_2=8 $lít,$ T_2=370 $K.$ p_2 $gần bằng — phát biểu 78,38 kPa là đúng.$}
\loigiai{
\begin{itemchoice}
\item (Đúng) Có 4 mol khí ở 280 $K$ chiếm 2 lít. Lấy R= $8,31\,\text{J}$ /(mol.K). Áp suất gần bằng — phát biểu đúng là: $4653,6\,\text{kPa}$. (Vì): $p=nRT/V=4653,6$ kPa. b) Sai: trái với hệ thức hoặc điều kiện đã nêu. c) Đúng: $p_2=p_1V_1T_2/(V_2T_1)=63,79$ kPa. d) Sai: trái với hệ thức hoặc điều kiện đã nêu
\item (Sai) Có 4 mol khí ở 280 $K$ chiếm 2 lít. Lấy R= $8,31\,\text{J}$ /(mol.K). Áp suất gần bằng — phát biểu $465,36\,\text{kPa}$ là đúng.
\item (Đúng) Khí chuyển từ $p_1=100$ kPa, $V_1=4 lít,$ T_1=290 $K đến$ V_2=8 $lít,$ T_2=370 $K.$ p_2 $gần bằng — phát biểu đúng là: 63,79 kPa.$.
\item (Sai) Khí chuyển từ $p_1=100$ kPa, $V_1=4 lít,$ T_1=290 $K đến$ V_2=8 $lít,$ T_2=370 $K.$ p_2 $gần bằng — phát biểu 78,38 kPa là đúng.$.
\end{itemchoice}
}
\end{ex}

% ===== Câu 206 =====
% ID: L12C2B11-57-TLN
% Mức: TH
\begin{ex}
% ID: L12C2B11-57-TLN
% Mức: TH
Trong dạng “Đồ thị và các quá trình biến đổi trạng thái liên tiếp”, Khí có $100\,\text{kPa}$, 4 lít ở 310 K. Đun đẳng tích đến 370 K. Tính áp suất sau theo kPa.
\shortans{119,35}
\loigiai{
$p_2=p_1T_2/T_1=119,35$ kPa.
}
\end{ex}

% ===== Câu 207 =====
% ID: L12C2B11-57-TN
% Mức: TH
\dangbt{Xác định sự biến thiên khối lượng riêng của chất khí}
\begin{ex}
% ID: L12C2B11-57-TN
% Mức: TH
Bóng thám không chỉ có thể bay lên cao trong khí quyển khi thỏa mãn điều kiện nào sau đây?
\choice
{\True Khối lượng riêng trung bình của toàn bộ quả bóng nhỏ hơn khối lượng riêng của không khí bên ngoài}
{Khối lượng riêng của khí dùng để bơm bóng lớn hơn khối lượng riêng của không khí bên ngoài}
{Áp suất khí bên trong vỏ bóng nhỏ hơn áp suất khí quyển ở bên ngoài}
{Áp suất khí bên trong vỏ bóng lớn hơn áp suất khí quyển ở bên ngoài}
\loigiai{
Để bóng thám không có thể tự bay lên cao, lực nâng Archimedes tác dụng lên nó phải lớn hơn trọng lượng của bóng, tức là khối lượng riêng trung bình của toàn bộ quả bóng (gồm vỏ bóng, thiết bị và khí bơm bên trong) phải nhỏ hơn khối lượng riêng của không khí bên ngoài.
 Chọn đáp án tương ứng.
}
\end{ex}

% ===== Câu 208 =====
% ID: L12C2B11-58-DS
% Mức: VD
\dangbt{Đồ thị và các quá trình biến đổi trạng thái liên tiếp}
\begin{ex}
% ID: L12C2B11-58-DS
% Mức: VD
Xét các phát biểu sau trong dạng “Đồ thị và các quá trình biến đổi trạng thái liên tiếp”.
\choiceTF
{\True Trong phương trình khí lí tưởng, nhiệt độ phải dùng — phát biểu đúng là: kelvin.}
{Trong phương trình khí lí tưởng, nhiệt độ phải dùng — phát biểu độ $C$ là đúng.}
{\True Ở thể tích không đổi, áp suất của lượng khí xác định tỉ lệ — phát biểu đúng là: thuận với nhiệt độ tuyệt đối.}
{Ở thể tích không đổi, áp suất của lượng khí xác định tỉ lệ — phát biểu thuận với bình phương nhiệt độ là đúng.}
\loigiai{
\begin{itemchoice}
\item (Đúng) Trong phương trình khí lí tưởng, nhiệt độ phải dùng — phát biểu đúng là: kelvin. (Vì): $T$ phải là nhiệt độ tuyệt đối. b) Sai: trái với hệ thức hoặc điều kiện đã nêu. c) Đúng: Từ $pV=nRT$, khi $V$ cố định thì $p/T$ không đổi. d) Sai: trái với hệ thức hoặc điều kiện đã nêu
\item (Sai) Trong phương trình khí lí tưởng, nhiệt độ phải dùng — phát biểu độ $C$ là đúng.
\item (Đúng) Ở thể tích không đổi, áp suất của lượng khí xác định tỉ lệ — phát biểu đúng là: thuận với nhiệt độ tuyệt đối.
\item (Sai) Ở thể tích không đổi, áp suất của lượng khí xác định tỉ lệ — phát biểu thuận với bình phương nhiệt độ là đúng.
\end{itemchoice}
}
\end{ex}

% ===== Câu 209 =====
% ID: L12C2B11-58-TLN
% Mức: VD
\begin{ex}
% ID: L12C2B11-58-TLN
% Mức: VD
Trong dạng “Đồ thị và các quá trình biến đổi trạng thái liên tiếp”, Khí có $110\,\text{kPa}$, 6 lít ở 320 $K$; chuyển đến $330\,\text{kPa}$ và 380 K. Tính thể tích sau theo lít.
\shortans{2,38}
\loigiai{
$V_2=p_1V_1T_2/(p_2T_1)=2,38$ lít.
}
\end{ex}

% ===== Câu 210 =====
% ID: L12C2B11-58-TN
% Mức: VD
\dangbt{Xác định sự biến thiên khối lượng riêng của chất khí}
\begin{ex}
% ID: L12C2B11-58-TN
% Mức: VD
Một khối khí lí tưởng xác định ở nhiệt độ $27\ ^\circ\mathrm{C}$ có khối lượng riêng là $1,20\ \mathrm{kg/m^3}$. Nếu đun nóng đẳng áp khối khí này đến nhiệt độ $127\ ^\circ\mathrm{C}$ thì khối lượng riêng của khối khí lúc này bằng bao nhiêu?
\choice
{$1,60\ \mathrm{kg/m^3}$}
{$1,00\ \mathrm{kg/m^3}$}
{\True $0,90\ \mathrm{kg/m^3}$}
{$0,80\ \mathrm{kg/m^3}$}
\loigiai{
Áp dụng công thức
  \begin{aligned}
 D_2 &= D_1 $\cdot\dfrac{T_1}{T_2}&= 1,20\cdot\dfrac{27 + 273}{127 + 273}&= 0,90\\mathrm{kg/m^3}$.
 \end{aligned} 
 Chọn đáp án tương ứng.
}
\end{ex}

% ===== Câu 211 =====
% ID: L12C2B11-59-DS
% Mức: VD
\dangbt{Đồ thị và các quá trình biến đổi trạng thái liên tiếp}
\begin{ex}
% ID: L12C2B11-59-DS
% Mức: VD
Xét các phát biểu sau trong dạng “Đồ thị và các quá trình biến đổi trạng thái liên tiếp”.
\choiceTF
{\True Số mol khí được tính từ khối lượng m và khối lượng mol $M$ bằng — phát biểu đúng là: $n=m/M$.}
{Số mol khí được tính từ khối lượng m và khối lượng mol $M$ bằng — phát biểu $n=mM$ là đúng.}
{\True Khối lượng riêng của khí lí tưởng thỏa — phát biểu đúng là: $\rho=pM/(RT)$.}
{Khối lượng riêng của khí lí tưởng thỏa — phát biểu $\rho=pRT/M$ là đúng.}
\loigiai{
\begin{itemchoice}
\item (Đúng) Số mol khí được tính từ khối lượng m và khối lượng mol $M$ bằng — phát biểu đúng là: $n=m/M$. (Vì): Theo định nghĩa số mol, $n=m/M$. b) Sai: trái với hệ thức hoặc điều kiện đã nêu. c) Đúng: Từ $pV=(m/M)RT$ và $\rho=m/V$. d) Sai: trái với hệ thức hoặc điều kiện đã nêu
\item (Sai) Số mol khí được tính từ khối lượng m và khối lượng mol $M$ bằng — phát biểu $n=mM$ là đúng.
\item (Đúng) Khối lượng riêng của khí lí tưởng thỏa — phát biểu đúng là: $\rho=pM/(RT)$.
\item (Sai) Khối lượng riêng của khí lí tưởng thỏa — phát biểu $\rho=pRT/M$ là đúng.
\end{itemchoice}
}
\end{ex}

% ===== Câu 212 =====
% ID: L12C2B11-59-TLN
% Mức: VD
\begin{ex}
% ID: L12C2B11-59-TLN
% Mức: VD
Trong dạng “Đồ thị và các quá trình biến đổi trạng thái liên tiếp”, Một khí có khối lượng $84\,\text{g}$, khối lượng mol $28\,\text{g}$ /mol. Tính số mol.
\shortans{3}
\loigiai{
$n=m/M=84/28=3$ mol.
}
\end{ex}

% ===== Câu 213 =====
% ID: L12C2B11-59-TN
% Mức: TH
\dangbt{Xác định sự biến thiên khối lượng riêng của chất khí}
\begin{ex}
% ID: L12C2B11-59-TN
% Mức: TH
Hằng số khí lí tưởng $R$ trong hệ SI có giá trị xấp xỉ bằng $8,31\ \mathrm{J/(mol\cdot K)}$. Đơn vị của $R$ được xác định dựa trên tích của các đơn vị đo nào sau đây?
\choice
{$\mathrm{Pa \cdot m^3 \cdot mol \cdot K}$}
{$\mathrm{Pa \cdot m^{-3} \cdot mol^{-1} \cdot K^{-1}}$}
{\True $\mathrm{Pa \cdot m^3 \cdot mol^{-1} \cdot K^{-1}}$}
{$\mathrm{Pa \cdot m^{-3} \cdot mol \cdot K^{-1}}$}
\loigiai{
Từ phương trình Clapeyron: $p \cdot V = n \cdot R \cdot T \Rightarrow R = \dfrac{p \cdot V}{n \cdot T}$. Do đó, đơn vị của $R$ là $\mathrm{Pa \cdot m^3/(mol\cdot K)}$ hay $\mathrm{Pa \cdot m^3 \cdot mol^{-1} \cdot K^{-1}}$.
 Chọn đáp án tương ứng.
}
\end{ex}

% ===== Câu 214 =====
% ID: L12C2B11-60-DS
% Mức: VDC
\dangbt{Đồ thị và các quá trình biến đổi trạng thái liên tiếp}
\begin{ex}
% ID: L12C2B11-60-DS
% Mức: VDC
Xét các phát biểu sau trong dạng “Đồ thị và các quá trình biến đổi trạng thái liên tiếp”.
\choiceTF
{\True Hai bình cùng p, $V$, $T$ chứa khí lí tưởng khác loại thì có — phát biểu đúng là: cùng số phân tử.}
{Hai bình cùng p, $V$, $T$ chứa khí lí tưởng khác loại thì có — phát biểu cùng khối lượng là đúng.}
{\True Khi giải bài toán khí lí tưởng, áp suất đưa vào phương trình phải là — phát biểu đúng là: áp suất tuyệt đối.}
{Khi giải bài toán khí lí tưởng, áp suất đưa vào phương trình phải là — phát biểu áp suất do chất lỏng בלבד là đúng.}
\loigiai{
\begin{itemchoice}
\item (Đúng) Hai bình cùng p, $V$, $T$ chứa khí lí tưởng khác loại thì có — phát biểu đúng là: cùng số phân tử. (Vì): Theo $n=pV/(RT)$, số mol và số phân tử bằng nhau. b) Sai: trái với hệ thức hoặc điều kiện đã nêu. c) Đúng: Phương trình trạng thái sử dụng áp suất tuyệt đối. d) Sai: trái với hệ thức hoặc điều kiện đã nêu
\item (Sai) Hai bình cùng p, $V$, $T$ chứa khí lí tưởng khác loại thì có — phát biểu cùng khối lượng là đúng.
\item (Đúng) Khi giải bài toán khí lí tưởng, áp suất đưa vào phương trình phải là — phát biểu đúng là: áp suất tuyệt đối.
\item (Sai) Khi giải bài toán khí lí tưởng, áp suất đưa vào phương trình phải là — phát biểu áp suất do chất lỏng בלבד là đúng.
\end{itemchoice}
}
\end{ex}

% ===== Câu 215 =====
% ID: L12C2B11-60-TLN
% Mức: VD
\begin{ex}
% ID: L12C2B11-60-TLN
% Mức: VD
Trong dạng “Đồ thị và các quá trình biến đổi trạng thái liên tiếp”, Ở cùng $T$ và $V$, áp suất tăng 2 lần. Số mol tăng bao nhiêu lần?
\shortans{2}
\loigiai{
Từ $pV=nRT$, khi $T$, $V$ cố định thì $p\sim n$.
}
\end{ex}

% ===== Câu 216 =====
% ID: L12C2B11-60-TN
% Mức: VD
\dangbt{Xác định sự biến thiên khối lượng riêng của chất khí}
\begin{ex}
% ID: L12C2B11-60-TN
% Mức: VD
Trong túi khí bảo vệ người lái của một ô tô có chứa $100\ \mathrm{g}$ chất rắn $\mathrm{NaN_3}$. Khi xe xảy ra va chạm mạnh, phản ứng phân hủy $\mathrm{NaN_3}$ xảy ra hoàn toàn tạo thành kim loại natri và giải phóng khí nitrogen ($\mathrm{N_2}$). Biết thể tích mol khí trong điều kiện phản ứng là $24,0\ \mathrm{l/mol}$. Thể tích khí $\mathrm{N_2}$ được giải phóng ra bằng bao nhiêu?
\choice
{$36,96\ \mathrm{l}$}
{$48,00\ \mathrm{l}$}
{\True $55,44\ \mathrm{l}$}
{$64,80\ \mathrm{l}$}
\loigiai{
Áp dụng công thức
  \begin{aligned}
 V_{ $\mathrm{N_2}} &=\dfrac{3}{2}\cdot\dfrac{m_{\mathrm{NaN_3}}}{M_{\mathrm{NaN_3}}}\cdotV_{\text{mol}} 
&=\dfrac{3}{2}\cdot\dfrac{100}{65}\cdot$ 24,0 
&= 55,44 $\\mathrm{l}$.
 \end{aligned} 
 Chọn đáp án tương ứng.
}
\end{ex}

% ===== Câu 217 =====
% ID: L12C2B11-61-TLN
% Mức: VDC
\dangbt{Đồ thị và các quá trình biến đổi trạng thái liên tiếp}
\begin{ex}
% ID: L12C2B11-61-TLN
% Mức: VDC
Trong dạng “Đồ thị và các quá trình biến đổi trạng thái liên tiếp”, Khí chuyển trạng thái sao cho p tăng 3 lần và $T$ tăng 3 lần. Thể tích thay đổi bao nhiêu lần?
\shortans{1}
\loigiai{
$V_2/V_1=(T_2/T_1)/(p_2/p_1)=1$.
}
\end{ex}

% ===== Câu 218 =====
% ID: L12C2B11-61-TN
% Mức: TH
\dangbt{Xác định sự biến thiên khối lượng riêng của chất khí}
\begin{ex}
% ID: L12C2B11-61-TN
% Mức: TH
Trong hiện tượng nào sau đây, cả ba thông số trạng thái của một lượng khí xác định đều thay đổi?
\choice
{Không khí bị đun nóng trong một bình kín}
{\True Không khí bên trong quả bóng bàn bị bẹp nhúng vào nước nóng và phồng lên}
{Nén không khí từ từ trong một xi lanh cách nhiệt tốt}
{Không khí bên trong quả bóng bay bị bóp bẹp dị dạng nhưng chưa bị vỡ}
\loigiai{
Khi nhúng quả bóng bàn bị bẹp vào nước nóng, không khí trong bóng nóng lên (nhiệt độ $T$ tăng), nở ra làm bóng phồng lên (thể tích $V$ tăng), và áp suất khí bên trong tăng để thắng lực đàn hồi của vỏ bóng (áp suất $p$ tăng). Do đó cả ba thông số trạng thái đều thay đổi.
 Chọn đáp án tương ứng.
}
\end{ex}

% ===== Câu 219 =====
% ID: L12C2B11-62-TN
% Mức: VDC
\begin{ex}
% ID: L12C2B11-62-TN
% Mức: VDC
Để một túi khí xe ô tô có thể tích khi phồng lên bằng $49,86\ \mathrm{l}$ đạt áp suất $1,2 \cdot 10^5\ \mathrm{Pa}$ ở nhiệt độ $27\ ^\circ\mathrm{C}$ thì cần nạp tối thiểu bao nhiêu gam chất rắn $\mathrm{NaN_3}$ nguyên chất vào thiết bị?
\choice
{$85\ \mathrm{g}$}
{\True $104\ \mathrm{g}$}
{$115\ \mathrm{g}$}
{$130\ \mathrm{g}$}
\loigiai{
Áp dụng công thức
  \begin{aligned}
 n_{ $\mathrm{N_2}} &=\dfrac{p \cdot V}{R \cdot T}&=\dfrac{1,2 \cdot 10^5 \cdot 49,86 \cdot 10^{-3}}{8,31 \cdot (27 + 273)}&= 2,40\\mathrm{mol}$.
 \end{aligned} 
 Số mol $\mathrm{NaN_3}$ cần dùng và khối lượng của nó là:
  \begin{aligned}
 n_{ $\mathrm{NaN_3}$ } &= $\dfrac{2}{3}\cdot$ n_{ $\mathrm{N_2}$ } = 1,60 $\\mathrm{mol}$ m_{ $\mathrm{NaN_3}$ } &= n_{ $\mathrm{NaN_3}$ } \cdot M_{ $\mathrm{NaN_3}$ } 
&= 1,60 \cdot 65 
&= 104 $\\mathrm{g}$.
 \end{aligned} 
 Chọn đáp án tương ứng.
}
\end{ex}

% ===== Câu 220 =====
% ID: L12C2B11-63-TN
% Mức: VD
\begin{ex}
% ID: L12C2B11-63-TN
% Mức: VD
Một bóng thám không được bơm ở mặt đất dưới áp suất $1,0 \cdot 10^5\ \mathrm{Pa}$ và nhiệt độ $300\ \mathrm{K}$ có thể tích ban đầu là $4,0\ \mathrm{m^3}$. Biết vỏ bóng làm bằng vật liệu đàn hồi chịu được thể tích dãn nở tối đa là $12\ \mathrm{m^3}$ trước khi bị vỡ. Khi bóng bay lên tầng khí quyển cao có nhiệt độ $225\ \mathrm{K}$, áp suất khí quyển giảm dần. Áp suất khí quyển tối thiểu tại độ cao mà bóng bắt đầu bị vỡ bằng bao nhiêu?
\choice
{$0,15 \cdot 10^5\ \mathrm{Pa}$}
{$0,20 \cdot 10^5\ \mathrm{Pa}$}
{\True $0,25 \cdot 10^5\ \mathrm{Pa}$}
{$0,30 \cdot 10^5\ \mathrm{Pa}$}
\loigiai{
Áp dụng công thức
  \begin{aligned}
 p_2 &= p_1 $\cdot\dfrac{V_1}{V_2}\cdot\dfrac{T_2}{T_1}&= 1,0\cdot10^5\cdot\dfrac{4,0}{12}\cdot\dfrac{225}{300}&= 0,25\cdot10^5\\mathrm{Pa}$.
 \end{aligned} 
 Chọn đáp án tương ứng.
}
\end{ex}

% ===== Câu 221 =====
% ID: L12C2B11-64-TN
% Mức: VD
\begin{ex}
% ID: L12C2B11-64-TN
% Mức: VD
Một bóng thám không có thể tích khi bay ở tầng khí quyển cao có áp suất $0,25 \cdot 10^5\ \mathrm{Pa}$ và nhiệt độ $220\ \mathrm{K}$ là $88,0\ \mathrm{m^3}$. Hỏi thể tích của quả bóng thám không này khi vừa bơm xong ở mặt đất dưới áp suất $1,00 \cdot 10^5\ \mathrm{Pa}$ và nhiệt độ $275\ \mathrm{K}$ là bao nhiêu?
\choice
{$18,5\ \mathrm{m^3}$}
{$22,0\ \mathrm{m^3}$}
{\True $27,5\ \mathrm{m^3}$}
{$35,0\ \mathrm{m^3}$}
\loigiai{
Áp dụng công thức
  \begin{aligned}
 V_1 &= V_2 $\cdot\dfrac{p_2}{p_1}\cdot\dfrac{T_1}{T_2}&= 88,0\cdot\dfrac{0,25 \cdot 10^5}{1,00 \cdot 10^5}\cdot\dfrac{275}{220}&= 27,5\\mathrm{m^3}$.
 \end{aligned} 
 Chọn đáp án tương ứng.
}
\end{ex}

% ===== Câu 222 =====
% ID: L12C2B11-65-TN
% Mức: VD
\begin{ex}
% ID: L12C2B11-65-TN
% Mức: VD
Một xilanh chứa $100\ \mathrm{cm^3}$ khí lí tưởng ở nhiệt độ $27\ ^\circ\mathrm{C}$ và áp suất $1,0 \cdot 10^5\ \mathrm{Pa}$. Khi pit-tông nén khí đến thể tích $25\ \mathrm{cm^3}$ thì nhiệt độ khí tăng lên đến $87\ ^\circ\mathrm{C}$. Áp suất của khí trong xilanh lúc này bằng bao nhiêu?
\choice
{$3,6 \cdot 10^5\ \mathrm{Pa}$}
{$4,0 \cdot 10^5\ \mathrm{Pa}$}
{\True $4,8 \cdot 10^5\ \mathrm{Pa}$}
{$5,4 \cdot 10^5\ \mathrm{Pa}$}
\loigiai{
Áp dụng công thức
  \begin{aligned}
 p_2 &= $\dfrac{p_1 \cdot V_1 \cdot T_2}{V_2 \cdot T_1}&=\dfrac{1,0 \cdot 10^5 \cdot 100 \cdot (87 + 273)}{25 \cdot (27 + 273)}&= 4,8\cdot10^5\\mathrm{Pa}$.
 \end{aligned} 
 Chọn đáp án tương ứng.
}
\end{ex}

% ===== Câu 223 =====
% ID: L12C2B11-66-TN
% Mức: VD
\begin{ex}
% ID: L12C2B11-66-TN
% Mức: VD
Một lốp xe máy có dung tích $0,020\ \mathrm{m^3}$ chứa không khí ở nhiệt độ $27\ ^\circ\mathrm{C}$ và áp suất $3,0 \cdot 10^5\ \mathrm{Pa}$. Biết khối lượng mol của không khí là $28,8\ \mathrm{g/mol}$. Khối lượng của lượng không khí chứa trong lốp xe bằng bao nhiêu?
\choice
{$58,5\ \mathrm{g}$}
{\True $69,3\ \mathrm{g}$}
{$72,4\ \mathrm{g}$}
{$81,2\ \mathrm{g}$}
\loigiai{
Áp dụng công thức
  \begin{aligned}
 m &= $\dfrac{p \cdot V \cdot M}{R \cdot T}&=\dfrac{3,0 \cdot 10^5 \cdot 0,020 \cdot 28,8 \cdot 10^{-3}}{8,31 \cdot (27 + 273)}$ &\approx 0,0693 $\\mathrm{kg}$ = 69,3 $\\mathrm{g}$.
 \end{aligned} 
 Chọn đáp án tương ứng.
}
\end{ex}

% ===== Câu 224 =====
% ID: L12C2B11-67-TN
% Mức: NB
\dangbt{Áp dụng phương trình Clapeyron pV = nRT}
\begin{ex}
% ID: L12C2B11-67-TN
% Mức: NB
Trong dạng “Áp dụng phương trình Clapeyron pV = nRT”, Phương trình Clapeyron của khí lí tưởng là
\choice
{\True $pV=nRT$}
{$pV=RT/n$}
{$p/T=nV$}
{$p=nR/(VT)$}
\loigiai{
Đáp án A. Một lượng khí lí tưởng thỏa $pV=nRT$.
}
\end{ex}

% ===== Câu 225 =====
% ID: L12C2B11-68-TN
% Mức: NB
\begin{ex}
% ID: L12C2B11-68-TN
% Mức: NB
Trong dạng “Áp dụng phương trình Clapeyron pV = nRT”, Với lượng khí xác định, phương trình trạng thái giữa hai trạng thái là
\choice
{\True $p_1V_1/T_1=p_2V_2/T_2$}
{$p_1V_1T_1=p_2V_2T_2$}
{$p_1/T_1=p_2/T_2$ trong mọi quá trình}
{$V_1/T_1=V_2/T_2$ trong mọi quá trình}
\loigiai{
Đáp án A. Do $pV/T=nR$ không đổi.
}
\end{ex}

% ===== Câu 226 =====
% ID: L12C2B11-69-TN
% Mức: TH
\begin{ex}
% ID: L12C2B11-69-TN
% Mức: TH
Trong dạng “Áp dụng phương trình Clapeyron pV = nRT”, Có 2 mol khí ở 260 $K$ chiếm 7 lít. Lấy R= $8,31\,\text{J}$ /(mol.K). Áp suất gần bằng
\choice
{\True $617,31\,\text{kPa}$}
{$61,73\,\text{kPa}$}
{$6173,14\,\text{kPa}$}
{$0,06\,\text{kPa}$}
\loigiai{
Đáp án A. $p=nRT/V=617,31$ kPa.
}
\end{ex}

% ===== Câu 227 =====
% ID: L12C2B11-70-TN
% Mức: TH
\begin{ex}
% ID: L12C2B11-70-TN
% Mức: TH
Trong dạng “Áp dụng phương trình Clapeyron pV = nRT”, Khí chuyển từ $p_1=90$ kPa, $V_1=2 lít,$ T_1=270 $K đến$ V_2=8 $lít,$ T_2=350 $K.$ p_2 $gần bằng$
\choice
{\True $29,17\,\text{kPa}$}
{$466,67\,\text{kPa}$}
{$69,43\,\text{kPa}$}
{$90\,\text{kPa}$}
\loigiai{
Đáp án A. $p_2=p_1V_1T_2/(V_2T_1)=29,17$ kPa.
}
\end{ex}

% ===== Câu 228 =====
% ID: L12C2B11-71-TN
% Mức: TH
\begin{ex}
% ID: L12C2B11-71-TN
% Mức: TH
Trong dạng “Áp dụng phương trình Clapeyron pV = nRT”, Trong phương trình khí lí tưởng, nhiệt độ phải dùng
\choice
{\True kelvin}
{độ $C$}
{độ $F$}
{thang tùy ý}
\loigiai{
Đáp án A. $T$ phải là nhiệt độ tuyệt đối.
}
\end{ex}

% ===== Câu 229 =====
% ID: L12C2B11-72-TN
% Mức: VD
\begin{ex}
% ID: L12C2B11-72-TN
% Mức: VD
Trong dạng “Áp dụng phương trình Clapeyron pV = nRT”, Ở thể tích không đổi, áp suất của lượng khí xác định tỉ lệ
\choice
{\True thuận với nhiệt độ tuyệt đối}
{nghịch với nhiệt độ}
{thuận với bình phương nhiệt độ}
{không phụ thuộc nhiệt độ}
\loigiai{
Đáp án A. Từ $pV=nRT$, khi $V$ cố định thì $p/T$ không đổi.
}
\end{ex}

% ===== Câu 230 =====
% ID: L12C2B11-73-TN
% Mức: VD
\begin{ex}
% ID: L12C2B11-73-TN
% Mức: VD
Trong dạng “Áp dụng phương trình Clapeyron pV = nRT”, Số mol khí được tính từ khối lượng m và khối lượng mol $M$ bằng
\choice
{\True $n=m/M$}
{$n=mM$}
{$n=M/m$}
{$n=m+M$}
\loigiai{
Đáp án A. Theo định nghĩa số mol, $n=m/M$.
}
\end{ex}

% ===== Câu 231 =====
% ID: L12C2B11-74-TN
% Mức: VD
\begin{ex}
% ID: L12C2B11-74-TN
% Mức: VD
Trong dạng “Áp dụng phương trình Clapeyron pV = nRT”, Khối lượng riêng của khí lí tưởng thỏa
\choice
{\True $\rho=pM/(RT)$}
{$\rho=RT/(pM)$}
{$\rho=pRT/M$}
{$\rho=p/(MRT)$}
\loigiai{
Đáp án A. Từ $pV=(m/M)RT$ và $\rho=m/V$.
}
\end{ex}

% ===== Câu 232 =====
% ID: L12C2B11-75-TN
% Mức: VDC
\begin{ex}
% ID: L12C2B11-75-TN
% Mức: VDC
Trong dạng “Áp dụng phương trình Clapeyron pV = nRT”, Hai bình cùng p, $V$, $T$ chứa khí lí tưởng khác loại thì có
\choice
{\True cùng số phân tử}
{cùng khối lượng}
{cùng khối lượng riêng}
{khác số mol}
\loigiai{
Đáp án A. Theo $n=pV/(RT)$, số mol và số phân tử bằng nhau.
}
\end{ex}

% ===== Câu 233 =====
% ID: L12C2B11-76-TN
% Mức: VDC
\begin{ex}
% ID: L12C2B11-76-TN
% Mức: VDC
Trong dạng “Áp dụng phương trình Clapeyron pV = nRT”, Khi giải bài toán khí lí tưởng, áp suất đưa vào phương trình phải là
\choice
{\True áp suất tuyệt đối}
{áp suất dư}
{áp suất do chất lỏng בלבד}
{áp suất tương đối tùy ý}
\loigiai{
Đáp án A. Phương trình trạng thái sử dụng áp suất tuyệt đối.
}
\end{ex}

% ===== Câu 234 =====
% ID: L12C2B11-77-TN
% Mức: NB
\dangbt{Đồ thị và các quá trình biến đổi trạng thái liên tiếp}
\begin{ex}
% ID: L12C2B11-77-TN
% Mức: NB
Trong dạng “Đồ thị và các quá trình biến đổi trạng thái liên tiếp”, Phương trình Clapeyron của khí lí tưởng là
\choice
{\True $pV=nRT$}
{$pV=RT/n$}
{$p/T=nV$}
{$p=nR/(VT)$}
\loigiai{
Đáp án A. Một lượng khí lí tưởng thỏa $pV=nRT$.
}
\end{ex}

% ===== Câu 235 =====
% ID: L12C2B11-78-TN
% Mức: NB
\begin{ex}
% ID: L12C2B11-78-TN
% Mức: NB
Trong dạng “Đồ thị và các quá trình biến đổi trạng thái liên tiếp”, Với lượng khí xác định, phương trình trạng thái giữa hai trạng thái là
\choice
{\True $p_1V_1/T_1=p_2V_2/T_2$}
{$p_1V_1T_1=p_2V_2T_2$}
{$p_1/T_1=p_2/T_2$ trong mọi quá trình}
{$V_1/T_1=V_2/T_2$ trong mọi quá trình}
\loigiai{
Đáp án A. Do $pV/T=nR$ không đổi.
}
\end{ex}

% ===== Câu 236 =====
% ID: L12C2B11-79-TN
% Mức: TH
\begin{ex}
% ID: L12C2B11-79-TN
% Mức: TH
Trong dạng “Đồ thị và các quá trình biến đổi trạng thái liên tiếp”, Có 4 mol khí ở 280 $K$ chiếm 2 lít. Lấy R= $8,31\,\text{J}$ /(mol.K). Áp suất gần bằng
\choice
{\True $4653,6\,\text{kPa}$}
{$465,36\,\text{kPa}$}
{$46536\,\text{kPa}$}
{$0,12\,\text{kPa}$}
\loigiai{
Đáp án A. $p=nRT/V=4653,6$ kPa.
}
\end{ex}

% ===== Câu 237 =====
% ID: L12C2B11-80-TN
% Mức: TH
\begin{ex}
% ID: L12C2B11-80-TN
% Mức: TH
Trong dạng “Đồ thị và các quá trình biến đổi trạng thái liên tiếp”, Khí chuyển từ $p_1=100$ kPa, $V_1=4 lít,$ T_1=290 $K đến$ V_2=8 $lít,$ T_2=370 $K.$ p_2 $gần bằng$
\choice
{\True $63,79\,\text{kPa}$}
{$255,17\,\text{kPa}$}
{$78,38\,\text{kPa}$}
{$100\,\text{kPa}$}
\loigiai{
Đáp án A. $p_2=p_1V_1T_2/(V_2T_1)=63,79$ kPa.
}
\end{ex}

% ===== Câu 238 =====
% ID: L12C2B11-81-TN
% Mức: TH
\begin{ex}
% ID: L12C2B11-81-TN
% Mức: TH
Trong dạng “Đồ thị và các quá trình biến đổi trạng thái liên tiếp”, Trong phương trình khí lí tưởng, nhiệt độ phải dùng
\choice
{\True kelvin}
{độ $C$}
{độ $F$}
{thang tùy ý}
\loigiai{
Đáp án A. $T$ phải là nhiệt độ tuyệt đối.
}
\end{ex}

% ===== Câu 239 =====
% ID: L12C2B11-82-TN
% Mức: VD
\begin{ex}
% ID: L12C2B11-82-TN
% Mức: VD
Trong dạng “Đồ thị và các quá trình biến đổi trạng thái liên tiếp”, Ở thể tích không đổi, áp suất của lượng khí xác định tỉ lệ
\choice
{\True thuận với nhiệt độ tuyệt đối}
{nghịch với nhiệt độ}
{thuận với bình phương nhiệt độ}
{không phụ thuộc nhiệt độ}
\loigiai{
Đáp án A. Từ $pV=nRT$, khi $V$ cố định thì $p/T$ không đổi.
}
\end{ex}

% ===== Câu 240 =====
% ID: L12C2B11-83-TN
% Mức: VD
\begin{ex}
% ID: L12C2B11-83-TN
% Mức: VD
Trong dạng “Đồ thị và các quá trình biến đổi trạng thái liên tiếp”, Số mol khí được tính từ khối lượng m và khối lượng mol $M$ bằng
\choice
{\True $n=m/M$}
{$n=mM$}
{$n=M/m$}
{$n=m+M$}
\loigiai{
Đáp án A. Theo định nghĩa số mol, $n=m/M$.
}
\end{ex}

% ===== Câu 241 =====
% ID: L12C2B11-84-TN
% Mức: VD
\begin{ex}
% ID: L12C2B11-84-TN
% Mức: VD
Trong dạng “Đồ thị và các quá trình biến đổi trạng thái liên tiếp”, Khối lượng riêng của khí lí tưởng thỏa
\choice
{\True $\rho=pM/(RT)$}
{$\rho=RT/(pM)$}
{$\rho=pRT/M$}
{$\rho=p/(MRT)$}
\loigiai{
Đáp án A. Từ $pV=(m/M)RT$ và $\rho=m/V$.
}
\end{ex}

% ===== Câu 242 =====
% ID: L12C2B11-85-TN
% Mức: VDC
\begin{ex}
% ID: L12C2B11-85-TN
% Mức: VDC
Trong dạng “Đồ thị và các quá trình biến đổi trạng thái liên tiếp”, Hai bình cùng p, $V$, $T$ chứa khí lí tưởng khác loại thì có
\choice
{\True cùng số phân tử}
{cùng khối lượng}
{cùng khối lượng riêng}
{khác số mol}
\loigiai{
Đáp án A. Theo $n=pV/(RT)$, số mol và số phân tử bằng nhau.
}
\end{ex}

% ===== Câu 243 =====
% ID: L12C2B11-86-TN
% Mức: VDC
\begin{ex}
% ID: L12C2B11-86-TN
% Mức: VDC
Trong dạng “Đồ thị và các quá trình biến đổi trạng thái liên tiếp”, Khi giải bài toán khí lí tưởng, áp suất đưa vào phương trình phải là
\choice
{\True áp suất tuyệt đối}
{áp suất dư}
{áp suất do chất lỏng בלבד}
{áp suất tương đối tùy ý}
\loigiai{
Đáp án A. Phương trình trạng thái sử dụng áp suất tuyệt đối.
}
\end{ex}
